% !TeX program = xelatex
% !TeX encoding = UTF-8
% !BIB program = bibtex
%
% =====================================================================
%  CS 6330 - Introduction to Data Science
%  개인 학습서: 개념 정리 · 과제 기록 · 프로젝트 로그
%
%  저자 : Taek Lim
%  빌드 : latexmk -xelatex main.tex   (또는 make)
%  규칙 : ../AUTHORING.md
%
%  구조 (AUTHORING.md 3절)
%    preamble.tex   패키지 · 폰트 · 스타일 · semantic 매크로
%    parts/         파트별 \input 묶음
%    chapters/      NN-slug.tex - 한 장에 한 파일
%    bib/           references.bib
% =====================================================================
\documentclass[11pt, oneside]{book}

\input{preamble}

% 그림은 복사하지 않고 과제 폴더를 직접 참조한다. 코드를 다시 돌려 그림을
% 갱신하면 다음 빌드에서 책에도 그대로 반영된다.
\graphicspath{%
  {figures/}%
  {../hw1/figures/}%
  {../hw1/00_lecture_charts/figures/}%
  {../hw2/figures/}%
}

\title{\Huge\bfseries 데이터 사이언스 입문\\[4pt]
       \Large \textit{Introduction to Data Science}\\[10pt]
       \large 개념 정리 $\cdot$ 과제 기록 $\cdot$ 프로젝트 로그}
\author{Taek Lim}
\date{CS 6330 \textperiodcentered\ \today}

\begin{document}

\frontmatter
\maketitle
\chapter*{머리말}
\addcontentsline{toc}{chapter}{머리말}
\markboth{머리말}{머리말}

이 책은 CS 6330 \textit{Introduction to Data Science} 수업을 들으면서
\textbf{내가 이해한 내용}을 정리한 개인 학습서다. 교재의 요약이 아니라, 과제를
직접 풀고 그림을 직접 그려 보면서 \emph{왜 그렇게 되는지} 납득한 것만 적는다.

\section*{이 책을 관통하는 비유: 지도}

책 전체에서 하나의 비유를 일관되게 쓴다. 데이터 분석은 \analogy{지도 제작}이다.

\begin{table}[ht]\centering
\caption{책 전체에서 유지하는 비유 대응표}
\label{tab:analogy}
\begin{tabularx}{\linewidth}{@{}l l X@{}}
\toprule
개념 & 비유 & 대응의 의미 \\
\midrule
데이터(Data) & \analogy{실제 지형} & 있는 그대로의 전부. 너무 커서 한눈에 못 본다 \\
요약통계(Summary Statistics) & \analogy{지도} & 줄인 것. 줄였으므로 반드시 뭔가를 버렸다 \\
평균/중앙값(Mean/Median) & \analogy{지도의 중심점} & 어디를 대표로 삼을지의 선택 \\
퍼짐(Spread) & \analogy{지도의 축척} & 얼마나 넓은 범위를 담고 있는가 \\
시각화(Visualization) & \analogy{항공사진} & 덜 줄인 것. 지도가 버린 것이 보인다 \\
분포 가정(Distribution) & \analogy{투영법} & 편리하지만 반드시 왜곡이 따른다 \\
가설검정(Hypothesis Testing) & \analogy{현장 측량으로 지도 검증} & 지도(귀무가설)가 맞다면 이만큼 어긋난 측량이 나올 법한가를 묻는다. 너무 어긋나면 지도를 고치고(기각), 오차 안이면 고칠 근거가 없을 뿐 지도가 옳다는 증명은 아니다(기각 실패) \\
\bottomrule
\end{tabularx}
\end{table}

이 비유의 핵심은 하나다. \textbf{지도는 지형이 아니다.} 요약통계는 데이터가
아니고, 줄이는 과정에서 무엇을 버렸는지 모르면 남은 숫자를 잘못 읽는다.

\section*{쓰는 방식}

\begin{itemize}
  \item \textbf{용어는 병기한다.} 새로 도입하는 개념은 \term{왜도}(Skewness)처럼
        한국어와 영어를 함께 적는다. 수업과 논문은 영어로 되어 있고, 한국어로만
        기억하면 나중에 원문을 읽을 때 연결이 끊긴다.
  \item \textbf{기호는 국제 표기를 그대로 쓴다.} $\mu$, $\sigma$, $\IQR$.
  \item \textbf{절은 두괄식으로 연다.} 각 절 첫머리의 \textit{In brief.} 블록이
        Why $\cdot$ When $\cdot$ Where $\cdot$ How 를 한 줄씩 먼저 답한다.
  \item \textbf{내가 틀렸던 것을 남긴다.} \textbf{함정} 상자의 내용은 대부분 실제로
        한 번 틀린 것들이다. 맞은 것보다 틀린 것이 나중에 더 쓸모 있다.
  \item \textbf{그림은 코드로 생성한다.} 손으로 그리거나 값을 타이핑하지 않는다.
        그림 파일은 과제 폴더에서 직접 가져오므로 코드를 다시 돌리면 책의 그림도
        함께 갱신된다.
\end{itemize}

\begin{aside}[읽는 사람에게]
이 책의 독자는 6개월 뒤의 나다. 그때의 내가 ``이걸 왜 이렇게 했지?''라고 물을
만한 지점마다 답을 적어 두는 것이 목표다.
\end{aside}

\vfill
\noindent\textbf{Taek Lim}\\
\textit{The University of Texas Rio Grande Valley}

\tableofcontents

\mainmatter
% 1부 - 개념. 수업에서 배운 것을 내 말로 다시 세운다.
\part{기초}

\chapter{속성 유형 --- 값의 종류를 먼저 가른다}
\label{ch:attributes}
% source: deck 3 slide 1 (제목만, 내용 없음)

\begin{inbrief}
  \item[Why] 값에 허용된 연산(척도)과 데이터가 놓이는 모양(행렬, 공간)을
             알아야 이후 모든 요약과 그림이 합법이 된다.
  \item[When] 데이터를 받자마자, 어떤 통계도 계산하기 전에.
  \item[Where] 표로 된 모든 데이터. 이미지, 텍스트, 네트워크도 행렬로
               바꾼 뒤 같은 틀에 들어온다.
  \item[How] 열마다 척도와 이산/연속을 가르고, 식별자를 뺀 나머지를
             $n \times d$ 행렬로 본다.
\end{inbrief}

같은 숫자라도 무엇을 뜻하느냐에 따라 허용되는 계산이 다르다. 우편번호
78539와 칼로리 540은 둘 다 숫자로 저장되지만, 우편번호의 평균은 아무 의미가
없다. \term{속성 유형}(Attribute Type)을 먼저 가르는 이유는 여기에 있다.
어떤 요약값이 \emph{합법}인지가 여기서 결정된다.

\begin{keyidea}
속성 유형은 데이터를 예쁘게 분류하려는 작업이 아니다. \textbf{그 값에 어떤
연산이 허용되는지}를 정하는 작업이다. 유형을 틀리면 이후의 평균 $\cdot$ 표준편차
$\cdot$ 그래프가 전부 조용히 틀린다 --- 계산은 성공하고 숫자도 나오기 때문이다.
\end{keyidea}

% =====================================================================
\section{배경}
\label{sec:attr-background}

이 장은 강의 자료 ``Data and Matrix''(덱 3,
\texttt{notes/3\_DataAndMatrix.pptx})를 옮긴 것이다. 책의 첫 본문 장이므로
앞서 다룬 개념은 없지만, 이 장이 바로 쓰는 기초 용어 세 가지는 여기서
한 줄씩 깔아 둔다.

\begin{itemize}
  \item \term{표}(Table): 값을 행과 열의 격자에 늘어놓은 것.
  \item \term{행}(Row)과 \term{열}(Column): 표에서 가로줄이 행, 세로줄이
        열이다.
  \item \term{좌표평면}(Coordinate Plane): 점 하나가 값 두 개(가로 좌표,
        세로 좌표)로 정해지는 평면.
\end{itemize}

이 장은 두 부분으로 나뉜다. 앞부분(\Cref{sec:attr-flow}부터
\Cref{sec:attr-stats}까지)은 열 하나의 성질, 곧 값 하나에 어떤 연산이
허용되는가를 다룬다. 뒷부분(\Cref{sec:attr-discrete}부터)은 그렇게 정리한
열들을 모아 표 전체, 곧 행렬로 보는 법을 다룬다. 이 장에서 정한 유형은
이후 요약값(\Cref{ch:descriptive})과 그래프(\Cref{ch:visualization})를
고르는 전제가 된다.

% =====================================================================
\section{두 번의 질문으로 가른다}
\label{sec:attr-flow}
% source: deck 3 slide 3

\begin{inbrief}
  \item[Why] 값에 허용된 연산을 알아야 요약값과 그래프를 고를 수 있기 때문이다.
  \item[When] 데이터를 받자마자, 어떤 통계도 계산하기 전에 한다.
  \item[Where] 모든 변수에 적용한다. 특히 숫자로 저장된 범주형(우편번호, ID,
               코드)에서 판단이 갈린다.
  \item[How] 질문 두 개를 순서대로 던진다. ``숫자로 계산이 되는가'' $\to$
             ``0이 없음을 뜻하는가'' 또는 ``값들 사이에 순서가 있는가''.
\end{inbrief}

가름은 2단계로 끝난다. 1단계에서 \term{수치형}(Numeric)과
\term{범주형}(Categorical)을 나누고, 2단계에서 각각을 다시 둘로 나눈다
(\Cref{fig:attr-flow}).

\begin{figure}[ht]\centering
\begin{tikzpicture}[
  qbox/.style    = {draw = bookaccent, fill = bookaccentbg, rounded corners = 2pt,
                    minimum width = 46mm, minimum height = 8mm, align = center,
                    font = \small, inner sep = 2pt},
  leafbox/.style = {draw = bookrule, fill = booksurface, rounded corners = 2pt,
                    minimum width = 34mm, minimum height = 11mm, align = center,
                    font = \small, inner sep = 2pt},
  arrow/.style   = {-{Stealth[length=2.2mm]}, draw = bookmuted},
  elab/.style    = {font = \scriptsize\color{bookmuted}, inner sep = 1.5pt}]

  \node[qbox] (q1) at (0, 0) {1단계: 숫자로 계산이 되는가?};

  \node[qbox] (q2) at (-3.9, -1.6) {2단계: 0이 ``없음''인가?};
  \node[qbox] (q3) at ( 3.9, -1.6) {2단계: 순서가 있는가?};

  \node[leafbox] (ratio)    at (-5.8, -3.6) {비율(Ratio)\\[-2pt] \scriptsize 칼로리 $\cdot$ 금액};
  \node[leafbox] (interval) at (-2.0, -3.6) {구간(Interval)\\[-2pt] \scriptsize 섭씨온도 $\cdot$ 연도};
  \node[leafbox] (ordinal)  at ( 2.0, -3.6) {순서(Ordinal)\\[-2pt] \scriptsize 학력 $\cdot$ 만족도};
  \node[leafbox] (nominal)  at ( 5.8, -3.6) {명목(Nominal)\\[-2pt] \scriptsize 혈액형 $\cdot$ 우편번호};

  \draw[arrow] (q1.south) -- ++(0,-0.4) -| (q2.north)
        node[elab, pos=0.45, above] {수치형};
  \draw[arrow] (q1.south) -- ++(0,-0.4) -| (q3.north)
        node[elab, pos=0.45, above] {범주형};

  \draw[arrow] (q2.south) -- ++(0,-0.5) -| (ratio.north)
        node[elab, pos=0.45, above] {0 = 없음};
  \draw[arrow] (q2.south) -- ++(0,-0.5) -| (interval.north)
        node[elab, pos=0.45, above] {0 = 기준점};
  \draw[arrow] (q3.south) -- ++(0,-0.5) -| (ordinal.north)
        node[elab, pos=0.45, above] {서열 있음};
  \draw[arrow] (q3.south) -- ++(0,-0.5) -| (nominal.north)
        node[elab, pos=0.45, above] {서열 없음};
\end{tikzpicture}
\caption{속성 유형을 가르는 두 번의 질문. 1단계에서 수치형/범주형을 나누고,
2단계에서 각각을 다시 둘로 나눈다.}
\label{fig:attr-flow}
\end{figure}

% =====================================================================
\section{수치형 --- 구간과 비율}
\label{sec:attr-numeric}
% source: deck 3 slide 3

\begin{inbrief}
  \item[Why] 같은 수치형이라도 ``몇 배''라는 말이 성립하는 것과 아닌 것이 있기
             때문이다.
  \item[When] 1단계에서 수치형으로 가른 뒤에 던지는 질문이다.
  \item[Where] 온도 $\cdot$ 연도 $\cdot$ 지능지수는 구간, 무게 $\cdot$ 금액
               $\cdot$ 개수는 비율이다.
  \item[How] ``값 0이 정말로 아무것도 없다는 뜻인가''를 묻는다. 그렇다면 비율,
             기준점일 뿐이라면 구간이다.
\end{inbrief}

\begin{definition}[구간 척도와 비율 척도]
\label{def:interval-ratio}
\term{구간 척도}(Interval Scale)는 값의 \emph{간격}은 의미가 있지만 0이
임의의 기준점인 척도다. 덧셈과 뺄셈은 되고, 나눗셈(비율)은 되지 않는다.

\term{비율 척도}(Ratio Scale)는 0이 \term{절대 영점}(Absolute Zero), 즉
``그 속성이 전혀 없음''을 뜻하는 척도다. 덧셈 $\cdot$ 뺄셈에 더해
곱셈 $\cdot$ 나눗셈까지 의미를 가진다.
\end{definition}

\begin{example}[``두 배''가 성립하는가]
\label{ex:twice}
섭씨온도는 구간 척도다. $0^\circ\mathrm{C}$ 는 열이 없다는 뜻이 아니라 물이 어는
지점을 0으로 잡은 약속일 뿐이다. 그래서 $20^\circ\mathrm{C}$ 는
$10^\circ\mathrm{C}$ 보다 ``$10^\circ$ 더 덥다''라고는 말할 수 있어도
``두 배 덥다''라고는 말할 수 없다. 실제로 같은 온도를 화씨로 바꾸면 50\,\textdegree F
와 68\,\textdegree F 가 되어 비율이 $2.0$ 에서 $1.36$ 으로 바뀐다 --- 단위를
바꿨을 뿐인데 답이 달라진다면 그 계산은 애초에 의미가 없었던 것이다.

칼로리는 비율 척도다. 0\,kcal 는 열량이 실제로 없다는 뜻이므로,
540\,kcal 는 270\,kcal 의 ``두 배''가 맞다. 단위를 kJ 로 바꿔도 비율은 그대로다.
\end{example}

% =====================================================================
\section{범주형 --- 명목과 순서}
\label{sec:attr-categorical}
% source: deck 3 slide 3

\begin{inbrief}
  \item[Why] 범주에 순서가 있으면 중앙값과 정렬이 허용되고, 없으면 금지되기
             때문이다.
  \item[When] 1단계에서 범주형으로 가른 뒤에 던지는 질문이다.
  \item[Where] 혈액형 $\cdot$ 지역 $\cdot$ 음식 종류는 명목, 학력 $\cdot$ 만족도
               $\cdot$ 학점은 순서다.
  \item[How] ``값들 사이에 자연스러운 서열이 있는가''를 묻는다. 나열 순서를
             뒤바꿔도 아무 문제가 없다면 명목이다.
\end{inbrief}

\begin{definition}[명목 척도와 순서 척도]
\label{def:nominal-ordinal}
\term{명목 척도}(Nominal Scale)는 ``같다/다르다''만 판정할 수 있는 척도다.
서열이 없으므로 나열 순서는 임의다.

\term{순서 척도}(Ordinal Scale)는 값들 사이에 서열이 있는 척도다. 단, 등급
사이의 \emph{간격이 같다는 보장은 없다}.
\end{definition}

순서 척도에서 가장 자주 나오는 실수가 이 간격 문제다. 만족도를
``매우 불만족 $=1$, 불만족 $=2$, 보통 $=3$, 만족 $=4$, 매우 만족 $=5$''로
숫자를 붙이면 평균을 계산할 수 있게 되지만, 그 평균 3.4는
``1$\to$2 의 차이''와 ``4$\to$5 의 차이''가 같다고 가정한 값이다. 그런 가정을
한 적은 없다.

\begin{pitfall}[숫자로 저장되었다고 수치형이 아니다]
우편번호, 학번, 제품 코드, 등번호는 전부 정수로 저장되지만 \term{명목
척도}(Nominal Scale)다. \texttt{pandas} 는 이들을 \texttt{int64} 로 읽고
\texttt{.mean()} 도 아무 불평 없이 계산해 준다. 오류가 나지 않는다는 것이
가장 위험한 지점이다 --- 유형은 저장 타입(dtype)이 아니라 \textbf{의미}로
정한다.
\end{pitfall}

% =====================================================================
\section{유형이 통계를 결정한다}
\label{sec:attr-stats}

\begin{inbrief}
  \item[Why] 속성 유형을 가르는 목적 자체가 어떤 요약값이 합법인지 정하는
             것이기 때문이다.
  \item[When] \Cref{ch:descriptive}의 요약값을 고르기 직전에 확인한다.
  \item[Where] 중심값 선택(\Cref{sec:center})과 그래프 선택
               (\Cref{ch:visualization})에 그대로 이어진다.
  \item[How] \Cref{tab:attr-ops} 의 행을 위에서 아래로 훑어 허용된 연산까지만
             쓴다.
\end{inbrief}

허용되는 연산은 \term{명목}$\to$\term{순서}$\to$\term{구간}$\to$\term{비율}
순으로 누적된다. 위 단계에서 허용된 것은 아래 단계에서도 전부 허용된다
(\Cref{tab:attr-ops}).

\begin{table}[ht]\centering
\caption{속성 유형별로 허용되는 연산과 대표 통계. 아래로 갈수록 허용 범위가
넓어진다.}
\label{tab:attr-ops}
\begin{tabularx}{\linewidth}{@{}l l l X@{}}
\toprule
척도 & 허용 연산 & 대표 통계 & 예시 \\
\midrule
명목(Nominal)  & $=,\ \neq$                    & 최빈값(Mode)
               & 혈액형 $\cdot$ 우편번호 $\cdot$ 음식 종류 \\
순서(Ordinal)  & $+\ <,\ >$                    & $+$ 중앙값(Median)
               & 학력 $\cdot$ 만족도 $\cdot$ 학점 \\
구간(Interval) & $+\ -$                        & $+$ 평균(Mean) $\cdot$ 표준편차
               & 섭씨온도 $\cdot$ 연도 $\cdot$ 지능지수 \\
비율(Ratio)    & $+\ \times,\ \div$            & $+$ 기하평균 $\cdot$ 변동계수
               & 칼로리 $\cdot$ 무게 $\cdot$ 금액 $\cdot$ 개수 \\
\bottomrule
\end{tabularx}
\end{table}

이 표는 \Cref{ch:descriptive} 전체의 전제다. 최빈값은 네 유형 모두에서 쓸 수
있고, 중앙값은 순서 척도부터, 평균과 표준편차는 구간 척도부터 허용된다. 즉
\Cref{sec:center}에서 ``평균이냐 중앙값이냐''를 고민하는 일 자체가 \textbf{구간
척도 이상일 때만} 성립하는 질문이다. 명목 척도에서는 그 질문이 아예 없다.

\begin{example}[과제 데이터의 여덟 변수]
\label{ex:hw1-types}
HW1 의 패스트푸드 영양정보 데이터
(\Cref{ch:hw1})는 여덟 개 변수를 가진다. 유형을 가르면 이렇게 된다.

\begin{center}
\begin{tabularx}{0.95\linewidth}{@{}l l X@{}}
\toprule
변수 & 유형 & 그렇게 가른 이유 \\
\midrule
\texttt{restaurant}   & 명목(Nominal) & 식당 이름. 순서가 없다 \\
\texttt{item}         & 명목(Nominal) & 메뉴 이름. 고유 식별자에 가깝다 \\
\texttt{type}         & 명목(Nominal) & 음식 종류. 우열이 없다 \\
\texttt{serving\_size}& 비율(Ratio)   & 0\,g 은 양이 없다는 뜻 \\
\texttt{calories}     & 비율(Ratio)   & 0\,kcal 은 열량이 없다는 뜻 \\
\texttt{fat}          & 비율(Ratio)   & 0\,g 은 지방이 없다는 뜻 \\
\texttt{sodium}       & 비율(Ratio)   & 0\,mg 은 나트륨이 없다는 뜻 \\
\texttt{sugars}       & 비율(Ratio)   & 0\,g 은 당류가 없다는 뜻 \\
\bottomrule
\end{tabularx}
\end{center}

수치형 다섯 개가 전부 비율 척도이므로 평균 $\cdot$ 표준편차 $\cdot$ 비율 계산이
모두 합법이다. 반면 명목 변수 셋에는 최빈값과 빈도만 쓸 수 있다 ---
\texttt{restaurant} 의 ``평균''은 존재하지 않는다.
\end{example}

\begin{aside}[이 데이터에 구간 척도가 없는 이유]
구간 척도는 사람이 기준점을 약속해서 만든 눈금(온도, 연도, 시험 점수)에서
주로 나온다. 물리적 양을 측정한 값은 대개 0 이 곧 ``없음''이므로 비율 척도가
된다. 영양정보처럼 측정값만 모인 데이터에 구간 척도가 없는 것은 우연이 아니다.
\end{aside}

% =====================================================================
\section{이산형과 연속형}
\label{sec:attr-discrete}
% source: deck 3 slide 2, 4

\begin{inbrief}
  \item[Why] 척도(명목$\cdot$순서$\cdot$구간$\cdot$비율)만으로는 값 사이에
             다른 값이 들어갈 수 있는지 알 수 없다. 이 축이 막대그래프
             (Bar Chart) 대 히스토그램(Histogram), 확률질량함수(PMF)
             대 확률밀도함수(PDF)의 선택을 가른다.
  \item[When] 척도를 정한 뒤, 그래프 종류나 확률 모형을 고르기 전에
              확인한다.
  \item[Where] 수치형과 범주형 모두에 적용되는 두 번째 축이다.
  \item[How] 두 값 사이에 중간값을 가진 다른 값이 실제로 존재할 수
             있는지 묻는다. 있으면 연속형, 셀 수 있게 끊겨 있으면
             이산형이다.
\end{inbrief}

값이 놓인 척도(\Cref{sec:attr-numeric}, \Cref{sec:attr-categorical})와
값이 끊겨 있는지는 \textbf{다른 질문}이다. 태어난 연도(구간 척도)와
키(비율 척도)는 둘 다 수치형이지만, 태어난 연도는 정수로만 끊겨 있고
키는 두 값 사이에 항상 더 세밀한 값이 있다. 이 차이를 놓치면 이산형
값을 연속형 그래프(히스토그램)로, 연속형 값을 이산형 그래프(막대그래프)로
그리는 실수가 나온다.

\begin{definition}[이산형과 연속형]
\label{def:discrete-continuous}
\term{이산형}(Discrete)은 값이 서로 떨어져 있어 셀 수 있는 변수다(예:
태어난 연도, 참/거짓). \term{연속형}(Continuous)은 실수 구간의 값을
갖는 변수다(예: 길이). 두 값 사이에 항상 더 세밀한 값을 넣을 수 있으면
연속형이다.
\end{definition}

척도(\Cref{tab:attr-ops})와 이산/연속은 서로 다른 축이므로, 실제 변수는
둘을 \textbf{함께} 답해야 완전히 분류된다. \Cref{tab:attr-discrete-2x2}가
이 두 축을 겹쳐 보여준다.

\begin{table}[ht]\centering
\caption{척도(행)와 이산/연속(열)을 겹친 분류. ``없음''은 그 조합이
존재하지 않는다는 뜻이다.}
\label{tab:attr-discrete-2x2}
\begin{tabularx}{\linewidth}{@{}l X X@{}}
\toprule
 & 이산형(Discrete) & 연속형(Continuous) \\
\midrule
수치형(Numeric) & 태어난 연도(구간), 고객 수(비율)
  & 길이 $\cdot$ 무게(비율), 섭씨온도 $\cdot$ 경도(구간) \\
범주형(Categorical) & 참/거짓(명목), 전공(명목), 만족도 1$\sim$5(순서)
  & 없음 \\
\bottomrule
\end{tabularx}
\end{table}

범주형 열이 연속형 칸에서 비는 이유는 한 문장으로 정리된다. 범주형
값은 정해진 기호의 집합에서 하나를 고르므로, 두 기호 사이에 들어갈
``중간 기호''가 없다. 그래서 범주형은 언제나 이산형이다.

\begin{example}[다섯 항목의 척도와 이산/연속을 함께 답하기]
\label{ex:attr-slide4}
다섯 개 항목에 척도와 이산/연속을 함께 매긴다(만족도 문항의 실제 모양은
1$\sim$5점 라디오 버튼이다). 척도 답은 슬라이드와 같고, 이산/연속은
이 책이 추가했다.

\begin{center}
\begin{tabularx}{\linewidth}{@{}l l l X@{}}
\toprule
항목 & 척도 & 이산/연속 & 근거 \\
\midrule
색(Color) & 명목(Nominal) & 이산형 & 색 이름 사이에 서열이 없다 \\
경도(Longitude) & 구간(Interval) & 연속형 & $0^\circ$는 Greenwich를
  고른 약속일 뿐이다 \\
강사 만족도 1$\sim$5 & 순서(Ordinal) & 이산형 & 등급 사이 간격이
  같다는 보장이 없다(\Cref{def:nominal-ordinal}) \\
무게(Weight) & 비율(Ratio) & 연속형 & 0\,kg은 무게가 없다는 뜻이다 \\
전공(College Major) & 명목(Nominal) & 이산형 & 전공 이름 사이에
  서열이 없다 \\
\bottomrule
\end{tabularx}
\end{center}

경도는 $180^\circ$와 $-180^\circ$가 같은 지점을 가리키는 원형 값이라,
두 경도의 ``차이''를 구할 때도 이 점을 조심해야 한다.
\end{example}

% =====================================================================
\section{표 데이터와 데이터 행렬}
\label{sec:data-matrix}
% source: deck 3 slide 5-8

\begin{inbrief}
  \item[Why] 열 하나씩 따로 보는 것만으로는 ``객체 하나가 여러 속성을
             동시에 가진다''는 사실을 다룰 수 없다. 행과 열을 한 번에
             다루는 틀이 필요하다.
  \item[When] 척도와 이산/연속을 다 가른 뒤, 표 전체를 계산에 넣기
              직전에.
  \item[Where] 행마다 하나의 대상(사람, 제품, 관측)을 담고, 열마다
               하나의 속성을 담는 모든 표.
  \item[How] 표를 $n \times d$ 행렬로 본다. $n$은 행(객체) 수, $d$는
             열(속성) 수다.
\end{inbrief}

열 하나씩 값을 따로 읽으면, 예를 들어 급여 열과 나이 열을 각각 볼 수는
있어도 ``이 사람의 급여와 나이''처럼 한 대상이 가진 여러 속성을 한꺼번에
계산에 넣을 방법이 없다. \term{표 데이터}(Tabular Data)는 행 하나가
대상 하나, 열 하나가 속성 하나를 담도록 정리해 이 문제를 없앤다. 이렇게
정리된 표는 $n \times d$ \term{데이터 행렬}(Data Matrix) $\mathbf{D} \in
\mathbb{R}^{n \times d}$로 쓸 수 있다. $n$은 행의 수(대상 수), $d$는
열의 수(속성 수)다.

행렬의 행 하나는
\[
  \mathbf{x}_i = (x_{i1}, x_{i2}, \dots, x_{id})^T \in \mathbb{R}^d
\]
로 쓴다. 대상 $i$번째의 속성 값 $d$개를 순서대로 늘어놓은 것이다.
행렬의 열 하나는
\[
  X_j = (x_{1j}, x_{2j}, \dots, x_{nj})^T \in \mathbb{R}^n
\]
로 쓴다. 속성 $j$번째의 값을 모든 대상에서 모은 것이다. \textbf{$\mathbf{x}_i$
(굵은 소문자, 행)와 $X_j$(대문자, 열)는 다른 것이다.} 첨자가 대상을
가리키면 행, 속성을 가리키면 열이다. 분야마다 행과 열을 부르는 이름이
다른데, \Cref{tab:attr-synonyms}가 그 대응을 정리한다.

\begin{table}[ht]\centering
\caption{데이터 행렬의 행과 열을 부르는 여러 이름. 같은 것을 분야마다
다르게 부른다.}
\label{tab:attr-synonyms}
\begin{tabularx}{\linewidth}{@{}l X@{}}
\toprule
대상 & 동의어 \\
\midrule
행($\mathbf{x}_i$) & 표본(Sample), 인스턴스(Instance), 레코드(Record),
  객체(Object), 점(Point), 특징 벡터(Feature Vector) \\
열($X_j$) & 속성(Attribute), 특징(Feature), 차원(Dimension),
  변수(Variable), 필드(Field) \\
\bottomrule
\end{tabularx}
\end{table}

\begin{example}[이름, 급여, 나이 표를 행렬로 보기]
\label{ex:attr-name-salary-age}
\Cref{tab:attr-name-salary-age}는 이름, 급여, 나이 세 열로 된 표다.

\begin{table}[ht]\centering
\caption{이름, 급여, 나이 표.}
\label{tab:attr-name-salary-age}
\begin{tabularx}{0.7\linewidth}{@{}l X X@{}}
\toprule
이름(Name) & 급여(Salary, \$) & 나이(Age, Years) \\
\midrule
Jane    & 90000 & 52 \\
John    & 85000 & 48 \\
Delilah & 75000 & 32 \\
Dave    & 90000 & 53 \\
Ellen   & 82000 & 44 \\
\bottomrule
\end{tabularx}
\end{table}

세는 순서는 한 줄씩이다.
\begin{align*}
  &\text{머리글 1행 제외} \to 5\text{행}
    && \text{Jane부터 Ellen까지 5명} \\
  &\text{열 3개}
    && \text{이름, 급여, 나이} \\
  &\text{이름(식별자) 열 제외} \to \text{속성 2개}
    && \text{급여, 나이만 속성이다} \\
  &\mathbf{D} \in \mathbb{R}^{5 \times 2}
    && \text{데이터 행렬 크기}
\end{align*}
\end{example}

\begin{pitfall}[식별자 열은 속성이 아니다]
\Cref{tab:attr-name-salary-age}의 이름(Name) 열처럼 대상을 가리키기만
하는 열은 속성으로 세지 않는다. 이 규칙을 놓치면 이 표와, 뒤에 나올
급여$\cdot$나이$\cdot$근속연수 표인 \Cref{tab:attr-3d-data}가 헷갈린다.
두 표 모두 ``5행 3열''이라 부를 수 있지만, 두 슬라이드가 세는 대상이
다르다.
이 표의 3열은 이름(식별자) 포함 3열이라 속성은 2개이고,
\Cref{tab:attr-3d-data}의 3열은 이름 열이 빠지고 근속연수가 더해진
3열이라 속성이 3개다. 규칙 하나로 정리하면 이렇다. \textbf{식별자는
속성이 아니다. 차원은 식별자를 뺀 열의 수다.} \Cref{ex:hw1-types}의
\texttt{item}도 고유값이 126행 중 107개뿐이라 완전한 식별자는 아니지만,
원칙은 같다. 대상 하나하나를 실제로 가리키는 것은 행 번호이지 그 열의
값이 아니다.
\end{pitfall}

% =====================================================================
\section{차원의 두 가지 뜻}
\label{sec:dimension}
% source: deck 3 slide 9-12

\begin{inbrief}
  \item[Why] 같은 낱말 ``차원''이 프로그래밍과 데이터 과학에서 다른
             것을 센다. 섞어 쓰면 ``속성이 100개인 데이터는 100차원
             배열이 필요하다''는 오해가 생긴다.
  \item[When] 데이터 행렬(\Cref{sec:data-matrix})의 크기를 말할 때마다.
  \item[Where] 저장 방식(배열)을 말할 때와 데이터가 놓인 공간을 말할
               때 모두에 쓰이는 말이라 헷갈린다.
  \item[How] ``지금 세는 것이 배열의 축인가, 속성의 개수인가''를 먼저
             구분한다.
\end{inbrief}

``차원''이라는 말 하나로 표의 축 개수와 속성 개수를 동시에 가리키면,
속성이 100개인 표를 ``100차원 배열''로 저장해야 한다는 오해가 생긴다.
표는 속성이 몇 개든 행과 열, 축 2개짜리 배열일 뿐이다. 이 오해를
없애려면 두 뜻을 다른 이름으로 나눠야 한다. \Cref{tab:attr-dimension-two-meanings}가
이 두 뜻을 나눈다. \term{배열 차원}(Array Dimension)은 값에 접근할
때 쓰는 인덱스(축)의 개수다. \term{공간 차원}(Space Dimension)은
데이터 행렬의 속성 수 $d$, 즉 각 점이 놓인 공간의 차원이다.

\begin{table}[ht]\centering
\caption{프로그래밍의 차원과 데이터 과학의 차원.}
\label{tab:attr-dimension-two-meanings}
\begin{tabularx}{\linewidth}{@{}l X@{}}
\toprule
뜻 & 세는 대상 \\
\midrule
배열 차원(Array Dimension) & 값에 접근하는 인덱스(축)의 수. 표(행,
  열)는 축이 2개이므로 항상 배열 차원 2다 \\
공간 차원(Space Dimension) & 속성의 수 $d$. 속성이 3개면 그 점들은
  3차원 공간에 놓인다 \\
\bottomrule
\end{tabularx}
\end{table}

\begin{pitfall}[슬라이드 9: 차원 = 셀의 수?]
슬라이드 값은 ``프로그래밍의 차원(Dimension) = 셀의 수''이고, 다시
계산하면 ``차원 = 인덱스(축)의 수''다. 표 하나는 셀이 아무리 많아도
축은 행과 열 둘뿐이라 배열 차원은 항상 2다. 슬라이드 11 자신도 아래에서
같은 표를 ``2차원 배열(2-dimensional array)''이라 부른다.
\end{pitfall}

\paragraph{2차원 예.} 이름/급여/나이 표인 \Cref{tab:attr-name-salary-age}에서
이름 열을 빼면 급여$\cdot$나이 두 속성만 남고, 각 사람은 2차원 평면의
점 하나가 된다(\Cref{fig:attr-2d-scatter}). Jane(52세, 90000)과
Dave(53세, 90000)는 나이만 한 살 다를 뿐이라 그림에서 거의 겹쳐
보인다.

\begin{figure}[ht]\centering
  \includegraphics[width=0.75\linewidth]{03_01_salary_age_2d.png}
  \caption{급여$\cdot$나이 2차원 산점도(Scatter Plot). 5명이 각각 점
  하나에 대응한다.}
  \label{fig:attr-2d-scatter}
\end{figure}

\paragraph{3차원 예.} 같은 다섯 사람에게 근속연수(Years in Service)를
더하면 \Cref{tab:attr-3d-data}가 된다. 이 표는 이름 없이 급여, 나이,
근속연수 세 열뿐이다.

\begin{table}[ht]\centering
\caption{급여, 나이, 근속연수 표.}
\label{tab:attr-3d-data}
\begin{tabularx}{0.6\linewidth}{@{}X X X@{}}
\toprule
급여(Salary, \$) & 나이(Age, Years) & 근속연수(Years in Service) \\
\midrule
90000 & 52 & 10 \\
85000 & 48 & 20 \\
75000 & 32 & 30 \\
90000 & 53 & 40 \\
82000 & 44 & 20 \\
\bottomrule
\end{tabularx}
\end{table}

\begin{align*}
  &5 \times 3 = 15 && \text{배열의 셀 수} \\
  &\text{인덱스(축) 2개(행, 열)} && \text{배열 차원 2} \\
  &\text{속성 3개(급여, 나이, 근속연수)} && \text{공간 차원 3}
\end{align*}

이름은 행렬에는 들어가지 않지만, 그림에서 점의 라벨로는 남는다
(\Cref{fig:attr-3d-scatter}).

\begin{figure}[ht]\centering
  \includegraphics[width=0.8\linewidth]{03_02_salary_age_service_3d.png}
  \caption{급여$\cdot$나이$\cdot$근속연수 3차원 산점도. 이름은 점의
  라벨로만 남는다.}
  \label{fig:attr-3d-scatter}
\end{figure}

\paragraph{속성이 100개라도 배열은 2차원이다.} 표 하나는 속성이 몇
개든 $n \times d$ 꼴이므로, $d = 100$이어도 여전히 축은 행과 열
둘뿐이라 배열 차원은 2다. 반면 그 점들이 놓인 공간 차원은 100이다.

\begin{pitfall}[슬라이드 12: 고차원 데이터 = 행이 많음?]
슬라이드 값은 ``고차원 데이터(High Dimension Data)는 행이 많은 것도
포함한다''이고, 다시 계산하면 ``고차원은 열(속성) 수 $d$가 클 때만
해당한다''다. 슬라이드 9 자신의 정의(차원 수 = 특징 수)를 그대로
적용하면, 행이 많은 것은 대상 수 $n$이 큰 대용량(Large-Scale) 데이터
이지 고차원 데이터가 아니다.
\end{pitfall}

\begin{pitfall}[슬라이드 9, 12: 마지막 열 머리글]
슬라이드 값은 표의 마지막 열 머리글 ``Feature $n$''이고, 다시
계산하면 ``Feature $d$''다. \Cref{sec:data-matrix}의 기호에서 $n$은
행(대상) 수이고 열(속성) 수를 가리키는 기호는 $d$이므로, 이 머리글은
기호가 뒤바뀐 것이다.
\end{pitfall}

% =====================================================================
\section{점, 행, 벡터}
\label{sec:vector}
% source: deck 3 slide 13-15

\begin{inbrief}
  \item[Why] 행을 그냥 ``숫자 목록''으로만 보면 두 대상 사이의 거리나
             평균 점 같은 기하 연산을 떠올릴 수 없다.
  \item[When] 거리, 평균, 방향처럼 기하학적으로 데이터를 다룰 때.
  \item[Where] 데이터 행렬의 모든 행. 물리의 벡터와 같은 방식으로
               계산한다.
  \item[How] 행 $\mathbf{x}_i$를 원점에서 그 점까지 뻗은 화살표로
             본다.
\end{inbrief}

행 $\mathbf{x}_i$를 숫자가 나열된 목록으로만 보면, 두 대상 사이의
거리나 다섯 대상의 평균 점 같은 기하 연산을 떠올릴 방법이 없다.
데이터 과학의 \term{벡터}(Vector)는 그 목록을 물리의 벡터와 똑같이,
표의 행 하나이자 공간의 점 하나로, 원점에서 그 점까지 뻗은 화살표로
본다.
\Cref{tab:attr-vectors-source}의 다섯 점(나이, 급여)을 화살표로 그리면
\Cref{fig:attr-vectors}가 된다.

\begin{table}[ht]\centering
\caption{나이, 급여 표(벡터 예).}
\label{tab:attr-vectors-source}
\begin{tabularx}{0.5\linewidth}{@{}X X@{}}
\toprule
나이(Age) & 급여(Salary) \\
\midrule
20 & 90000 \\
30 & 85000 \\
28 & 40000 \\
40 & 95000 \\
35 & 42000 \\
\bottomrule
\end{tabularx}
\end{table}

\begin{figure}[ht]\centering
  \includegraphics[width=0.75\linewidth]{03_03_vectors_origin.png}
  \caption{원점에서 각 점까지 뻗은 화살표. 점 하나가 벡터 하나다.}
  \label{fig:attr-vectors}
\end{figure}

\paragraph{붓꽃 예.} 이 데이터는 Fisher의 붓꽃 데이터(Iris Data)다.
자세한 출처는 \Cref{sec:qq-iris-source}에서 다룬다. 이 책은
\texttt{sklearn.datasets.load\_iris()}로 값만 읽는다.
\Cref{tab:attr-iris-rows}가 그중 10행을 옮긴다.

\begin{table}[ht]\centering
\caption{붓꽃 데이터 10행. $\mathbf{x}_1 \sim \mathbf{x}_8$은 setosa,
$\mathbf{x}_{149}, \mathbf{x}_{150}$은 virginica다.}
\label{tab:attr-iris-rows}
\begin{tabularx}{\linewidth}{@{}l X X X X l@{}}
\toprule
행 & 꽃받침 길이 & 꽃받침 너비 & 꽃잎 길이 & 꽃잎 너비 & 종(Class) \\
 & (Sepal Length) & (Sepal Width) & (Petal Length) & (Petal Width) & \\
\midrule
$\mathbf{x}_1$   & 5.1 & 3.5 & 1.4 & 0.2 & setosa \\
$\mathbf{x}_2$   & 4.9 & 3.0 & 1.4 & 0.2 & setosa \\
$\mathbf{x}_3$   & 4.7 & 3.2 & 1.3 & 0.2 & setosa \\
$\mathbf{x}_4$   & 4.6 & 3.1 & 1.5 & 0.2 & setosa \\
$\mathbf{x}_5$   & 5.0 & 3.6 & 1.4 & 0.2 & setosa \\
$\mathbf{x}_6$   & 5.4 & 3.9 & 1.7 & 0.4 & setosa \\
$\mathbf{x}_7$   & 4.6 & 3.4 & 1.4 & 0.3 & setosa \\
$\mathbf{x}_8$   & 5.0 & 3.4 & 1.5 & 0.2 & setosa \\
$\mathbf{x}_{149}$ & 6.2 & 3.4 & 5.4 & 2.3 & virginica \\
$\mathbf{x}_{150}$ & 5.9 & 3.0 & 5.1 & 1.8 & virginica \\
\bottomrule
\end{tabularx}
\end{table}

이 표 전체는 150행 5열이다. 수치 열은 꽃받침 길이$\cdot$너비, 꽃잎
길이$\cdot$너비 4개뿐이다. 종(Class) 열 $X_5$는 이름(setosa,
versicolor, virginica)을 담은 명목 범주형(\Cref{sec:attr-categorical})
이라, 수치 데이터 행렬 $\mathbf{D} \in \mathbb{R}^{150 \times 4}$에서는
뺀다.

\paragraph{평균 점(Mean Point).} 꽃받침 길이와 너비 두 열만 놓고 150개 점을 찍은
뒤 평균 점을 표시하면 \Cref{fig:attr-iris-mean}이 된다. 평균 점은
열마다 따로 평균을 낸 벡터 $(\bar X_1, \bar X_2)$다.

\begin{align*}
  &\bar X_1 = 5.8433 && \text{꽃받침 길이 150개의 평균} \\
  &\bar X_2 = 3.0573 && \text{꽃받침 너비 150개의 평균}
\end{align*}

\begin{figure}[ht]\centering
  \includegraphics[width=0.75\linewidth]{03_04_iris_points_mean.png}
  \caption{붓꽃 꽃받침 길이$\times$너비 150개 점과 평균 점
  $(5.8433,\ 3.0573)$.}
  \label{fig:attr-iris-mean}
\end{figure}

% =====================================================================
\section{무엇이든 행렬로}
\label{sec:everything-matrix}
% source: deck 3 slide 16-19

지금까지는 이미 표로 정리된 데이터를 다뤘다. 그런데 이미지, 텍스트,
네트워크처럼 표 모양이 아닌 데이터도 있다. 이런 데이터도 행과 열을
정해 같은 $n \times d$ 데이터 행렬 틀에 넣을 수 있다.

\paragraph{이미지.} 이미지 하나를 여러 영역 $r_1, r_2, \dots$로 나누고,
영역마다 RGB 색이나 질감(Texture) 같은 특징을 뽑으면, 영역이 행이
되고 특징이 열이 되는 행렬이 만들어진다(\Cref{fig:attr-image-matrix}).
슬라이드에 실린 원본 사진 두 장은 출처가 밝혀지지 않아 이 책에는
싣지 않는다.

\begin{figure}[ht]\centering
\begin{tikzpicture}[
  imgbox/.style = {draw=bookaccent, fill=bookaccentbg, rounded corners=2pt,
                    minimum width=20mm, minimum height=20mm, align=center,
                    font=\small},
  rbox/.style   = {draw=bookrule, fill=booksurface, rounded corners=2pt,
                    minimum width=16mm, minimum height=8mm, align=center,
                    font=\scriptsize},
  arrow/.style  = {-{Stealth[length=2.2mm]}, draw=bookmuted}]

  \node[imgbox] (img) at (0,0) {이미지\\(Image)};
  \node[rbox] (r1) at (3.0, 1.0) {영역 $r_1$};
  \node[rbox] (r2) at (3.0, -1.0) {영역 $r_2$};
  \draw[arrow] (img.east) -- (r1.west);
  \draw[arrow] (img.east) -- (r2.west);

  \node[align=left, font=\scriptsize] (tbl) at (7.6, 0) {%
    \begin{tabular}{@{}l l l@{}}
    \toprule
    영역 & RGB 색 & 질감 \\
    \midrule
    $r_1$ & \dots & \dots \\
    $r_2$ & \dots & \dots \\
    \bottomrule
    \end{tabular}};
  \draw[arrow] (r1.east) -- ++(0.5,0) |- ($(tbl.west)+(0,0.55)$);
  \draw[arrow] (r2.east) -- ++(0.5,0) |- ($(tbl.west)+(0,-0.55)$);
\end{tikzpicture}
\caption{이미지를 영역으로 나누고 영역마다 특징(RGB 색, 질감)을 뽑으면
영역이 행, 특징이 열인 데이터 행렬이 된다.}
\label{fig:attr-image-matrix}
\end{figure}

\paragraph{텍스트.} 단어 하나마다 그 단어와 함께 나타나는(동시출현,
Co-occurrence) 다른 단어들의 통계를 모아 벡터로 만든다. 슬라이드의
``King'' 벡터는 값이 50개다(앞 5개만 보이면 $0.50451$, $0.68607$,
$-0.59517$, $-0.022801$, $0.60046$, 뒤에 45개 더, 모두 50개). 슬라이드는 이
벡터 연산으로 king $-$ man $+$ woman $\approx$ queen이 성립한다고
보이지만, 슬라이드에는 king 벡터 하나만 실려 있어 이 책에서 다시
계산할 수는 없다. 이 기법에 붙는 모형 이름은 슬라이드에 나오지
않으므로 책이 임의로 붙이지 않는다.

\paragraph{네트워크.} 노드(Node) 여러 개와 그 사이를 잇는 간선(Edge)
으로 이뤄진 네트워크는 노드를 행과 열로 삼은 \term{인접 행렬}
(Adjacency Matrix)로 바꾼다. 두 노드 사이에 간선이 있으면 그 칸이 1,
없으면 0이다. \Cref{fig:attr-network}가 노드 5개짜리 네트워크를
인접 행렬로 바꾼 예다.

\begin{figure}[ht]\centering
\begin{tikzpicture}[
  netdot/.style = {draw=bookaccent, fill=bookaccentbg, circle,
                    minimum size=7mm, font=\small},
  eline/.style  = {draw=bookmuted, thick},
  arrow/.style  = {-{Stealth[length=2.5mm]}, draw=bookaccent, thick}]

  \node[netdot] (A) at (0, 1.3) {A};
  \node[netdot] (B) at (1.6, 1.3) {B};
  \node[netdot] (C) at (0, 0) {C};
  \node[netdot] (D) at (1.6, 0) {D};
  \node[netdot] (E) at (3.0, 0.65) {E};

  \draw[eline] (A) -- (B);
  \draw[eline] (A) -- (C);
  \draw[eline] (A) -- (D);
  \draw[eline] (B) -- (D);
  \draw[eline] (B) -- (E);
  \draw[eline] (C) -- (D);
  \draw[eline] (D) -- (E);
  \draw[eline] (D) to[out=300, in=330, looseness=8] (D);

  \draw[arrow] (4.0, 0.65) -- (5.0, 0.65);

  \node[font=\scriptsize] at (8.6, 0.65) {%
    $\begin{matrix}
       & A & B & C & D & E \\
      A & 0 & 1 & 1 & 1 & 0 \\
      B & 1 & 0 & 0 & 1 & 1 \\
      C & 1 & 0 & 0 & 1 & 0 \\
      D & 1 & 1 & 1 & 1 & 1 \\
      E & 0 & 1 & 0 & 1 & 0
     \end{matrix}$};
\end{tikzpicture}
\caption{노드 5개짜리 네트워크(왼쪽)를 인접 행렬(오른쪽)로 바꾼다.
D의 자기 고리(Self-loop)가 대각 성분 $A_{DD}=1$이 된다.}
\label{fig:attr-network}
\end{figure}

\begin{align*}
  &3+3+2+5+2 = 15 && \text{행 합(A, B, C, D, E 순)} \\
  &2 \times 7 + 1 = 15 && \text{간선 7개는 두 번, 자기 고리는 한 번
    센 값과 일치}
\end{align*}

인접 행렬은 대칭이다($\mathbf{A} = \mathbf{A}^T$). 간선에 방향이
없으므로 $i$행 $j$열 값과 $j$행 $i$열 값이 같기 때문이다.

% =====================================================================
\section{어떻게 적용하는가}
\label{sec:attr-apply}

HW1 의 패스트푸드 영양정보 데이터(\Cref{ch:hw1}, \Cref{ex:hw1-types})를
이 장의 틀로 세 번 센다.

\begin{align*}
  &126 \times 8 && \text{저장된 표 크기(행 126, 열 8)} \\
  &8 - 1 = 7 && \text{\texttt{item} 열을 빼면 속성 7개} \\
  &\mathbf{D} \in \mathbb{R}^{126 \times 5} && \text{수치 데이터
    행렬(비율 척도 5개 열)}
\end{align*}

\begin{pitfall}[\texttt{item} 열은 식별자가 아니다]
\texttt{item} 열은 식별자처럼 보이지만 126행에 고유값이 107개뿐이다.
같은 메뉴 이름이 여러 식당에 있기 때문이다(예: ``Cheeseburger''는
McDonald's와 Burger King 모두에 있다). 행을 하나로 가리키는 것은
행 번호이지 \texttt{item} 값이 아니다. 그래도 수치 데이터 행렬은
그대로 $\mathbf{D} \in \mathbb{R}^{126 \times 5}$다. \texttt{item}은
애초에 수치 열이 아니었기 때문이다.
\end{pitfall}

\paragraph{거리 하나로 보는 단위의 문제.} 2차원 예
(\Cref{fig:attr-2d-scatter})의 나이$\cdot$급여로 \term{유클리드
거리}(Euclidean Distance)를 재보면, 단위가 큰 급여가 거리를
압도한다.

\begin{align*}
  &\sqrt{1^2 + 0^2} = 1.0000 && \text{Jane–Dave 원래 척도 거리(나이
    1살, 급여 \$0 차이)} \\
  &\sqrt{4^2 + 5000^2} = 5000.0016 && \text{Jane–John 원래 척도 거리
    (나이 4살, 급여 \$5000 차이)}
\end{align*}

나이 차이 4는 급여 차이 \$5000 앞에서 거리에 거의 반영되지 않는다.
두 축을 각 축의 범위(나이 범위 21, 급여 범위 15000)로 나누는
\term{정규화}(Normalization)를 한 뒤 다시 재면, 어느 쪽이 더
가까운지(Jane–Dave가 Jane–John보다 가깝다)는 그대로다. 바뀌는 것은
나이가 거리에 차지하는 몫이다.

\begin{align*}
  &4 / 21 = 0.1905 && \text{Jane–John 정규화 나이 몫} \\
  &5000 / 15000 = 0.3333 && \text{Jane–John 정규화 급여 몫} \\
  &\sqrt{0.1905^2 + 0.3333^2} = 0.3839 && \text{Jane–John 정규화 거리} \\
  &1 / 21 = 0.0476 && \text{Jane–Dave 정규화 거리}
\end{align*}

정규화 전에는 급여 몫이 거리를 거의 전부 차지해(원래 척도 거리
5000.0016 중 급여가 5000) 나이는 있으나 마나였다. 정규화 뒤에는
나이 몫 0.1905와 급여 몫 0.3333이 같은 척도(0에서 1)에서 거리에 들어간다.
Jane–Dave가 Jane–John보다 가깝다는 순서는 정규화 전후로 바뀌지
않는다(0.0476 < 0.3839, 1.0000 < 5000.0016). 정규화를 언제,
어떻게 하는지는 \Cref{sec:normalization}에서 다룬다.

% =====================================================================
\section{앞으로}
\label{sec:attr-future}

\begin{itemize}
  \item \textbf{고차원(High Dimension)과 차원의 저주(Curse of
        Dimensionality).} 속성 수 $d$가 커지면 점들 사이의 거리가
        비슷해지고 계산량도 늘어난다. 이 문제와 그 처방인 차원 축소
        (Dimensionality Reduction)는 이 덱에 없다. 강의계획서(덱 1)
        슬라이드 4의 주제 목록에만 이름이 올라 있다.
  \item \textbf{이산화(Discretization).} 연속형 속성을 이산형 속성으로 바꾸는 다리는
        \Cref{sec:discretization}에서 다룬다. \Cref{tab:attr-discrete-2x2}
        의 열을 오른쪽에서 왼쪽으로 옮기는 것과 같다.
  \item \Cref{ch:hypothesis}의 가설검정은 변수가 어떤 척도인지 이
        장에서 정한 판단을 그대로 전제로 삼는다.
\end{itemize}

\chapter{기술통계 --- 숫자로 데이터를 요약하기}
\label{ch:descriptive}

\begin{inbrief}
  \item[Why] 데이터를 몇 개의 숫자로 줄이려면, 먼저 ``값이 나올 가능성''을
             적는 언어인 확률변수(Random Variable), 확률질량함수(PMF),
             확률밀도함수(PDF), 누적분포함수(CDF)가 필요하다.
  \item[When] 새 데이터를 받아 첫 요약표를 만들 때, 그리고 분포 모형이나
              가설검정(Hypothesis Testing)으로 넘어가기 전에 쓴다.
  \item[Where] 주사위, 붓꽃(Iris) 꽃받침 길이(Sepal Length), 과제 데이터의
               \texttt{sugars}(당류) 변수에 적용한다.
  \item[How] 확률변수로 값을 정의하고, PMF/PDF와 CDF로 분포를 적은 뒤,
             중심(Center, 평균, 중앙값, 최빈값(Mode))과
             퍼짐(Spread, 범위(Range), 표준편차, IQR)을
             저항성(Resistance) 기준으로 고른다.
\end{inbrief}
% source: deck 4 slide 1

\term{기술통계}(Descriptive Statistics)는 데이터를 몇 개의 숫자로 줄이는
도구다. \analogy{지도}를 그리는 일이다(\Cref{tab:analogy}). 줄인다는 것은
반드시 무언가를 버린다는 뜻이고, 무엇을 버렸는지 모르면 남은 숫자를 잘못 읽는다.

이 장은 네 가지 요약값을 다룬다. 중심(Center), 퍼짐(Spread), 비대칭(Skewness),
그리고 값이 아예 없는 경우(Missing Value)다.

\begin{keyidea}
요약통계와 그림은 \textbf{한 쌍으로 다녀야 한다}. 요약값은 ``무엇이 일어났는가''
를 알려 주고, 그림은 ``왜 그런 값이 나왔는가''를 알려 준다. 둘 중 하나만으로
내린 결론은 이 장에서 계속 반복해서 틀린다.
\end{keyidea}

% =====================================================================
\section{배경}
\label{sec:desc-background}
% source: deck 4 slide 1-3

이 장의 출처는 \texttt{notes/4\_BasicStats.pptx}(23 슬라이드) 하나다.
슬라이드가 외부 문헌을 밝히지 않는 자리는 슬라이드 자체를 출처로 적는다.

기술통계를 계산하려면 그 전에 데이터를 어떻게 표로 담고, 각 값이 얼마나
자주 나오는지 어떻게 적을지부터 정해야 한다. 이 장은 다음 두 가지를
그대로 가져와 쓴다.

\begin{itemize}
  \item 속성(Attribute)과 척도: 값이 숫자인지 범주인지, 순서가 있는지를
        가른다(\Cref{ch:attributes}, \Cref{sec:attr-numeric}).
  \item 데이터 행렬(Data Matrix) $n \times d$: 행이 객체, 열이 속성인 표
        (\Cref{sec:data-matrix}).
\end{itemize}

이 장에서 계속 쓰는 데이터는 붓꽃(Iris) 데이터셋이다. 150송이의 붓꽃마다
꽃받침 길이(Sepal Length)를 포함한 네 가지 속성을 측정했고, 꽃받침
길이는 4.3\,cm에서 7.9\,cm 사이에 분포한다.

여러 속성 중 하나씩만 골라 한 번에 한 속성만 들여다보는 일을
\term{일변량 분석}(Univariate Analysis)이라고 부른다.
% source: deck 4 slide 14
이때 표의 한 열, 즉 속성 하나가 이 장에서 다룰 확률변수 하나에
대응한다.

이 장의 비유는 \Cref{ch:descriptive}의 \analogy{지도}를 그대로 잇는다.
확률분포(Probability Distribution)는 지형 전체이고, 이 장에서 배우는
평균, 중앙값, 표준편차 같은 요약값은 그 지형 위에 찍은 몇 개의 표지판이다.
표지판만 보고 지형 전체를 상상할 수 없듯, 요약값 몇 개만으로 분포
전체를 복원할 수는 없다.

% =====================================================================
\section{확률변수}
\label{sec:random-variable}
% source: deck 4 slide 4-5

\begin{inbrief}
  \item[Why] 관측값을 계산에 넣을 수 있는 숫자로 바꾸는 규칙이 필요하다.
  \item[When] 속성을 확률 계산이나 수식에 넣기 전에 정의한다.
  \item[Where] 수치형 속성은 있는 그대로 쓰고, ``조건을 만족하는가''
               같은 판정은 0과 1로 바꾼다.
  \item[How] 표본공간의 각 결과에 실수 하나를 대응시킨다.
\end{inbrief}

\paragraph{고치는 문제.} ``이 붓꽃은 꽃받침이 길다''라는 말은 그 자체로는
계산할 수 없다. 평균을 내거나 확률을 구하려면 먼저 이 말을 숫자로 바꿔야
한다.

\begin{definition}[확률변수]
\label{def:random-variable}
\term{표본공간}(Sample Space) $\mathcal{O}$는 실험에서 나올 수 있는 모든
결과의 집합이다. \term{확률변수}(Random Variable) $X$는 표본공간의 각
결과를 실수 하나로 대응시키는 함수다.
\[
  X : \mathcal{O} \to \mathbb{R}
\]
$X$가 유한 개 또는 셀 수 있는(Countably Infinite) 값만 가지면
\term{이산 확률변수}(Discrete Random Variable), $\mathbb{R}$의 어떤
구간에서도 값을 가질 수 있으면 \term{연속 확률변수}(Continuous Random
Variable)라고 부른다.
\end{definition}

수치형 속성 $X_j$는 별다른 변환 없이도 이미 확률변수다. 결과 $v$를
그 값 자신으로 보내는 \term{항등 확률변수}(Identity Random Variable)
\[
  X(v) = v, \qquad v \in \mathcal{O} = \mathbb{R}
\]
이기 때문이다.

\begin{example}[꽃받침 길이를 0/1로 바꾸기]
\label{ex:iris-long-sepal}
``꽃받침이 길다''를 ``꽃받침 길이가 7\,cm 이상이다''로 정하면, 확률변수
$A$를
\[
  A(v) = \begin{cases} 1 & v \ge 7 \\ 0 & v < 7 \end{cases}
\]
로 정의할 수 있다. $A$가 가지는 값은 $\{0,1\}$뿐이므로 $A$는 이산
확률변수다.
\end{example}

연속 확률변수의 예(삼각형 모양 밀도)는 \Cref{sec:pmf-pdf}의
\Cref{ex:triangle-pdf}에서 이어 다룬다. 한 절에서는 한 개념만 새로
들여온다.

표본공간을 적을 때 주의할 점이 있다. \Cref{def:random-variable}은
수치형 속성의 표본공간을 $\mathcal{O} = \mathbb{R}$이라고 적지만, 실제로
관측되는 꽃받침 길이는 $[4.3, 7.9]$ 구간에만 있다. $[4.3, 7.9]$는
\textbf{이번에 관측된 범위}이고, 표본공간은 정의상 \textbf{나올 수 있는
모든 결과}다. 관측 범위가 좁다고 표본공간 자체가 좁아지는 것은 아니다.

% =====================================================================
\section{확률질량함수와 확률밀도함수}
\label{sec:pmf-pdf}
% source: deck 4 slide 6-8, 15

\begin{inbrief}
  \item[Why] 확률변수의 각 값이 얼마나 자주 나오는지를 적는 함수가
             필요하다.
  \item[When] 이산 확률변수는 PMF, 연속 확률변수는 PDF를 쓴다.
  \item[Where] 주사위 눈, 붓꽃의 0/1 판정에는 PMF를, 꽃받침 길이 같은
               연속 측정값에는 PDF를 쓴다.
  \item[How] PMF는 값마다의 확률을 더하고, PDF는 구간에서 적분한다.
\end{inbrief}

\paragraph{고치는 문제.} 연속 확률변수에서는 한 점 $X=v$가 나올 확률이
정확히 0이다(\Cref{def:pdf} 뒤에서 다시 본다). 그러므로 ``이 값이 나올
확률이 얼마인가''라는 질문이 이산형에서는 성립해도 연속형에서는 성립하지
않는다. 그래서 이산형과 연속형은 서로 다른 함수로 분포를 적는다.

\begin{definition}[확률질량함수]
\label{def:pmf}
이산 확률변수 $X$의 \term{확률질량함수}(Probability Mass Function, PMF)
$f$는
\[
  f(x) = P(X = x)
\]
로 정의되고, 다음 두 성질을 만족한다. $f(x) \ge 0$이고
$\sum_x f(x) = 1$이다.
\end{definition}

\Cref{ex:iris-long-sepal}의 $A$로 PMF를 계산해 보자. 150송이 중
꽃받침 길이가 7\,cm 이상인 것은 13송이다.

\begin{align*}
  f(1) &= P(A=1) = 13/150 = 0.0867  && \text{길다고 판정된 비율}\\
  f(0) &= P(A=0) = 137/150 = 0.9133 && \text{나머지 비율}
\end{align*}

$f(0) + f(1) = 1$로 \Cref{def:pmf}의 성질을 만족한다. $A$처럼 값이
$\{0,1\}$뿐인 확률변수의 분포를 \term{베르누이 분포}(Bernoulli
Distribution)라고 부른다. 이름 있는 분포로서의 자세한 성질은
\Cref{sec:binomial}에서 다시 다룬다.

\begin{definition}[확률밀도함수]
\label{def:pdf}
연속 확률변수 $X$의 \term{확률밀도함수}(Probability Density Function,
PDF) $f$는 구간 확률을 다음과 같이 넓이로 정의한다.
\[
  P(X \in [a,b]) = \int_a^b f(x)\,dx, \qquad \int_{-\infty}^{\infty} f(x)\,dx = 1
\]
한 점의 확률은 폭이 0인 구간의 넓이이므로 항상 $P(X=v) = 0$이다.
\end{definition}

\begin{example}[삼각형 모양의 밀도]
\label{ex:triangle-pdf}
\[
  f(x) = \begin{cases}
    x     & 0 \le x \le 1 \\
    2 - x & 1 < x \le 2 \\
    0     & \text{그 밖}
  \end{cases}
\]
그래프로 그리면 $(0,0)$-$(1,1)$-$(2,0)$을 잇는 삼각형이다
(\Cref{fig:triangle-pdf-cdf}). 넓이가 1인지 먼저 확인한다.

\begin{align*}
  &\tfrac{1}{2}\times 1 \times 1 = 0.5 && \text{왼쪽 삼각형(밑변 1, 높이 1) 넓이}\\
  &\tfrac{1}{2}\times 1 \times 1 = 0.5 && \text{오른쪽 삼각형 넓이}\\
  &0.5 + 0.5 = 1                       && \text{전체 넓이, \Cref{def:pdf}의 조건과 일치}
\end{align*}

구간 확률도 넓이로 구한다. $P(0.5 \le X \le 1.5)$는 전체 넓이 1에서
양끝의 작은 삼각형 두 개를 뺀 값이다. 양끝 삼각형은 각각 밑변 0.5,
높이 0.5다.

\begin{align*}
  &2 \times \left(\tfrac{1}{2}\times 0.5 \times 0.5\right) = 0.25 && \text{양끝 두 삼각형의 넓이 합}\\
  &1 - 0.25 = 0.75                                                 && \text{$P(0.5 \le X \le 1.5)$}
\end{align*}
\end{example}

\begin{figure}[ht]\centering
  \includegraphics[width=\linewidth]{04_02_triangle_pdf_cdf.png}
  \caption{삼각형 밀도(PDF, 왼쪽)와 그 누적분포함수(CDF, 오른쪽). 색칠한
  구간 $[0.5,1.5]$의 넓이가 0.75이고, CDF에서는 $F(1)=0.5$인 지점이
  표시되어 있다.}
  \label{fig:triangle-pdf-cdf}
\end{figure}

\begin{pitfall}[PDF 값 자체는 확률이 아니다]
\label{pit:pdf-not-prob}
$f(x)$가 1을 넘어도 PDF의 정의를 어기지 않는다. 예를 들어 폭 0.5인
구간에 고르게 퍼진 밀도는 $f(x) = 2$다($0.5 \times 2 = 1$로 넓이가
1이어야 하기 때문이다). 확률은 언제나 \Cref{def:pdf}처럼 \textbf{넓이}로
읽는다. 높이 하나만 보고 확률이라고 읽으면 틀린다.
\end{pitfall}

\subsection{기댓값}
\label{subsec:expectation}

\term{기댓값}(Expected Value) $\E[X]$는 확률변수가 평균적으로 어떤 값을
갖는지를 나타낸다. 이산형은 PMF로 가중 평균을 내고, 연속형은 PDF로
적분한다.
\[
  \E[X] = \sum_x x f(x) \qquad \text{(이산)}, \qquad
  \E[X] = \int_{-\infty}^{\infty} x f(x)\,dx \qquad \text{(연속)}
\]

주사위에서는 여섯 눈이 각각 $1/6$의 확률을 가지므로

\begin{align*}
  \E[X] &= \sum_{k=1}^{6} k \cdot \tfrac{1}{6} = \tfrac{1+2+3+4+5+6}{6}
    && \text{PMF로 가중 평균}\\
  \E[X] &= \tfrac{21}{6} = 3.5 && \text{약분}
\end{align*}

\Cref{ex:triangle-pdf}의 삼각형 밀도에서는

\begin{align*}
  \E[X] &= \int_0^1 x \cdot x\,dx + \int_1^2 x(2-x)\,dx
    && \text{구간별로 나눠 적분}\\
  \E[X] &= \tfrac{1}{3} + \tfrac{2}{3} && \text{각 적분값}\\
  \E[X] &= 1 && \text{합}
\end{align*}

두 계산 모두 도형의 무게중심과 일치한다. 표본에서 이 기댓값을
추정한 값이 \Cref{def:center}의 평균 $\hat\mu$다.

% =====================================================================
\section{누적분포함수}
\label{sec:cdf-def}
% source: deck 4 slide 9-11, 18

\begin{inbrief}
  \item[Why] 이산과 연속을 한 함수로 다루고, ``이 값 이하가 나올
             확률''을 바로 읽기 위해서다.
  \item[When] 분포를 비교하거나, 분위수(Quantile)나 중앙값을 구할 때
              쓴다.
  \item[Where] 주사위, 붓꽃 데이터의 이항·정규 근사, 뒤에서 다룰
               분포 비교와 가설검정에 쓰인다.
  \item[How] PMF는 누적합으로, PDF는 누적적분으로 CDF를 만든다.
\end{inbrief}

\paragraph{고치는 문제.} PMF와 PDF는 이산형과 연속형에서 서로 다른
종류의 함수라서 같은 축 위에 겹쳐 놓고 비교할 수 없다. 이산형과
연속형을 같은 잣대로 비교하려면, 둘 다 $[0,1]$ 사이의 값을 내놓는
함수가 필요하다.

\begin{definition}[누적분포함수]
\label{def:cdf}
확률변수 $X$의 \term{누적분포함수}(Cumulative Distribution Function,
CDF) $F$는
\[
  F(x) = P(X \le x), \qquad F : \mathbb{R} \to [0,1]
\]
로 정의된다. 이산형에서는 $F(x) = \sum_{u \le x} f(u)$이고, 연속형에서는
$F(x) = \int_{-\infty}^{x} f(u)\,du$다. $F$는 네 가지 성질을 항상
만족한다. 감소하지 않는다, 오른쪽에서 연속이다(Right-continuous, 즉
$\lim_{h \to 0^+} F(x+h) = F(x)$), $F(x) \to 0$ ($x \to -\infty$),
$F(x) \to 1$ ($x \to \infty$).
\end{definition}

\begin{example}[주사위 눈의 CDF]
\label{ex:dice-cdf}
\Cref{tab:dice-cdf}은 주사위 눈 $k$마다 $P(X=k)$와 $F(k)=P(X\le k)$를
정리한다(\Cref{fig:dice-pmf-cdf}).

정수가 아닌 값에서도 $F$는 정의된다. 계단이 정수 사이에서는 평평하게
이어지므로 $F(3.5) = F(3) = 3/6 = 0.5$다.
\end{example}

\begin{table}[ht]\centering
\begin{tabular}{lrr}
\toprule
$k$ & $P(X=k)$ (확률) & $F(k)=P(X\le k)$ (누적확률) \\
\midrule
1 & 1/6 & 1/6 = 0.167 \\
2 & 1/6 & 2/6 = 0.333 \\
3 & 1/6 & 3/6 = 0.5 \\
4 & 1/6 & 4/6 = 0.667 \\
5 & 1/6 & 5/6 = 0.833 \\
6 & 1/6 & 6/6 = 1 \\
\midrule
$a < 1$ & 없음 & $F(a) = 0$ \\
$a > 6$ & 없음 & $F(a) = 1$ \\
\bottomrule
\end{tabular}
\caption{주사위 눈의 확률(Probability)과 누적확률(Cumulative Probability).}
\label{tab:dice-cdf}
\end{table}

\begin{figure}[ht]\centering
  \includegraphics[width=\linewidth]{04_01_dice_pmf_cdf.png}
  \caption{주사위 눈의 확률질량함수(PMF, 왼쪽)와 누적분포함수(CDF,
  오른쪽). CDF는 각 정수에서 계단으로 뛰어오르고, 뛰어오르기 직전의
  값은 열린 점, 뛰어오른 뒤의 값은 닫힌 점으로 표시된다.}
  \label{fig:dice-pmf-cdf}
\end{figure}

\begin{example}[F(3)과 P(X>3)]
\label{ex:cdf-fifteenths}
확률변수 $X$가 $P(X=x) = x/15$, $x \in \{1,2,3,4,5\}$를 만족한다고
하자. 먼저 이것이 PMF의 조건을 만족하는지 확인한다.

\begin{align*}
  &1+2+3+4+5 = 15               && \text{분자의 합}\\
  &15/15 = 1                    && \text{확률의 합, \Cref{def:pmf}의 조건과 일치}
\end{align*}

$F(3)$은 $x=1,2,3$의 확률을 누적한다.

\begin{align*}
  F(3) &= 1/15 + 2/15 + 3/15 = 6/15 = 0.4
\end{align*}

$P(X>3)$은 여집합으로 구한다.

\begin{align*}
  P(X>3) &= 1 - F(3) = 1 - 0.4 = 0.6
\end{align*}

교차 확인은 $x=4,5$의 확률을 직접 더해서 한다.

\begin{align*}
  P(4) + P(5) &= 4/15 + 5/15 = 9/15 = 0.6
\end{align*}

두 계산이 일치한다.
\end{example}

\begin{pitfall}[이산형에서는 $>$와 $\ge$가 다르다]
\label{pit:cdf-strict}
\Cref{ex:cdf-fifteenths}에서 $P(X>3) = 0.6$이지만 $P(X \ge 3)$은
다르다. $X=3$을 포함하기 때문이다.
\begin{align*}
  P(X \ge 3) &= 1 - F(2) && \text{여집합}\\
  &= 1 - \tfrac{1+2}{15} && \text{$F(2)$를 PMF로 전개}\\
  &= 1 - \tfrac{3}{15} = \tfrac{12}{15} = 0.8 && \text{계산}
\end{align*}
$0.8 \ne 0.6$이다. 이산형에서 경계값을 포함하는지는 결과를 바꾼다.
연속형에서는 $P(X=v)=0$이므로(\Cref{def:pdf}) $>$와 $\ge$의 차이가
사라진다.
\end{pitfall}

\begin{pitfall}[식의 좌우가 다른 변수를 쓰면 안 된다]
\label{pit:cdf-slide-kx}
슬라이드 값은 $P(X>k)=1-P(X\le x)$, 다시 계산하면
$P(X>x)=1-F(x)$, 즉 $P(X>3)=0.6$. 왼쪽 변에 쓴 $k$와 오른쪽 변에
쓴 $x$가 같은 변수를 가리켜야 식이 성립하는데, 슬라이드는 이 둘을
다른 기호로 적었다. 변수 하나를 통일해 $P(X>x) = 1 - F(x)$로 써야
\Cref{def:cdf}와 맞는 식이 된다. 숫자 결과 0.6 자체는 맞다.
\end{pitfall}

\paragraph{중앙값을 CDF로 읽기.} 중앙값의 정의 자리는
\Cref{def:center}에 하나만 둔다. 여기서는 CDF로 중앙값을 찾는 성질만
본다. 중앙값 $m$은
\[
  P(X \le m) \ge \tfrac{1}{2}, \qquad P(X \ge m) \ge \tfrac{1}{2}
\]
을 만족하는 값이다. 주사위에서는 $F(3) = 0.5$이고 $F(4) = 0.667$이므로
$3 \le m \le 4$인 모든 값이 이 조건을 만족하고, 관례로 가운데 값인
3.5를 쓴다. 정규분포(Normal Distribution)처럼 CDF가 매끈하게 증가하는
경우에는 $F(\mu) = 0.5$이므로 중앙값과 평균이 같다.
\Cref{fig:iris-cdfs}의 붓꽃 정규 근사 CDF에서 $(\mu, 0.5) = (5.843,
0.5)$인 점이 바로 이 성질을 보여 준다. 중앙값을 실제로 어떻게
계산하는지는 \Cref{sec:center}에서 이어진다.

붓꽃 데이터에서 서로 다른 두 확률변수의 이론 CDF를 비교할 수 있다
(\Cref{fig:iris-cdfs}). 하나는 꽃받침 길이 자체의 정규
CDF($\mu=5.843$, $\sigma^2=0.681$)이고, 다른 하나는 꽃받침 길이가
아니라 \Cref{ex:iris-long-sepal}의 0/1 판정을 $m=10$번 반복했을 때
``길다''로 판정되는 \textbf{횟수}의 이항(Binomial) CDF($p=0.087$,
$m=10$)다.

\begin{align*}
  F(0) &= 0.913^{10} = 0.4024 && \text{이항, $P(\text{10번 중 0번 ``길다''})$}\\
  F(1) &= 0.7859              && \text{이항, 누적}\\
  F(2) &= 0.9504               && \text{이항, 누적}\\
  F(3) &= 0.9922               && \text{이항, 누적}
\end{align*}

이항분포와 정규분포 각각의 자세한 정의와 모수는
\Cref{sec:binomial}과 \Cref{sec:normal}에서 다룬다.

\begin{figure}[ht]\centering
  \includegraphics[width=\linewidth]{04_03_iris_cdfs.png}
  \caption{붓꽃 데이터에서 나온 이론 CDF 두 개. 정규 CDF($\mu=5.843$,
  $\sigma^2=0.681$, 점 $(\mu, 0.5)$ 표시)는 꽃받침 길이 자체의 분포이고,
  이항 CDF($p=0.087$, $m=10$)는 꽃받침 길이가 아니라 ``길다''로
  판정되는 횟수의 분포다.}
  \label{fig:iris-cdfs}
\end{figure}

\begin{pitfall}[분산은 $n$과 $n-1$ 중 무엇으로 나눴는지가 다르다]
\label{pit:variance-n}
슬라이드 값은 0.681($n$으로 나눔), $n-1$로 다시 계산하면 0.686.
이것은 오류가 아니라 \textbf{관례 차이}다. $n$으로 나누면 모집단
분산의 정의를 그대로 쓴 것이고, $n-1$로 나누면 표본에서 모분산을
추정할 때 생기는 편향(Bias)을 보정한 것이다(\Cref{def:spread},
\Cref{pit:bias}). 이 책의 \Cref{def:spread}는 $n-1$ 쪽 정의를 쓴다.
\end{pitfall}

두 CDF를 겹쳐 놓고 최대 격차를 재는 방법은 \Cref{sec:cdf}에서, 데이터
자체로 계단 모양 CDF를 만드는 경험적 CDF(Empirical CDF)는
\Cref{def:ecdf}에서 다룬다.

% =====================================================================
\section{중심 --- 평균과 중앙값}
\label{sec:center}
% source: deck 4 slide 13, 15-18

\begin{inbrief}
  \item[Why] 데이터 전체를 대표하는 값 하나가 필요하기 때문이다.
  \item[When] \term{평균}(Mean)은 분포가 대칭일 때, \term{중앙값}(Median)은
              한쪽으로 치우쳤을 때 쓴다.
  \item[Where] 소득·가격·반응시간처럼 극단값이 존재하는 데이터에서는 거의 항상
               중앙값 쪽이 안전하다.
  \item[How] 왜도를 먼저 계산해 어느 쪽을 쓸지 정한다(\Cref{sec:skewness}).
\end{inbrief}

\begin{definition}[중심의 두 정의]
\label{def:center}
\term{평균}(Mean)은 모든 값을 더해서 개수로 나눈 값이다.
\[
  \mu = \frac{1}{n}\sum_{i=1}^{n} x_i
\]
\term{중앙값}(Median)은 값을 크기순으로 정렬했을 때 가운데에 오는 값이다.
\end{definition}

두 값의 차이는 계산식이 아니라 \textbf{극단값(Extreme Value)에 대한 반응}에
있다. 평균은 모든 값을 실제로 더하므로, 값 하나가 커지면 그 크기만큼 끌려간다.
중앙값은 ``몇 번째 값인가''만 보므로, 가장 큰 값이 두 배가 되어도 순위가 그대로면
꿈쩍하지 않는다.

이 성질을 \term{저항성}(Resistance)이라고 한다. 중앙값은 저항성이 있고 평균은
없다.

\begin{example}[값 하나가 평균을 움직인다]
\label{ex:resistance}
다섯 개의 값 $\{1,\ 2,\ 3,\ 4,\ 5\}$ 에서 평균과 중앙값은 모두 3이다.
마지막 값만 $5 \to 100$ 으로 바꾸면:

\begin{center}
\begin{tabular}{lrr}
\toprule
데이터 & 평균 & 중앙값 \\
\midrule
$\{1,2,3,4,5\}$   & 3.0  & 3 \\
$\{1,2,3,4,100\}$ & 22.0 & 3 \\
\bottomrule
\end{tabular}
\end{center}

평균은 7배 넘게 뛰었지만 중앙값은 그대로 3이다. 그리고 바뀐 평균 22는
\textbf{다섯 개 값 중 네 개보다 크다} --- 대표값으로서 아무 의미가 없다.
\end{example}

% =====================================================================
\section{퍼짐 --- 표준편차와 IQR}
\label{sec:spread}
% source: deck 4 slide 21-22

\begin{inbrief}
  \item[Why] 중심만으로는 부족하다. 값들이 중심 주변에 모여 있는지 흩어져 있는지를
             알아야 한다.
  \item[When] \term{표준편차}(Standard Deviation)는 대칭 분포에서, \term{사분위
              범위}(Interquartile Range)는 극단값이 있을 때 쓴다.
  \item[Where] 두 값의 차이가 클수록 그 데이터에 극단값이 많다는 신호다.
  \item[How] $\sd$ 와 $\IQR$ 을 \emph{둘 다} 계산해 비교한다. 차이 자체가 정보다.
\end{inbrief}

\begin{definition}[퍼짐의 두 정의]
\label{def:spread}
\term{표준편차}(Standard Deviation)는 각 값이 평균에서 떨어진 거리를 제곱해
평균낸 뒤 제곱근을 취한 값이다.
\[
  \sigma = \sqrt{\frac{1}{n-1}\sum_{i=1}^{n}(x_i - \mu)^2}
\]
\term{사분위 범위}(Interquartile Range)는 가운데 50\%가 차지하는 폭이다.
\[
  \IQR = Q_3 - Q_1
\]
\end{definition}

여기서도 대비는 \Cref{sec:center}와 같다. 거리를 \textbf{제곱}하기 때문에
표준편차는 극단값에 가장 민감하고, 상하위 25\%를 아예 보지 않는 $\IQR$ 은
극단값에 면역이다.

과제 데이터의 \texttt{sugars}(당류) 변수가 이 차이를 극단적으로 보여준다.

\begin{center}
\begin{tabularx}{0.85\linewidth}{@{}l r X@{}}
\toprule
 & 값 & 이 숫자가 실제로 가리키는 것 \\
\midrule
$\sd$  & 21.02 & 밀크셰이크 12개에 끌려 올라간 값 \\
$\IQR$ &  8.00 & 일반 메뉴가 실제로 차지하는 폭 \\
\bottomrule
\end{tabularx}
\end{center}

$\sd = 21.02$ 는 \textbf{실존하지 않는 폭}이다. 대부분의 메뉴는 3--11\,g 사이에
있고, 21이라는 숫자는 소수의 밀크셰이크가 만들어 낸 것이다.

\subsection{IQR이 이상치를 결정한다}
\label{subsec:fence}

박스플롯(Box Plot)에서 점으로 찍히는 \term{이상치}(Outlier)는 눈대중이 아니라
$\IQR$ 로부터 기계적으로 계산된다.

\begin{figure}[ht]\centering
\begin{tikzpicture}[
  node distance = 7mm,
  flowbox/.style = {draw = bookrule, fill = booksurface, rounded corners = 2pt,
                   minimum width = 34mm, minimum height = 9mm, align = center,
                   font = \small},
  arrow/.style  = {-{Stealth[length=2.2mm]}, draw = bookmuted}]
  \node[flowbox] (q)     {$Q_1 = 3$,\quad $Q_3 = 11$};
  \node[flowbox, below = of q]  (iqr)   {$\IQR = 11 - 3 = 8$};
  \node[flowbox, below = of iqr] (fence) {울타리(Fence)\\ $= 11 + 1.5 \times 8 = 23$};
  \node[flowbox, below = of fence, fill = bookwarnbg, draw = bookwarn] (out)
        {23을 넘는 12개\\ = 전부 밀크셰이크};
  \draw[arrow] (q)     -- (iqr);
  \draw[arrow] (iqr)   -- (fence);
  \draw[arrow] (fence) -- (out);
\end{tikzpicture}
\caption{이상치는 $\IQR$ 에서 기계적으로 유도된다. \texttt{sugars} 변수의 실제
계산 과정이다.}
\label{fig:fence-flow}
\end{figure}

\begin{figure}[ht]\centering
  \includegraphics[width=\linewidth]{q1_fig2_iqr_boxplot_anatomy.png}
  \caption{박스플롯의 해부. 상자가 곧 $\IQR$ 이고, 그 $\IQR$ 이
  \Cref{fig:fence-flow}의 울타리를 통해 어떤 점이 이상치인지를 결정한다. x축은
  의도적으로 끊었다 --- 실제 비율로 그리면 상자가 전체 폭의 1/5로 찌그러져 각
  부분에 이름을 붙일 공간이 없다.}
  \label{fig:iqr-anatomy}
\end{figure}

\begin{pitfall}[$1.5 \times \IQR$ 은 법칙이 아니다]
계수 1.5는 통계적으로 유도된 값이 아니라 Tukey가 제안한 \textbf{관례}다
\citep{tukey1977}. 따라서 ``울타리를 넘었으니 이상치다''가 아니라 "이 관례로는
이상치로 분류된다"가 정확한 표현이다. 실제로 \Cref{fig:iqr-anatomy}의 12개는
오류가 아니라 정상적인 밀크셰이크다. \textbf{이상치는 버릴 대상이 아니라 설명할
대상이다.}
\end{pitfall}

% =====================================================================
\section{최빈값과 범위}
\label{sec:mode-range}
% source: deck 4 slide 13, 19-20

\begin{inbrief}
  \item[Why] 평균과 중앙값 말고도, 범주형 데이터에 쓸 수 있는 중심값과
             가장 단순한 퍼짐값이 필요하다.
  \item[When] 값이 범주형이면 최빈값(Mode)을 쓰고, 계산이 빠른 1차
              점검이 필요하면 범위(Range)를 쓴다.
  \item[Where] 최빈값은 명목 척도(Nominal Scale)에도 정의된다
               (\Cref{tab:attr-ops}).
  \item[How] 최빈값은 가장 많이 나온 값을 세고, 범위는 최댓값에서
             최솟값을 뺀다.
\end{inbrief}

\begin{definition}[최빈값과 범위]
\label{def:mode}
확률변수의 \term{최빈값}(Mode)은 PMF 또는 PDF가 가장 큰 값이다. 표본에서는
관측 빈도 $\hat f(x)$가 가장 큰 값을 쓴다.
\[
  \operatorname{mode}(X) = \arg\max_x \hat f(x)
\]
\end{definition}

\begin{definition}[범위]
\label{def:range}
확률변수 $X$의 \term{범위}(Range) $r$은 최댓값과 최솟값의 차이다.
\[
  r = \max\{X\} - \min\{X\}
\]
\end{definition}

\begin{example}[붓꽃 꽃받침 길이의 최빈값과 범위]
\label{ex:mode-range}
꽃받침 길이는 소수 한 자리로 반올림되어 기록되어 있어, 같은 값이
여러 번 나온다. 가장 많이 나온 값은 5.0\,cm로 10송이가 이 값을
가진다(\Cref{fig:iris-mode-range}). 이 값은 \Cref{def:center}의 평균
5.843보다 작다. 범위는

\begin{align*}
  r &= 7.9 - 4.3 = 3.6 \qquad \text{(cm)}
\end{align*}
\end{example}

\begin{figure}[ht]\centering
  \includegraphics[width=\linewidth]{04_04_iris_mode_range.png}
  \caption{붓꽃 꽃받침 길이의 빈도. 최빈값(Mode) 5.0\,cm(10송이)을
  강조하고, 평균 5.843\,cm에 선을 그었으며, 최소 4.3\,cm과 최대
  7.9\,cm 사이의 범위(Range) 3.6\,cm을 화살표로 표시했다.}
  \label{fig:iris-mode-range}
\end{figure}

\Cref{ex:resistance}의 데이터를 다시 쓰면, 최빈값과 범위가 저항성에서
어떻게 갈리는지 보인다. $\{1,2,3,4,5\}$는 모든 값이 한 번씩만 나와
최빈값이 정의되지 않으므로, 값 하나를 중복시킨 $\{1,2,2,3,5\}$로
바꿔서 비교한다.

\begin{center}
\begin{tabular}{lrr}
\toprule
데이터 & 최빈값(Mode) & 범위(Range) \\
\midrule
$\{1,2,2,3,5\}$   & 2 & 4 \\
$\{1,2,2,3,100\}$ & 2 & 99 \\
\bottomrule
\end{tabular}
\end{center}

마지막 값을 $5 \to 100$으로 바꿔도 최빈값은 2로 그대로지만, 범위는
$4 \to 99$로 25배 가까이 뛴다.

\Cref{tab:robustness}은 이 장에서 다룬 여섯 가지 통계량을 저항성
기준으로 정리한다.

\begin{table}[ht]\centering
\begin{tabularx}{\linewidth}{@{}l l X@{}}
\toprule
통계량(Statistic) & 극단값 하나에 흔들리는가 & 근거 \\
\midrule
평균(Mean) & 흔들림 & \Cref{ex:resistance} \\
중앙값(Median) & 안 흔들림 & \Cref{ex:resistance} \\
최빈값(Mode) & 안 흔들림(단, 아래 함정 참고) & \Cref{ex:mode-range} \\
범위(Range) & 흔들림 & \Cref{ex:mode-range} \\
표준편차(Standard Deviation) & 흔들림 & \Cref{sec:spread} \\
사분위범위(IQR) & 안 흔들림 & \Cref{sec:spread} \\
\bottomrule
\end{tabularx}
\caption{여섯 가지 통계량의 저항성(Resistance) 성적표.}
\label{tab:robustness}
\end{table}

\begin{pitfall}[최빈값은 극단값이 아닌 값 하나에도 뒤집힌다]
\label{pit:mode-unstable}
최빈값이 저항성이 있다는 말은 \textbf{극단값}에 대해서만 참이다. 극단값이
\emph{아닌} 값 하나가 바뀌면 최빈값 자체가 뒤집힐 수 있다. $\{1,2,2,3,3,3\}$의
최빈값은 3인데, 극단값이 아닌 값 하나를 $3 \to 2$로 바꾼
$\{1,2,2,2,3,3\}$에서는 최빈값이 2로 바뀐다. 게다가 연속 측정값에서는
모든 값이 한 번씩만 나오면 최빈값이 아예 정의되지 않는다. 붓꽃
꽃받침 길이처럼 소수 한 자리로 반올림되어야 비로소 값이 겹치고
최빈값이 의미를 가진다.
\end{pitfall}

% =====================================================================
\section{비대칭 --- 왜도}
\label{sec:skewness}

\begin{inbrief}
  \item[Why] 평균을 쓸지 중앙값을 쓸지 \emph{객관적으로} 정하기 위해서다.
  \item[When] 요약통계표를 만들 때 항상 함께 계산한다.
  \item[Where] $|\gamma| > 1$ 이면 평균 대신 중앙값을 보고한다.
  \item[How] 부호는 꼬리의 방향, 크기는 치우친 정도를 뜻한다.
\end{inbrief}

\begin{definition}[왜도]
\label{def:skewness}
\term{왜도}(Skewness)는 분포의 비대칭 정도다.
\[
  \gamma = \frac{\E\!\left[(x-\mu)^3\right]}{\sigma^3}
\]
\end{definition}

핵심은 \textbf{세제곱}이다. 제곱이면 부호가 사라지므로 표준편차는 퍼진
\emph{정도}만 잰다. 세제곱은 부호가 살아남는다. 그래서 왜도는 "어느 쪽으로
퍼졌는가"를 잰다.

\begin{pitfall}[왜도(Skewness)는 편향(Bias)이 아니다]
\label{pit:bias}
두 단어가 일상어로는 모두 ``치우침''이지만, 통계학에서는 \textbf{완전히 다른
개념}이고 서로 바꿔 쓸 수 없다.

\begin{center}
\begin{tabularx}{\linewidth}{@{}l l X@{}}
\toprule
용어 & 무엇의 성질인가 & 정의 \\
\midrule
\term{왜도}(Skewness) & \textbf{분포}의 모양 & 분포가 좌우 대칭이 아닌 정도.
  $\gamma = \E[(x-\mu)^3]/\sigma^3$ \\
\term{편향}(Bias) & \textbf{추정량}의 정확도 & 추정값의 기댓값이 참값에서
  벗어난 크기. $\mathrm{Bias}(\hat\theta) = \E[\hat\theta] - \theta$ \\
\bottomrule
\end{tabularx}
\end{center}

왜도는 \emph{데이터가 어떻게 생겼는지}를 말하고, 편향은 \emph{내 추정 방법이
틀렸는지}를 말한다. 완벽하게 편향 없는 방법으로 추정해도 왜도가 큰 분포는
여전히 왜도가 크다. 반대로 완벽하게 대칭인 정규분포에서도 잘못된 추정량을 쓰면
편향이 생긴다.

\textbf{혼동이 실제로 일어나는 지점:} 표준편차의 $n-1$ 나눗셈이 바로 \emph{편향}
보정(Bias Correction)이다. $n$ 으로 나누면 분산이 과소추정되어 편향이 생긴다.
이것은 왜도와 아무 관련이 없다.

번역어로는 \term{왜도} 또는 \term{비대칭도}가 표준이며, \textbf{편향은 오역}이다.
\end{pitfall}

\subsection{실제로 계산하는 법 --- $g_1$ 과 $G_1$}
\label{sec:skew-compute}

\Cref{def:skewness}의 $\gamma$ 는 모집단(Population)의 정의다. 손에 있는 것은 표본이므로
실제로 계산하는 값은 두 단계를 거친다.

\textbf{1단계. 보정 없는 표본 왜도 $g_1$.} 정의식의 기댓값을 표본 평균으로 바꾼다.
$k$차 \term{중심적률}(Central Moment)을 $m_k = \frac{1}{n}\sum_{i=1}^{n}(x_i - \bar{x})^k$
라 하면
\[
  g_1 = \frac{m_3}{m_2^{3/2}} .
\]

\textbf{2단계. 표본 보정 $G_1$.} $g_1$ 은 표본에서 왜도를 과소추정하므로 계수를 곱한다.
\[
  G_1 = \frac{\sqrt{n(n-1)}}{n-2}\; g_1 .
\]
이 $G_1$ 이 \term{조정된 Fisher-Pearson 표준화 적률계수}(Adjusted Fisher-Pearson
Standardized Moment Coefficient)이고, \texttt{pandas} 의 \texttt{.skew()} 가 돌려주는 값이다.

손으로 구하는 절차는 다섯 줄이다.

\begin{enumerate}
  \item 평균 $\bar{x}$ 를 구한다.
  \item 편차의 제곱을 모두 더해 $n$ 으로 나눈다. 이것이 $m_2$ 다.
  \item 편차의 세제곱을 모두 더해 $n$ 으로 나눈다. 이것이 $m_3$ 다.
  \item $g_1 = m_3 / m_2^{1.5}$ 를 계산한다.
  \item $\sqrt{n(n-1)}/(n-2)$ 를 곱해 $G_1$ 을 얻는다.
\end{enumerate}

\begin{example}[다섯 개 값으로 끝까지 계산하기]
\label{ex:skew-by-hand}
데이터 $\{1,\ 2,\ 3,\ 4,\ 10\}$, $n = 5$ 이다.

\begin{center}
\begin{tabular}{lrrr}
\toprule
$x_i$ & $x_i - \bar{x}$ & $(x_i-\bar{x})^2$ & $(x_i-\bar{x})^3$ \\
\midrule
1  & $-3$ & 9  & $-27$ \\
2  & $-2$ & 4  & $-8$  \\
3  & $-1$ & 1  & $-1$  \\
4  & $0$  & 0  & $0$   \\
10 & $6$  & 36 & $216$ \\
\midrule
합 & $0$ & 50 & 180 \\
\bottomrule
\end{tabular}
\end{center}

평균은 $\bar{x} = 20/5 = 4$ 다. 이어서

\begin{align*}
  m_2 &= 50/5 = 10, \qquad m_3 = 180/5 = 36, \\
  g_1 &= \frac{36}{10^{1.5}} = \frac{36}{31.6228} = 1.1384, \\
  G_1 &= \frac{\sqrt{5 \cdot 4}}{3}\, g_1 = 1.4907 \times 1.1384 = 1.6971 .
\end{align*}

\texttt{pd.Series([1,2,3,4,10]).skew()} 가 정확히 $1.6971$ 을 돌려준다.

편차 세제곱 열을 보면 왜도가 어디서 오는지 한눈에 보인다. 값 10 하나가 $+216$ 을
만들고 나머지 넷이 합쳐서 $-36$ 밖에 못 만든다. \textbf{세제곱이 먼 값에 압도적인
가중치를 준다.}
\end{example}

보정 계수 $\sqrt{n(n-1)}/(n-2)$ 는 $n$ 이 커지면 1로 수렴한다. \Cref{ex:skew-by-hand}
처럼 $n=5$ 면 1.49로 크지만, 이 책의 데이터처럼 $n=126$ 이면 1.0121이라 $g_1$ 과 $G_1$ 의
차이는 1\% 남짓이다.

\begin{center}
\begin{tabular}{lr}
\toprule
\texttt{sugars} 의 왜도 계산 방식 & 값 \\
\midrule
$g_1$ (보정 전)                                  & 2.6765 \\
$G_1$ (보정 후)                                  & 2.7088 \\
$m_3 / s^3$ (분모만 $\mathrm{ddof}=1$)           & 2.6447 \\
\texttt{pandas .skew()}                          & 2.7088 \\
\bottomrule
\end{tabular}
\end{center}

\begin{pitfall}[분자와 분모의 나눗셈을 섞지 않는다]
위 표의 셋째 줄 $m_3/s^3$ 는 흔한 오답이다. 분모에는 표본표준편차($\mathrm{ddof}=1$,
$n-1$ 로 나눔)를 쓰면서 분자 $m_3$ 는 $n$ 으로 나눈 값이라, \textbf{어느 정의와도
일치하지 않는다}. $m_2$ 와 $m_3$ 는 둘 다 $n$ 으로 나누고, 보정은 마지막에 계수로
한 번만 한다.
\end{pitfall}

\subsection{부호와 방향}
\label{sec:skew-sign}

\begin{center}
\begin{tabularx}{\linewidth}{@{}l l l X@{}}
\toprule
부호 & 이름 & 모양 & 평균과 중앙값의 관계 \\
\midrule
$\gamma > 0$ & 우측 왜도(Right-skewed) & 오른쪽 꼬리가 김 & 평균 $>$ 중앙값 \\
$\gamma = 0$ & 대칭(Symmetric)         & 종 모양         & 평균 $=$ 중앙값 \\
$\gamma < 0$ & 좌측 왜도(Left-skewed)  & 왼쪽 꼬리가 김  & 평균 $<$ 중앙값 \\
\bottomrule
\end{tabularx}
\end{center}

왜 평균만 끌려가는지는 \Cref{ex:resistance}에서 이미 봤다. 그래서 왜도는
그림에서 \textbf{두 세로선 사이의 틈}으로 나타난다.

\begin{figure}[ht]\centering
  \includegraphics[width=\linewidth]{q1_fig1_skewness_mean_vs_median.png}
  \caption{왜도를 눈으로 보는 법. 왜도가 커질수록($0.24 \to 2.71$) 중앙값(검정
  실선)과 평균(주황 점선) 사이의 틈이 벌어진다.}
  \label{fig:skew-gap}
\end{figure}

\begin{pitfall}[틈은 반드시 SD 단위로 비교한다]
원래 단위로 재면 거짓말이 된다. \texttt{sodium}(나트륨)의 틈은 43.75\,mg이고
\texttt{sugars}(당류)의 틈은 6.29\,g이지만, 왜도는 0.42 대 2.71로 정반대다.
단위가 다르기 때문이다. $(\mu - \mathrm{median})/\sigma$ 로 표준화해야
0.08 대 0.30이 되어 왜도와 순서가 맞는다. \Cref{fig:skew-gap}의 각 패널에
적힌 \texttt{gap} 값이 이 표준화된 값이다.
\end{pitfall}

같은 표본을 강의 5강(Basic Statistics and Visualization)에서 다룬 다섯 가지 그림,
즉 히스토그램(Histogram, \Cref{sec:histogram}), 선 그래프(Line Chart, 밀도 곡선
Density Line, \Cref{sec:density-line}), 박스플롯(Box Plot, \Cref{sec:boxplot}),
바이올린 플롯(Violin Plot, \Cref{sec:violin}), 바코드 차트(Barcode Chart,
\Cref{sec:barcode})로 나란히 그리면, 왜도의 부호가 그림마다 어떤 흔적을 남기는지
한눈에 비교할 수 있다(\Cref{tab:skew-by-chart}). \Cref{fig:skew-right}부터
\Cref{fig:skew-left}까지 세 그림의 데이터는
\textbf{실제 강의 자료가 아니라 합성(synthetic) 표본}이다. 우측 왜도
표본은 고정 시드(seed)로 뽑은 감마(Gamma) 분포 난수이고, 좌측 왜도 표본은 그 값을
$60$에서 뺀 거울상($60-x$)이라 $G_1$ 의 부호가 정확히 뒤집히며, 대칭 표본은 평균과
표준편차를 맞춘 정규(Normal) 분포 난수다. 각 패널의 검정 실선은 중앙값(median),
주황 점선은 평균(mean)으로 \Cref{fig:skew-gap}과 같은 인코딩이다.

\begin{example}[세 표본에서 왜도의 부호를 손으로 확인하기]
\label{ex:skew-three-cases}
세 표본 모두 $n=300$ 이다. 순서는 (a) $G_1$ 의 부호를 본다, (b) 그 부호가
평균$-$중앙값의 부호로 이어지는지 확인한다, (c) 그 부호로부터 긴 꼬리가 어느
쪽인지 결정한다.

\textbf{우측 왜도 표본.} $G_1 = 1.48 > 0$ 이므로 평균이 중앙값보다 커야 한다.
\begin{align*}
  &16.32 - 13.95 = 2.37       && \text{평균 $-$ 중앙값 (양수, 예상과 일치)}\\
  &2.37 / 11.15 = 0.21        && \text{SD 단위로 표준화한 틈}\\
  &22.06 - 13.95 = 8.11       && \text{$Q_3 - $ 중앙값}\\
  &13.95 - 8.40 = 5.55        && \text{중앙값 $- Q_1$}
\end{align*}
위쪽 거리(8.11)가 아래쪽 거리(5.55)보다 커서, 중앙값 위로 값이 더 퍼져 있다.
즉 긴 꼬리는 \textbf{오른쪽}이다. $G_1$ 의 부호, 평균$-$중앙값의 부호, 꼬리의
방향이 셋 다 같은 이야기를 하고 있다.

\textbf{좌측 왜도 표본.} $G_1 = -1.48 < 0$ 이므로 평균이 중앙값보다 작아야 한다.
\begin{align*}
  &43.68 - 46.05 = -2.37      && \text{평균 $-$ 중앙값 (음수, 예상과 일치)}\\
  &-2.37 / 11.15 = -0.21      && \text{SD 단위로 표준화한 틈}\\
  &51.60 - 46.05 = 5.55       && \text{$Q_3 - $ 중앙값}\\
  &46.05 - 37.94 = 8.11       && \text{중앙값 $- Q_1$}
\end{align*}
이번에는 아래쪽 거리(8.11)가 위쪽 거리(5.55)보다 커서, 긴 꼬리는
\textbf{왼쪽}이다. 좌측 왜도 표본은 우측 왜도 표본을 $60-x$ 로 뒤집어 만들었으므로,
틈($0.21 \leftrightarrow -0.21$)과 두 분위수 거리($8.11 \leftrightarrow 5.55$)가
정확히 자리를 바꾼다. 이것이 \Cref{fig:skew-left}의 표본이 \Cref{fig:skew-right}의 표본을
좌우로 뒤집은 모습으로 보이는 이유다.

\textbf{대칭 표본.} $G_1 = -0.07$ 로 부호는 음수지만 크기가 매우 작다.
\begin{align*}
  &16.48 - 16.78 = -0.30      && \text{평균 $-$ 중앙값}\\
  &-0.30 / 10.95 = -0.03      && \text{SD 단위로 표준화한 틈}\\
  &23.49 - 16.78 = 6.71       && \text{$Q_3 - $ 중앙값}\\
  &16.78 - 9.55 = 7.23        && \text{중앙값 $- Q_1$}
\end{align*}
두 분위수 거리(6.71과 7.23)가 서로 가까워서 어느 쪽으로도 뚜렷한 꼬리가
없다. $G_1$ 이 정확히 $0$ 이 아닌 이유는 유한한 표본이 완벽하게 대칭일 수는
없기 때문이다. $|G_1|=0.07 < 0.5$ 이므로 \Cref{sec:skew-rule}의 기준으로도
``거의 대칭''으로 분류된다.
\end{example}

\begin{figure}[p]\centering
  \includegraphics[width=\linewidth,height=0.85\textheight,keepaspectratio]{02_skew_right.png}
  \caption{우측 왜도(Right-skewed, $G_1=1.48$) 표본을 다섯 가지 그림으로 본 모습.
  히스토그램(Histogram)과 선 그래프(Line Chart, 밀도 곡선 Density Line)는 왼쪽에
  봉우리가 몰리고 오른쪽으로 긴 꼬리가 남으며, 박스플롯(Box Plot)은
  중앙값(Median, 검정 실선)이 상자 왼쪽으로 치우치고 오른쪽에
  이상치(Outlier, 주황 점)가 늘어선다. 바이올린 플롯(Violin Plot)은 오른쪽으로
  갈수록 가늘어지고, 바코드 차트(Barcode Chart)는 막대가 왼쪽에 몰렸다가
  오른쪽으로 갈수록 성겨진다.}
  \label{fig:skew-right}
\end{figure}

\begin{figure}[p]\centering
  \includegraphics[width=\linewidth,height=0.85\textheight,keepaspectratio]{02_skew_symmetric.png}
  \caption{대칭(Symmetric, $G_1=-0.07$) 표본을 다섯 가지 그림으로 본 모습. 중앙값과
  평균이 거의 겹치고, 박스플롯의 두 수염과 바이올린 플롯의 양쪽 폭, 바코드 차트의
  양쪽 밀도가 서로 비슷하다.}
  \label{fig:skew-symmetric}
\end{figure}

\begin{figure}[p]\centering
  \includegraphics[width=\linewidth,height=0.85\textheight,keepaspectratio]{02_skew_left.png}
  \caption{좌측 왜도(Left-skewed, $G_1=-1.48$) 표본을 다섯 가지 그림으로 본 모습.
  \Cref{fig:skew-right}를 좌우로 뒤집은 모양과 같다. 히스토그램(Histogram)과
  선 그래프(Line Chart, 밀도 곡선 Density Line)는 오른쪽에 봉우리가 몰리고
  왼쪽으로 긴 꼬리가 남으며, 박스플롯(Box Plot)은 중앙값(Median)이 상자
  오른쪽으로 치우치고 왼쪽에 이상치(Outlier)가 늘어선다.}
  \label{fig:skew-left}
\end{figure}

\begin{table}[htbp]\centering
\begin{tabularx}{\linewidth}{@{}l X X X@{}}
\toprule
그림(Chart) & 우측 왜도(Right-skewed) & 대칭(Symmetric) & 좌측 왜도(Left-skewed) \\
\midrule
히스토그램(Histogram) & 봉우리가 왼쪽, 막대가 오른쪽으로 길게 늘어짐
  & 봉우리가 가운데, 양쪽 막대 길이가 비슷함
  & 봉우리가 오른쪽, 막대가 왼쪽으로 길게 늘어짐 \\
선 그래프(Line Chart, 밀도 곡선 Density Line) & 봉우리가 평균 왼쪽, 오른쪽 꼬리가 김
  & 봉우리가 평균과 거의 겹침
  & 봉우리가 평균 오른쪽, 왼쪽 꼬리가 김 \\
박스플롯(Box Plot) & 중앙값 선이 상자 왼쪽에 치우침, 오른쪽 수염이 길고
  오른쪽에 이상치
  & 중앙값 선이 상자 가운데 근처, 양쪽 수염 길이가 비슷함
  & 중앙값 선이 상자 오른쪽에 치우침, 왼쪽 수염이 길고 왼쪽에 이상치 \\
바이올린 플롯(Violin Plot) & 오른쪽으로 가늘고 긴 꼬리
  & 좌우가 거의 대칭인 모양
  & 왼쪽으로 가늘고 긴 꼬리 \\
바코드 차트(Barcode Chart) & 막대가 왼쪽에 몰리고 오른쪽으로 갈수록 성겨짐
  & 막대가 가운데를 중심으로 고르게 퍼짐
  & 막대가 오른쪽에 몰리고 왼쪽으로 갈수록 성겨짐 \\
\bottomrule
\end{tabularx}
\caption{다섯 가지 그림(Chart)에서 왜도(Skewness)가 남기는 흔적.}
\label{tab:skew-by-chart}
\end{table}

\begin{pitfall}[박스플롯에는 평균이 기본으로 그려지지 않는다]
박스플롯은 원래 $Q_1$과 $Q_3$ 사이의 상자, 중앙값 선, 그리고 $1.5\times\IQR$
이내의 가장 먼 관측값까지 뻗는 수염만 그리며(\Cref{sec:boxplot}), 그 밖의 값은
점으로 찍히는 이상치다. 평균은 기본값으로 표시되지 않는다. \Cref{fig:skew-right}와
\Cref{fig:skew-left}에서는 설명을 위해 평균 점선을 일부러 덧그린 것이다. 평균 선이
없는 보통의 박스플롯에서 왜도를 읽으려면, 중앙값 선이 상자의 가운데에서 어느 쪽으로
치우쳤는지, 두 수염의 길이 차이, 그리고 이상치가 어느 쪽에 몰려 있는지를 본다.
\end{pitfall}

\subsection{판단 기준}
\label{sec:skew-rule}

\begin{center}
\begin{tabular}{@{}ll@{}}
\toprule
$|\gamma| < 0.5$        & 거의 대칭 --- 평균 + 표준편차 \\
$0.5 \le |\gamma| \le 1$ & 중간 --- 평균도 쓸 만하나 주의 \\
$|\gamma| > 1$          & 강한 왜도 --- \textbf{중앙값 + $\IQR$} \\
\bottomrule
\end{tabular}
\end{center}

\texttt{sugars} 는 $\gamma = 2.71$ 이고, 결정적으로 \textbf{평균 13.29가
$Q_3 = 11$ 보다 크다}. 즉 전체의 82.5\%가 ``평균 이하''다. 평균이라는 단어의
일상적 의미가 깨진 상태다(\Cref{ex:resistance}의 22와 같은 상황이다).

\begin{pitfall}[왜도의 정의는 하나가 아니다]
과목에 따라 \term{Pearson 제2 왜도계수}(Pearson's Second Skewness Coefficient)를
가르치기도 한다.
\[
  \mathrm{Sk}_2 = \frac{3(\bar{x} - \mathrm{median})}{s}
\]
같은 데이터에서 이 값은 \texttt{serving\_size} 0.19, \texttt{calories} 0.21,
\texttt{sodium} 0.25, \texttt{sugars} 0.90 이다. 부호는 $G_1$ 과 같지만 크기가 전혀
다르다. \texttt{sugars} 가 2.71이 아니라 0.90으로 나오므로, 위의 ``$|\gamma| > 1$ 이면
중앙값'' 기준을 그대로 적용하면 \textbf{판정이 뒤집힌다}. 두 지표는 스케일이 다르니
임계값을 섞어 쓰면 안 된다.
\end{pitfall}

\begin{aside}
R로 교차검증한다면 패키지마다 정의가 다르다는 점을 예상해야 한다.
\texttt{e1071::skewness(type = 2)} 가 \texttt{pandas} 와 같은 $G_1$ 이고,
\texttt{moments::skewness()} 는 보정 없는 $g_1$ 을 돌려준다. \texttt{sugars} 에서
2.7088과 2.6765의 차이로 나타난다.
\end{aside}

\begin{pitfall}[왜도는 모양을 알려 주지 않는다]
$\gamma = 2.71$ 만 보면 ``완만하게 늘어진 긴 꼬리''를 상상하게 된다. 실제로는
0--20\,g에 114개, 56--93\,g에 12개가 있고 \textbf{그 사이 20--56\,g은 완전히
비어 있다}. 왜도는 ``평균을 쓸까 중앙값을 쓸까''를 정해 주는 스위치일 뿐,
분포의 모양은 반드시 그림으로 확인해야 한다 --- \analogy{지도}로는 보이지 않고
\analogy{항공사진}이 필요한 지점이다.
\end{pitfall}

\begin{figure}[ht]\centering
  \includegraphics[width=\linewidth]{q1_fig3_skew_cause_two_populations.png}
  \caption{왜도 2.71의 정체. 긴 꼬리가 아니라 구멍을 사이에 둔 두 집단
  (Mixture of Two Populations)이다.}
  \label{fig:two-populations}
\end{figure}

% =====================================================================
\section{결측값 --- 0은 결측이 아니다}
\label{sec:missing}

\begin{inbrief}
  \item[Why] 0과 결측을 섞으면 계산은 조용히 성공하고 결과만 틀린다.
  \item[When] 데이터를 처음 받아 \texttt{isna()} 를 돌린 직후에 판단한다.
  \item[Where] 센서 데이터, 설문 응답, 의료 기록처럼 0이 ``없음''의 코드로 쓰이는
               분야에서 특히 위험하다.
  \item[How] ``이 변수에서 0이 \emph{물리적으로 가능한 값인가}''를 묻는다.
\end{inbrief}

\term{결측값}(Missing Value)은 \textbf{값이 기록되지 않았다}는 뜻이다. 0은
\textbf{측정했더니 0이었다}는 뜻이다. 둘은 완전히 다른 정보이고, 섞으면 복구할
수 없다.

\begin{example}[0을 결측으로 처리하면]
\label{ex:zero-as-missing}
과제 데이터에서 \texttt{sugars == 0} 인 항목은 17개이고, 전부 감자튀김(10개)과
치킨너겟(7개)이다. 이들은 실제로 당류가 없는 음식이다. 이를 결측으로 바꾸면:

\begin{center}
\begin{tabular}{@{}lrrrr@{}}
\toprule
 & $n$ & 평균 & 중앙값 & $\IQR$ \\
\midrule
0을 값으로 (옳음)   & 126 & 13.29 & 7.0 & 8.0 \\
0을 결측으로 (틀림) & 109 & 15.36 & 8.0 & 6.0 \\
\bottomrule
\end{tabular}
\end{center}

감자튀김의 당류 중앙값이 0\,g에서 1.15\,g으로 바뀐다. \textbf{당류가 없는
음식이 당류를 가진 것처럼 보이게 된다.} 게다가 $n$ 이 126에서 109로 줄어
``결측 없는 데이터''라는 사실 자체가 사라진다.
\end{example}

\begin{keyidea}
판단 기준은 하나다. \textbf{그 0이 물리적으로 가능한 값인가?}

\begin{itemize}
  \item 당류 0\,g, 강수량 0\,mm, 사고 건수 0건 $\rightarrow$ \textbf{가능하다. 값이다.}
  \item 혈압 0, 키 0\,cm, 성인 나이 0세 $\rightarrow$ \textbf{불가능하다.} 결측을
        0으로 채워 넣은 흔적일 가능성이 높다.
\end{itemize}

후자의 경우에만 \verb|df.replace(0, np.nan)| 이 정당하다.
\end{keyidea}

% =====================================================================
\section{어떻게 적용하는가}
\label{sec:desc-apply}

이 장의 도구는 과제와 다른 장에서 바로 쓰인다.

\begin{itemize}
  \item HW1의 버거 대 밀크셰이크 비교(\Cref{tab:hw1-burger-vs-shake})는
        이 장의 평균과 퍼짐(표준편차, 최솟값, 최댓값)으로 두 음식 유형을 비교한 결과다.
        중앙값과 $\IQR$은 HW1 표에 나오지 않는다.
  \item HW2의 추정량(\Cref{sec:hw2-estimators})은
        \Cref{subsec:expectation}의 기댓값을 표본에서 추정하는
        방법을 다룬다. HW2의 분포 비교(\Cref{sec:hw2-compare-how})는
        이 장의 CDF를 겹쳐 보는 방식을 쓴다.
  \item 경험적 CDF로 후보 분포를 눈으로 확인하는 절차는
        \Cref{sec:dist-check-cdf}에서 이어진다.
  \item \Cref{ex:iris-long-sepal}의 0/1 확률변수 $A$와 같은 발상은
        이산화(Discretization, \Cref{sec:discretization})로 일반화된다.
  \item 결측값을 최빈값으로 채우는 방법은 \Cref{sec:impute}에서 다룬다.
  \item 히스토그램(\Cref{sec:histogram})은 표본의 경험적 PDF고,
        분위수(\Cref{sec:quantile})는 CDF의 역함수다.
\end{itemize}

파이썬 \texttt{pandas}로는 이 절의 값들을 각각 \texttt{.mode()},
\texttt{.max() - .min()}, \texttt{(x >= 7).mean()}으로 구한다. 마지막
식의 결과는 $13/150 = 0.0867$로 \Cref{sec:pmf-pdf}의 $f(1)$과 같다.

% =====================================================================
\section{앞으로}
\label{sec:desc-future}

이 장에서 다룬 확률변수, PMF/PDF, CDF는 이름 있는 분포를 다루는 다음
장으로 이어진다(\Cref{ch:distributions}). 특히 이항분포와 정규분포는
\Cref{sec:binomial}과 \Cref{sec:normal}에서 모수를 추정하고 적합도를
검정하는 법까지 다룬다.

두 CDF를 겹쳐 최대 격차를 재는 방법(\Cref{sec:cdf})은 콜모고로프--스미르노프
검정(Kolmogorov--Smirnov Test, \Cref{sec:cd-ks})으로 정식화된다.

표본 평균이 모평균을 얼마나 잘 맞히는지, 즉 \term{불편추정량}(Unbiased
Estimator)이라는 성질이 실제로 무엇을 보장하는지는 가설검정
(\Cref{ch:hypothesis})에서 이어진다.

% =====================================================================
\section{요약}

\begin{itemize}
  \item 평균과 표준편차는 극단값에 끌려가고, 중앙값과 $\IQR$ 은 저항한다
        (\Cref{def:center}, \Cref{def:spread}).
  \item 왜도가 어느 쪽을 쓸지 결정한다. $|\gamma| > 1$ 이면 중앙값이다.
  \item 왜도(Skewness)와 편향(Bias)은 다른 개념이다(\Cref{pit:bias}).
  \item 왜도는 ``얼마나 치우쳤나''만 알려 준다. \emph{왜} 치우쳤는지는 그림의 몫이다
        (\Cref{fig:two-populations}).
  \item 0과 결측은 다르다. 0이 물리적으로 가능한지를 먼저 묻는다
        (\Cref{ex:zero-as-missing}).
\end{itemize}

\chapter{시각화 --- 요약이 숨긴 것을 보기}
\label{ch:visualization}
% source: deck 5 slide 1

\begin{inbrief}
  \item[Why] 평균이나 표준편차 같은 \term{요약통계}(Summary Statistics) 하나만으로는
             분포의 \emph{모양}(봉우리(Peak) 개수, 치우침(Skewness),
             이상치(Outlier))이 지워진다. 시각화는 그 지워진 정보를 되찾는 도구다.
  \item[When] 변수 하나의 분포를 처음 확인할 때, 또는 여러 집단·변수의 분포를
              \emph{비교}할 때.
  \item[Where] 이 책 전체의 탐색 단계에 쓰인다. 다만 형태마다 전제가 다르다
               (예: 박스플롯(Box Plot)은 이봉분포(Bimodal Distribution)에 쓰면
               안 된다).
  \item[How] 먼저 ``이 그림으로 무엇을 묻는가''를 정한다 $\to$ 그 질문에 맞는
             형태를 고른다(\Cref{tab:form-choice}) $\to$ 형태마다의 함정
             (구간 개수, 에러바(Error Bar)가 무엇을 뜻하는지, 표본 크기)을
             확인한다.
\end{inbrief}

\Cref{ch:descriptive}은 \analogy{지도}가 반드시 무언가를 버린다는 결론으로
끝났다. 이 장은 그 버려진 것을 되찾는 \analogy{항공사진}을 다룬다.

강의의 표현을 그대로 옮기면 이렇다. \emph{``통계값은 숫자이고 유용하지만, 이야기
전체를 들려주지는 않는다.''} 시각화는 한 변수의 분포를 이해할 때 유용하고,
여러 변수의 분포를 \textbf{비교}할 때는 거의 필수다.

\begin{keyidea}
차트는 장식이 아니라 \textbf{질문에 답하는 도구}다. 먼저 ``이 그림으로 무엇을
묻는가''를 정하고, 그 다음에 형태를 고른다. 순서를 뒤집으면 예쁘지만 아무것도
답하지 않는 그림이 나온다.
\end{keyidea}

% =====================================================================
\section{배경}
\label{sec:vis-background}
% source: deck 5 "Basic Statistics and Visualization" slide 1

이 장은 강의 노트 \texttt{5\_StatandVisualization.pptx}(``Basic Statistics and
Visualization'')를 옮긴다. 이 장을 읽는 데 필요한 개념을 먼저 깐다. 이미
다룬 것은 한 줄 요약과 링크로, 이 장 안에서 처음 정의하는 것은 뒤의 각 절에서
다룬다.

\begin{itemize}
  \item \term{평균}(Mean)과 \term{표준편차}(Standard Deviation, 퍼짐의 크기)는
        \Cref{sec:center,sec:spread}에서 다뤘다. \Cref{sec:ci}의 신뢰구간
        공식 $\bar{x} \pm z \cdot s/\sqrt{n}$이 이 둘 위에 선다.
  \item \term{확률밀도함수}(Probability Density Function, PDF)는
        \Cref{sec:pmf-pdf}에서 다뤘다. \Cref{sec:histogram}의 히스토그램은
        이 PDF를 데이터로부터 경험적으로 근사한 것이다.
  \item \term{정규분포}(Normal Distribution)와 그 표준편차 배수(예:
        $z=1.96$)는 \Cref{sec:normal}에서 다룬다. \Cref{sec:ci}의 공식에
        쓰이는 $z$ 값이 여기서 나온다.
\end{itemize}

% =====================================================================
\section{분위수와 백분위수}
\label{sec:quantile}

\begin{inbrief}
  \item[Why] 평균 하나로는 부족하고, 분포의 \emph{여러 지점}을 알아야 하기 때문이다.
  \item[When] 박스플롯, $\IQR$, 신뢰구간 --- 이 장의 거의 모든 도구가 분위수 위에 선다.
  \item[Where] 왜도가 있는 데이터에서 특히 중요하다. 평균은 무너져도 분위수는 버틴다.
  \item[How] 정렬 $\rightarrow$ 각 값에 $0$에서 $1$까지의 위치를 매김 $\rightarrow$
             원하는 위치에서 보간(Interpolation).
\end{inbrief}

\begin{definition}[분위수와 백분위수]
\label{def:quantile}
\term{분위수}(Quantile)는 데이터를 \textbf{같은 크기의 그룹으로 나누는 경계선}이다.
\term{백분위수}(Percentile)는 그중 데이터를 \textbf{100개} 그룹으로 나누는 분위수다.

$k$번째 백분위수는 값들을 나누어, $k\%$ 가 그 아래에 놓이고 $(100-k)\%$ 가 그
위에 놓이게 하는 값이다.
\end{definition}

\begin{aside}[둘의 관계]
백분위수는 분위수의 \textbf{부분집합}이다. 분위수 \contains{} 백분위수.
같은 지점을 부르는 이름이 둘인 셈이라 혼용되지만, 엄밀히는 \emph{눈금의 개수}가
다르다 --- 사분위수(Quartile)는 4등분, 십분위수(Decile)는 10등분, 백분위수는
100등분이다. 소프트웨어에서는 보통 $0 \sim 1$ 을 받는 \texttt{quantile()} 하나로
통일되어 있다.
\end{aside}

\begin{center}
\begin{tabularx}{\linewidth}{@{}l l l X@{}}
\toprule
이름 & 몇 등분 & 표기 & 비고 \\
\midrule
사분위수(Quartile) & 4 & $Q_1, Q_2, Q_3$ & $Q_1$ = 하사분위수(Lower Quartile) \\
십분위수(Decile) & 10 & $D_1 \dots D_9$ & \\
백분위수(Percentile) & 100 & $P_1 \dots P_{99}$ & $P_{50} = Q_2 = $ 중앙값 \\
\bottomrule
\end{tabularx}
\end{center}

\begin{pitfall}[0\% 분위수는 왜 최솟값인가]
강의가 던지는 질문이다. $0\%$ 분위수(0번째 백분위수)는 \textbf{최솟값}이다.
정의상 ``그 값보다 작은 값이 $0\%$''여야 하는데, 최솟값보다 작은 값은 데이터에
하나도 없기 때문이다. 같은 이유로 $100\%$ 분위수는 최댓값이다.
\end{pitfall}

\subsection{직접 계산해 보기}
\label{subsec:quantile-calc}

\begin{example}[25\% 분위수 구하기]
\label{ex:quantile}
강의 예제다. 여섯 개의 값이 있다: $3.7,\ 2.7,\ 3.3,\ 1.3,\ 2.2,\ 3.1$.

\textbf{1단계 --- 정렬한다.}
\[
  1.3,\quad 2.2,\quad 2.7,\quad 3.1,\quad 3.3,\quad 3.7
\]

\textbf{2단계 --- 각 값에 $0$부터 $1$까지 등간격의 위치를 매긴다.}
$n = 6$ 이므로 $i$번째 값의 위치는 $(i-1)/(n-1) = (i-1)/5$ 다.

\begin{center}
\begin{tabular}{@{}lrrrrrr@{}}
\toprule
$i$ & 1 & 2 & 3 & 4 & 5 & 6 \\
값 & 1.3 & 2.2 & 2.7 & 3.1 & 3.3 & 3.7 \\
위치 & 0.0 & 0.2 & 0.4 & 0.6 & 0.8 & 1.0 \\
\bottomrule
\end{tabular}
\end{center}

\textbf{3단계 --- 0.25는 눈금 사이에 있으므로 보간한다.}
$0.25$ 는 $0.2$(값 2.2)와 $0.4$(값 2.7) 사이에 있다. $0.2$ 에서
$0.05$ 만큼 갔으므로 전체 간격 $0.2$ 의 $1/4$ 지점이다.
\[
  0.75 \times 2.2 + 0.25 \times 2.7 = 1.65 + 0.675 = \mathbf{2.325}
\]

즉 하사분위수는 2.325다. \textbf{데이터에 없는 값}이라는 점에 주의하자 ---
분위수는 관측값 중 하나일 필요가 없다.
\end{example}

\begin{lstlisting}[language=Python]
import pandas as pd
s = pd.Series([3.3, 2.7, 2.2, 1.3, 3.7, 3.1])
s.quantile(0.25)    # 2.325  <- 위 손계산과 일치
\end{lstlisting}

\begin{figure}[ht]\centering
  \includegraphics[width=0.85\linewidth]{05_01_quantiles.png}
  \caption{정렬된 값에 위치를 매기고 원하는 지점을 읽는 과정. 곡선의 가로축이
  \Cref{ex:quantile}의 ``위치'', 세로축이 그 위치의 값이다.}
  \label{fig:quantile-curve}
\end{figure}

% =====================================================================
\section{히스토그램}
\label{sec:histogram}

\begin{inbrief}
  \item[Why] 확률밀도함수(Probability Density Function)를 데이터로부터 경험적으로
             그려 보기 위해서다.
  \item[When] 변수 하나의 분포 모양을 처음 확인할 때. 가장 먼저 그리는 그림이다.
  \item[Where] 연속형 변수에 쓴다. 범주형에는 막대그래프(Bar Chart)를 쓴다.
  \item[How] 데이터를 구간(Bin)으로 나누고, x축에 구간 중심을, y축에 그 구간의
             개수(또는 비율)를 놓는다.
\end{inbrief}

\term{히스토그램}(Histogram)은 확률밀도함수의 \textbf{경험적 표현}(Empirical
Representation)이다. 만드는 법은 단순하다.

\begin{enumerate}
  \item 데이터를 구간(Bin)으로 나눈다.
  \item x축 = 구간의 중심점.
  \item y축 = 그 구간에 들어간 데이터 개수(또는 전체 대비 비율).
\end{enumerate}

\subsection{구간 개수가 결론을 바꾼다}
\label{subsec:binning}

강의는 ``구간을 어떻게 나눌 것인가?''를 괄호 안에 던져 놓고 지나간다. 그런데
이것이 히스토그램의 유일한 자유도이자 가장 큰 함정이다.

\begin{figure}[ht]\centering
  \includegraphics[width=\linewidth]{05_02_histogram_bins.png}
  \caption{같은 데이터, 다른 구간 개수. 구간이 너무 적으면 구조가 뭉개지고, 너무
  많으면 표본 잡음(Sampling Noise)이 봉우리처럼 보인다.}
  \label{fig:bins}
\end{figure}

\begin{pitfall}[구간 개수는 보고해야 할 선택이다]
구간 개수를 바꾸면 봉우리가 생겼다 사라진다. 즉 \textbf{같은 데이터에서 다른
결론}이 나온다. 따라서 히스토그램에서 읽은 구조(``봉우리가 두 개다'')를
주장할 때는 다른 구간 개수에서도 그 구조가 살아남는지 확인해야 한다.
\Cref{fig:two-populations}의 두 집단은 구간을 어떻게 잡아도 사라지지 않기
때문에 믿을 수 있는 구조다.
\end{pitfall}

\subsection{여러 분포를 나란히 놓기}
\label{subsec:small-multiples}

히스토그램 여러 개를 같은 축으로 나란히 배치하면 집단 간 차이가 드러난다. 강의는
서브레딧별 순위 분포를 예로 든다.

\begin{figure}[ht]\centering
  \includegraphics[width=\linewidth]{05_03_small_multiples.png}
  \caption{작은 배수(Small Multiples). 축을 공유해야 비교가 성립한다 --- 패널마다
  축 범위가 다르면 눈은 그 차이를 읽지 못한다.}
  \label{fig:small-multiples}
\end{figure}

\begin{pitfall}[표본 크기를 같이 적는다]
\Cref{fig:small-multiples}에서 버거는 69개, 나머지는 11--12개다. 개수가 적은
패널은 모양이 들쭉날쭉한데, 이것은 분포의 성질이 아니라 \textbf{표본이 작아서}
생긴 것이다. 각 패널에 $n$ 을 적어 두지 않으면 독자가 잡음을 구조로 읽는다.
\end{pitfall}

% =====================================================================
\section{선 그래프 --- 밀도에 더 가까이}
\label{sec:density-line}

\begin{inbrief}
  \item[Why] 히스토그램의 계단을 매끄럽게 만들어 확률밀도함수에 더 가까이 가기 위해서다.
  \item[When] 두 개 이상의 분포를 \emph{겹쳐서} 비교할 때. 히스토그램은 겹치면 가린다.
  \item[Where] 데이터가 충분히 많을 때. 표본이 적으면 매끄러운 곡선이 없는 구조를
               지어낸다.
  \item[How] 구간을 충분히 잘게 나눠 곡선이 매끄러우면서도 분포를 왜곡하지 않는
             지점을 찾는다.
\end{inbrief}

강의의 기준은 명확하다. \emph{``곡선이 충분히 매끄러우면서도 분포를 정확히
표현하도록, 충분한 수의 구간을 쓴다.''} 두 조건이 서로 반대 방향으로 당기므로
절충점을 찾아야 한다 --- \Cref{subsec:binning}의 문제가 그대로 이어진다.

\begin{figure}[ht]\centering
  \includegraphics[width=\linewidth]{05_04_density_line.png}
  \caption{밀도 곡선(Density Line). 히스토그램의 계단을 매끄럽게 편 것이며,
  여러 분포를 겹쳐 비교할 때 서로를 가리지 않는다는 장점이 있다.}
  \label{fig:density}
\end{figure}

% =====================================================================
\section{에러바 --- 요약을 나란히 놓기}
\label{sec:errorbar}

\begin{inbrief}
  \item[Why] 비교할 범주가 너무 많아 분포를 전부 그릴 수 없을 때, 각 분포의
             \emph{요약}만 나란히 놓기 위해서다.
  \item[When] 연속 변수를 구간으로 나눠 그 함수로 볼 때, 또는 집단이 여럿일 때.
  \item[Where] 분포가 단봉이고 대칭일 때만 정직하다. 이봉분포는 요약이 거짓말을 한다.
  \item[How] 중심(평균)을 점이나 선으로, 퍼짐을 막대로 그린다.
\end{inbrief}

\begin{figure}[ht]\centering
  \includegraphics[width=\linewidth]{05_05_error_bar.png}
  \caption{집단별 평균 곡선과 퍼짐을 나타내는 막대. 강의의 예는 약물에 저항성인
  환자군과 민감한 환자군의 비교다.}
  \label{fig:errorbar}
\end{figure}

\begin{figure}[ht]\centering
  \includegraphics[width=\linewidth]{05_06_error_variants.png}
  \caption{에러바의 변형. 왼쪽은 에러 구름(Error Cloud), 오른쪽은 에러바를 얹은
  막대그래프(Bar Plot with Error Bars)다. 담는 정보는 같고 강조점이 다르다.}
  \label{fig:error-variants}
\end{figure}

\begin{pitfall}[막대가 무엇인지 캡션에 반드시 쓴다]
\label{pit:sd-se-ci}
에러바의 길이로 쓰이는 후보가 셋이고, \textbf{길이가 서로 몇 배씩 차이 난다}.

\begin{center}
\begin{tabularx}{\linewidth}{@{}l l X@{}}
\toprule
무엇 & 계산 & 무엇을 말하는가 \\
\midrule
표준편차 $\sd$ & $\sigma$ & \textbf{데이터}가 얼마나 퍼져 있는가 \\
표준오차 $\mathrm{SE}$ & $\sigma/\sqrt{n}$ & \textbf{평균 추정값}이 얼마나 불안정한가 \\
신뢰구간 $\mathrm{CI}$ & $\approx \bar{x} \pm 1.96\,\mathrm{SE}$ & 참평균이 있을 법한 범위 \\
\bottomrule
\end{tabularx}
\end{center}

버거 데이터로 보면 $\sd = 277.16$ 이지만 $\mathrm{SE} = 277.16/\sqrt{69} = 33.37$
로 \textbf{8배 넘게 차이} 난다. 캡션에 무엇인지 쓰지 않은 에러바는 읽을 수 없다.
\end{pitfall}

% =====================================================================
\section{바이올린 플롯}
\label{sec:violin}

\begin{inbrief}
  \item[Why] 에러바와 박스플롯이 버리는 \emph{분포의 모양}을 되살리기 위해서다.
  \item[When] 집단을 비교하는데 그 분포가 단봉이 아닐 가능성이 있을 때.
  \item[Where] 표본이 충분해야 한다. $n$ 이 작으면 밀도 추정이 없는 봉우리를 만든다.
  \item[How] 각 집단의 밀도 곡선을 세로로 세우고 좌우 대칭으로 펼친다.
\end{inbrief}

\begin{figure}[ht]\centering
  \includegraphics[width=\linewidth]{05_07_violin.png}
  \caption{바이올린 플롯. 폭이 곧 그 값 근처의 밀도이므로, 봉우리가 몇 개인지가
  즉시 보인다. 강의의 예는 보스턴 마라톤 완주 시간의 연도별 분포다.}
  \label{fig:violin}
\end{figure}

\begin{figure}[ht]\centering
  \includegraphics[width=\linewidth]{05_08_violin_combined.png}
  \caption{결합 바이올린 플롯. 두 집단을 좌우 반쪽씩 맞붙여 같은 축에서 비교한다.
  두 집단을 나란히 두 개 그리는 것보다 시선 이동이 짧다.}
  \label{fig:violin-combined}
\end{figure}

% =====================================================================
\section{바코드 차트}
\label{sec:barcode}

\begin{inbrief}
  \item[Why] 요약하지 않고 \textbf{모든 데이터 점}을 그대로 보여 주기 위해서다.
  \item[When] 데이터 개수가 적어 요약이 오히려 정보를 없앨 때.
  \item[Where] 각 행의 값 범위가 서로 비교 가능해야 한다. 아니면 행끼리 읽을 수 없다.
  \item[How] 데이터 점 하나를 선 위의 막대 하나로 그리고, 변수마다 행을 따로 준다.
\end{inbrief}

강의의 예는 워싱턴 메트로의 정시 운행률이다. 각 행이 한 노선, 각 행의 12개 막대가
월별 평균이고, 연평균은 굵은 막대로 표시된다.

\begin{figure}[ht]\centering
  \includegraphics[width=\linewidth]{05_09_barcode.png}
  \caption{바코드 차트. 점 하나하나가 보이므로 \emph{빈 구간}이 눈에 띈다 ---
  요약통계로는 절대 드러나지 않는 정보다.}
  \label{fig:barcode}
\end{figure}

\begin{keyidea}
바코드 차트는 이 장에서 유일하게 \textbf{아무것도 버리지 않는} 그림이다.
\analogy{지도}가 아니라 \analogy{지형 그 자체}다. 그래서 데이터가 조금만 많아져도
쓸 수 없지만, 적을 때는 가장 정직하다.
\end{keyidea}

% =====================================================================
\section{박스플롯으로 집단 비교하기}
\label{sec:boxplot}

\begin{inbrief}
  \item[Why] 여러 집단의 중심과 퍼짐을 한 화면에서 비교하기 위해서다.
  \item[When] 집단이 3개 이상이고 각 집단의 분포가 대체로 단봉일 때.
  \item[Where] 이봉분포에는 쓰면 안 된다 --- 박스플롯은 봉우리 개수를 표현하지 못한다.
  \item[How] $Q_1$ 에서 $Q_3$ 까지 상자를 그리고, $1.5 \times \IQR$ 이내의 가장 먼
             관측값까지 수염을 뻗고, 그 밖은 이상치로 찍는다.
\end{inbrief}

박스플롯을 만드는 규칙은 \Cref{subsec:fence}에서 이미 다뤘다. 여기서는 그것을
\emph{비교}에 쓰는 법을 본다.

\begin{figure}[ht]\centering
  \includegraphics[width=\linewidth]{05_11_box_compare.png}
  \caption{박스플롯으로 집단 비교하기. 중앙선의 높이 차이가 중심의 차이를,
  상자의 길이 차이가 퍼짐의 차이를 나타낸다.}
  \label{fig:box-compare}
\end{figure}

강의는 붓꽃 데이터에서 두 가지를 읽는다 --- 꽃받침(Sepal)의 중앙값이 꽃잎(Petal)
보다 높고, 꽃잎 쪽의 값의 퍼짐이 더 크다. 이 두 문장이 박스플롯에서 읽어야 할
것의 전부다. \textbf{중심의 차이}와 \textbf{퍼짐의 차이}.

과제 데이터에서도 같은 방식으로 읽힌다. 음식 유형별 칼로리 중앙값은
치킨너겟 240 $\rightarrow$ 감자튀김 315 $\rightarrow$ 그릴드치킨샌드위치 398
$\rightarrow$ 브레디드치킨샌드위치 510 $\rightarrow$ 밀크셰이크 585
$\rightarrow$ 버거 591 로 단조 증가한다.

\subsection{노치 박스플롯}
\label{subsec:notch}

\begin{figure}[ht]\centering
  \includegraphics[width=\linewidth]{05_12_box_notched.png}
  \caption{노치(Notch)를 넣은 박스플롯. 허리가 잘록한 부분이 중앙값의 신뢰구간을
  나타내며, 두 상자의 노치가 겹치지 않으면 중앙값이 다르다고 볼 근거가 된다.}
  \label{fig:box-notch}
\end{figure}

\term{노치}(Notch)는 중앙값 주변의 불확실성을 상자에 직접 새긴 것이다. 두 집단의
노치가 겹치지 않으면 중앙값이 통계적으로 다르다고 본다 --- \Cref{sec:ci}의
신뢰구간 비교를 박스플롯 안으로 가져온 것이다.

\begin{pitfall}[노치가 상자보다 커지면 표본이 부족한 것]
노치의 폭은 $\pm 1.58 \times \IQR/\sqrt{n}$ 으로 계산된다. $n$ 이 작으면 이 값이
$\IQR$ 자체보다 커져서 \textbf{노치가 상자 밖으로 뒤집혀 나온다}. 그림이
망가진 것이 아니라, ``이 표본으로는 중앙값을 말할 수 없다''는 경고다.
\end{pitfall}

% =====================================================================
\section{신뢰구간}
\label{sec:ci}

\begin{inbrief}
  \item[Why] 계산한 통계값에 \emph{얼마나 불확실성이 있는지}를 함께 말하기 위해서다.
  \item[When] 표본에서 계산한 값으로 모집단을 말하려 할 때 --- 즉 거의 항상.
  \item[Where] 두 집단이 ``정말로'' 다른지 판단할 때. 겹치는지만 보면 된다.
  \item[How] 공식($\bar{x} \pm z \cdot \sigma/\sqrt{n}$) 또는 부트스트랩
             (Bootstrap) 둘 중 하나로 구한다.
\end{inbrief}

\term{신뢰구간}(Confidence Interval)은 특정 통계값에 얼마나 불확실성이 있는지를
나타낸다. 설문 조사에서 ``오차 범위''(Margin of Error)라고 부르는 것과 같다 ---
전체 모집단을 조사했다면 얻었을 값이 이 범위 안에 있으리라는 뜻이다.

\subsection{공식으로 구하기}

\[
  \mathrm{CI} = \bar{x} \pm z \cdot \frac{s}{\sqrt{n}}
\]

\begin{center}
\begin{tabularx}{\linewidth}{@{}l X@{}}
\toprule
기호 & 의미 \\
\midrule
$s$ & 표준편차(Standard Deviation) \\
$n$ & 관측값의 개수 \\
$z$ & 표에서 읽는 $z$ 값. $95\%$ 신뢰구간이면 $1.96$ \\
$\bar{x}$ & 신뢰구간을 구하려는 대상 값 \\
\bottomrule
\end{tabularx}
\end{center}

$z = 1.96$ 의 의미는 이렇다. 평균에서 $-1.96$ 부터 $+1.96$ 표준편차 사이의 값이
무작위로 뽑힐 확률이 $95\%$ 다. 뒤집으면, 어떤 값이 무작위로 뽑힐 확률이 $5\%$
미만이면 그것을 \textbf{통계적으로 유의하다}(Statistically Significant)고 본다.

\subsection{부트스트랩으로 구하기}
\label{subsec:bootstrap}

공식은 분포에 대한 가정을 필요로 한다. \term{부트스트랩}(Bootstrapping)은 그
가정 없이, 컴퓨터로 재표집을 반복해서 같은 것을 구한다. 강의의 예는 암컷 생쥐
12마리의 몸무게다.

\begin{figure}[ht]\centering
\begin{tikzpicture}[
  node distance = 6mm,
  bsbox/.style = {draw = bookrule, fill = booksurface, rounded corners = 2pt,
                  minimum width = 62mm, minimum height = 8.5mm, align = center,
                  font = \small},
  arrow/.style = {-{Stealth[length=2.2mm]}, draw = bookmuted}]
  \node[bsbox] (a) {표본 평균을 계산한다};
  \node[bsbox, below = of a] (b)
        {12개를 \textbf{복원추출}(Replacement)로 다시 뽑는다\\[-1pt]
         \footnotesize 같은 개체가 여러 번 뽑혀도 된다};
  \node[bsbox, below = of b] (c) {재표집된 12개의 평균을 계산한다};
  \node[bsbox, below = of c] (d) {2--3 단계를 \textbf{10{,}000회 이상} 반복한다};
  \node[bsbox, below = of d, fill = bookaccentbg, draw = bookaccent] (e)
        {10{,}000개 평균의 2.5\%\,--\,97.5\% 구간\\[-1pt]
         \footnotesize 이것이 95\% 신뢰구간이다};
  \draw[arrow] (a) -- (b);
  \draw[arrow] (b) -- (c);
  \draw[arrow] (c) -- (d);
  \draw[arrow] (d) -- (e);
  \draw[arrow] (d.east) -- ++(9mm,0) |- (b.east);
\end{tikzpicture}
\caption{부트스트랩 절차. 핵심은 \textbf{복원추출}이다 --- 중복을 허용하지 않으면
매번 같은 표본이 나와 아무 변동도 생기지 않는다.}
\label{fig:bootstrap-flow}
\end{figure}

\begin{figure}[ht]\centering
  \includegraphics[width=\linewidth]{05_13_bootstrap_ci.png}
  \caption{버거 칼로리의 부트스트랩 신뢰구간. 재표집 10{,}000회의 평균 분포이며,
  2.5 백분위수와 97.5 백분위수가 신뢰구간의 양끝이 된다.}
  \label{fig:bootstrap}
\end{figure}

\begin{pitfall}[이 히스토그램은 버거의 분포가 아니다]
\Cref{fig:bootstrap}의 히스토그램은 \textbf{10{,}000개 ``평균''의 분포}이지
69개 버거의 분포가 아니다. 가로축 범위가 555--690 정도로 좁은 이유가 여기 있다 ---
개별 버거는 140부터 1240까지 퍼져 있지만, 69개의 평균은 그만큼 흔들리지 않는다.
이 둘을 혼동하면 ``버거는 전부 600칼로리 근처''라는 완전히 틀린 결론에 이른다.
\end{pitfall}

\begin{keyidea}
표본 크기를 키우면 신뢰구간은 \textbf{좁아지고}, 표준편차가 크면 신뢰구간은
\textbf{넓어진다}. 공식 $\bar{x} \pm z \cdot s/\sqrt{n}$ 에서 $n$ 은 분모에,
$s$ 는 분자에 있기 때문이다.
\end{keyidea}

\subsection{두 집단을 비교하기}
\label{subsec:ci-compare}

\begin{figure}[ht]\centering
  \includegraphics[width=\linewidth]{05_14_ci_compare.png}
  \caption{신뢰구간으로 두 집단 비교하기. 두 구간이 겹치지 않으므로 두 평균은
  통계적으로 다르다($p < 0.05$).}
  \label{fig:ci-compare}
\end{figure}

버거의 평균 칼로리는 620이고 $95\%$ 신뢰구간은 $[557,\ 686]$, 치킨너겟은 평균
275에 신뢰구간 $[211,\ 345]$ 다. 두 구간 사이에는 345와 557 사이의 빈틈이 있다.
\textbf{겹치지 않으므로} 두 평균이 우연히 다르게 나왔을 가능성은 $5\%$ 미만이다.

\begin{pitfall}[겹친다고 같은 것은 아니다]
이 판정법은 한 방향으로만 강하다. 구간이 \textbf{겹치지 않으면} 다르다고 말할 수
있지만, \textbf{겹친다고 해서 같다}고 말할 수는 없다. 겹치는 경우는 ``이 데이터
로는 판단할 수 없다''는 뜻이고, 정확한 판정에는 별도의 검정이 필요하다.
두 평균을 직접 비교하는 검정은 \Cref{sec:cd-t-two}에서 다룬다.
\end{pitfall}

% =====================================================================
\section{형태 고르기 --- 정리}
\label{sec:choosing-form}

\begin{table}[ht]\centering
\caption{질문과 형태의 대응. 위에서 아래로 갈수록 데이터를 덜 버린다.}
\label{tab:form-choice}
\begin{tabularx}{\linewidth}{@{}l l X@{}}
\toprule
묻는 것 & 형태 & 그 형태인 이유 \\
\midrule
한 변수의 분포 & 히스토그램(Histogram) & 구간별 개수를 직접 센다 \\
분포를 겹쳐서 비교 & 밀도 곡선(Density Line) & 겹쳐도 서로를 가리지 않는다 \\
여러 집단의 분포 & 작은 배수(Small Multiples) & 축을 공유해 비교를 성립시킨다 \\
집단이 많을 때의 요약 & 에러바(Error Bar) & 중심과 퍼짐만 남겨 자리를 아낀다 \\
집단의 중심과 퍼짐 & 박스플롯(Box Plot) & 다섯 개 숫자로 요약해 나란히 놓는다 \\
중앙값의 차이 판정 & 노치 박스플롯(Notched Box Plot) & 불확실성을 상자에 새긴다 \\
분포의 \emph{모양}까지 & 바이올린(Violin Plot) & 박스가 버린 봉우리를 되살린다 \\
모든 점을 그대로 & 바코드(Barcode Chart) & 아무것도 버리지 않는다 \\
추정값의 불확실성 & 신뢰구간(Confidence Interval) & 겹치는지로 차이를 판정한다 \\
\bottomrule
\end{tabularx}
\end{table}

\section{요약}

\begin{itemize}
  \item 분위수는 정렬 $\rightarrow$ 위치 부여 $\rightarrow$ 보간으로 계산되며,
        관측값 중 하나일 필요가 없다(\Cref{ex:quantile}).
  \item 히스토그램의 구간 개수는 \emph{선택}이고, 선택이 결론을 바꾼다.
  \item 에러바의 막대가 $\sd$인지 $\mathrm{SE}$인지 $\mathrm{CI}$인지 반드시
        밝힌다. 길이가 몇 배씩 차이 난다(\Cref{pit:sd-se-ci}).
  \item 박스플롯은 봉우리 개수를 표현하지 못한다. 의심되면 바이올린이나 바코드로
        확인한다.
  \item 신뢰구간이 겹치지 않으면 다르다. 겹친다고 같은 것은 아니다.
  \item 부트스트랩은 분포 가정 없이 신뢰구간을 구한다. 핵심은 복원추출이다.
\end{itemize}

% =====================================================================
\section{어떻게 적용하는가}
\label{sec:vis-apply}

이 장의 차트와 개념은 이 책의 실전 장에서 되풀이해 쓰인다.

\begin{itemize}
  \item \Cref{ch:hw1}(과제 1 사례집)은 버거와 밀크셰이크의 칼로리
        분포를 이 장의 차트가 아니라 평균·SD·최소·최대로 된 요약통계 표로
        비교한다(\Cref{tab:hw1-burger-vs-shake}). 버거와 치킨너겟을 신뢰구간으로
        직접 겹쳐 보는 비교는 이 장의 \Cref{subsec:ci-compare}에서 한다.
  \item \Cref{ch:hw2}(과제 2, 분포 비교)는 영화 평점과 나이 분포에는 이 장의
        히스토그램을 쓰지만, 공항 노선 수는 꼬리가 두 자릿수 넘게 퍼져
        히스토그램 대신 로그-로그 축의 CCDF(Complementary Cumulative
        Distribution Function)를 쓴다.
  \item \Cref{sec:qq}(Q-Q 플롯, 강의 7)은 \Cref{ex:quantile}과 같은 분위수
        비교 아이디어를 쓰지만, 플로팅 위치(Plotting Position) 규칙은 다르다.
        \Cref{ex:quantile}은 $(i-1)/(n-1)$을, \Cref{sec:qq}는 $p_i = i/(n+1)$을
        쓴다(\Cref{sec:qq}의 함정에서 둘을 대조한다).
  \item \Cref{sec:cd-t-two}(강의 11)의 두 표본 $t$검정(Two-Sample $t$-Test)은
        이 장의 신뢰구간 겹침 판정을 평균 비교로 정식화한 것이다.
        \Cref{sec:cd-ks-two}의 두 표본 KS검정(Kolmogorov-Smirnov Test)은
        평균이 아니라 분포 모양 전체를 비교한다.
\end{itemize}

% =====================================================================
\section{앞으로}
\label{sec:vis-future}

\Cref{subsec:ci-compare}의 신뢰구간 판정과 $5\%$ 유의 기준은 이 장에서는
직관으로만 제시된다. 그 직관을 귀무가설(Null Hypothesis)·유의수준
(Significance Level)·기각 규칙(Rejection Rule)으로 정식화한 것이
\Cref{ch:hypothesis}다. 또한 \Cref{subsec:ci-compare}의 함정 상자가 던진
``구간이 겹치면 같은 것인가''라는 질문도 \Cref{ch:hypothesis}가 답한다.

신뢰구간 겹침 판정은 평균 하나만 비교한다. 평균만 더 정식으로 비교하려면
\Cref{sec:cd-t-two}의 두 표본 $t$검정(Two-Sample $t$-Test)이 필요하고,
분포 모양 전체를 비교하려면 \Cref{sec:cd-ks-two}의 두 표본
KS검정(Kolmogorov-Smirnov Test)이 필요하다.

\chapter{분포 --- 데이터가 따르는 모양}
\label{ch:distributions}
% source: deck 6 slide 1 (내용 없음, 제목만)

\begin{keyidea}
분포를 가정하는 것은 \analogy{투영법}을 고르는 일이다. 편리해지지만 반드시
왜곡이 따르고, 어디가 왜곡되었는지는 검증해야만 알 수 있다.
\end{keyidea}

\begin{inbrief}
  \item[Why] 데이터의 모양을 이름 있는 \term{분포}(Distribution) 하나로
             요약하면, 모수(Parameter) 몇 개만으로 확률을 계산할 수 있다.
  \item[When] 히스토그램의 모양(대칭, 한쪽 꼬리, 평평함, 두 값)이 보이고,
              그 모양을 계산에 쓰고 싶을 때.
  \item[Where] 정규(Normal), 지수(Exponential), 멱법칙(Power Law, 연속),
               균등(Uniform, 이산과 연속), 베르누이(Bernoulli),
               이항(Binomial, 이산). 뒤의 검정(\Cref{ch:hypothesis})과
               정규화(\Cref{ch:qq-normalization})가 이 장의 가정 위에 선다.
  \item[How] 히스토그램 모양으로 후보를 고른다 $\to$ 모수를 추정한다
             $\to$ CDF를 겹쳐 얼마나 어긋나는지 눈으로, 또는 검정으로
             확인한다.
\end{inbrief}

이 장의 출처는 \texttt{notes/6\_DataDistribution.pptx}(30 슬라이드) 하나다.
슬라이드가 외부 문헌을 밝히지 않는 자리는 슬라이드 자체를 출처로 적는다.

% =====================================================================
\section{배경}
\label{sec:dist-background}
% source: deck 6 slide 2, 3

\paragraph{왜 분포인가.} 측정값을 하나씩 늘어놓으면 숫자 더미일 뿐이지만,
같은 값 근처가 얼마나 자주 나오는지를 세어 모양으로 그리면(히스토그램,
\Cref{sec:histogram}) 그 모양에서 규칙이 보인다. 모양이 가장 높이 솟은
자리가 가장 흔한 값이고, 낮은 자리는 드문 값이다. 이 장은 실제로 자주
나타나는 몇 가지 모양에 이름을 붙이고, 그 이름 하나와 모수 몇 개만으로
확률을 계산하는 법을 다룬다.

이 장은 다음 개념을 그대로 가져와 쓴다.

\begin{itemize}
  \item 평균과 표준편차(\Cref{sec:center}, \Cref{def:spread}): 분포의
        중심과 폭을 요약하는 숫자. 이 장의 정규분포, 균등분포, 지수분포
        모두 평균과 분산(또는 표준편차)으로 모양이 정해진다.
  \item 히스토그램과 구간 나누기(\Cref{sec:histogram}): 분포의 모양을
        눈으로 보는 도구.
  \item 왜도와 그 부호(\Cref{sec:skewness}, \Cref{sec:skew-sign}):
        정규분포처럼 대칭이 아닌 모양을 설명하는 말.
\end{itemize}

\term{확률밀도함수}(Probability Density Function, PDF)와
\term{확률질량함수}(Probability Mass Function, PMF)의 구분도 여기서
한 번 깐다. 연속형 변수(어떤 실수든 값이 될 수 있는 변수, 예: 키)는
PDF를 쓰고, 그 값 자체는 확률이 아니라 \textbf{밀도}다. 확률은 PDF
곡선 아래 \textbf{넓이}로만 정해진다. 이산형 변수(정해진 값만 가지는
변수, 예: 주사위 눈)는 PMF를 쓰고, PMF 값 자체가 그 값이 나올 확률이다.

슬라이드 3은 흔한 분포를 연속형과 이산형으로 나눈다. \Cref{tab:dist-map}이
그 목록을 이 장의 절과 짝짓는다. 슬라이드는 정규 $\to$ 균등 $\to$
지수·멱법칙 $\to$ 베르누이 $\to$ 이항 순서로 나열하지만, 이 장은 지수분포와
멱법칙을 먼저 묶고(둘 다 ``꼬리가 얼마나 빨리 줄어드는가''를 다룬다)
균등분포를 뒤로 옮겼다. 슬라이드 순서를 그대로 따르지 않는다.

\begin{table}[ht]\centering
\begin{tabularx}{\linewidth}{@{}l l X@{}}
\toprule
분포 & 연속/이산 & 이 장의 절 \\
\midrule
정규(Normal)     & 연속 & \Cref{sec:normal} \\
지수(Exponential) & 연속 & \Cref{sec:exponential} \\
멱법칙(Power Law) & 연속 & \Cref{sec:powerlaw} \\
균등(Uniform)     & 연속·이산 둘 다 & \Cref{sec:uniform} \\
베르누이(Bernoulli) & 이산 & \Cref{sec:binomial} \\
이항(Binomial)     & 이산 & \Cref{sec:binomial} \\
\bottomrule
\end{tabularx}
\caption{슬라이드 3의 분포 목록과 이 장의 절.}
\label{tab:dist-map}
\end{table}

% =====================================================================
\section{정규분포}
\label{sec:normal}

\begin{inbrief}
  \item[Why] 자연에서 반복해서 나타나는 종 모양의 분포를 평균과 분산,
             모수 둘만으로 완전히 표현하기 위해.
  \item[When] 측정값이 평균 주변에 좌우 대칭으로 몰리고, 평균에서 멀어질수록
              드물어질 때.
  \item[Where] 키, 몸무게, 성적처럼 자연·사회 현상 전반. 이 책에서는
               패스트푸드 \texttt{calories}.
  \item[How] 평균 $\mu$로 중심을, 표준편차 $\sigma$로 폭을 잡는다. PDF
             곡선 아래 넓이가 항상 1이 되도록 높이가 자동으로 맞춰진다.
\end{inbrief}

% source: deck 6 slide 4
정규분포(Normal Distribution)는 \term{가우스분포}(Gaussian Distribution)
또는 \term{종 모양 곡선}(Bell Curve)이라고도 부른다. 자연에서 특히 흔한
이유는 회사의 급여 분포, 키, 몸무게, 성적처럼 많은 독립적인 원인이 더해져
만들어지는 값들이 대체로 이 모양을 따르기 때문이다.

% source: deck 6 slide 5
세로축을 읽을 때 흔히 하는 오해가 하나 있다. PDF의 세로축 값은
확률이 \textbf{아니라} 상대적 가능성(밀도)이다. 아주 크거나 아주 작은
값일수록 곡선이 낮아 가능성이 작고, 평균 근처일수록 곡선이 높아
가능성이 크다는 뜻일 뿐, 특정 $x$ 하나가 나올 확률은 여전히 그 점에서의
넓이(사실상 0)로 정해진다. 실제 확률은 구간에 대해 곡선 아래 넓이를
적분해야 나온다.

% source: deck 6 slide 6
\Cref{sec:dist-background}의 배경 개념(연속 변수는 PDF, 확률은 넓이로만
읽는다)이 이 절 전체의 전제다. 예를 들어 아기 키 분포와 성인 키 분포를
겹쳐 그리면, 두 곡선 모두 각자의 평균을 중심으로 대칭이지만 폭(표준편차)이
다르다. 아기 키의 곡선이 성인 키의 곡선보다 뾰족하고 높은데, 이는 아기가
더 ``특별해서''가 아니라 \textbf{곡선 아래 넓이가 항상 1로 고정}되어
있기 때문이다. 아기 키는 성인 키보다 흩어짐(표준편차)이 작아 폭이
좁으므로, 넓이 1을 유지하려면 높이가 커져야 한다.

% source: deck 6 slide 10 (주석만 추가, 정의 자체는 이전 슬라이드와 동일)
\begin{definition}[정규분포]
\label{def:normal}
\term{정규분포}(Normal Distribution) $X \sim \mathcal{N}(\mu, \sigma^2)$ 의
확률밀도함수(Probability Density Function)는
\[
  f(x) = \frac{1}{\sigma\sqrt{2\pi}}
         \exp\!\left(-\frac{(x-\mu)^2}{2\sigma^2}\right)
\]
\end{definition}

% source: deck 6 slide 7
정규분포의 폭을 $\sigma$ 단위로 재면 어디까지가 ``흔한 범위''인지 규칙이
생긴다. 이를 68-95-99.7 규칙이라 부른다.

\begin{pitfall}[95\%가 아니라 95.45\%]
슬라이드 값은 $\pm 2$ 표준편차 안에 $95\%$, 다시 계산하면 $95.45\%$다. 표준정규분포의
누적분포함수로 $P(|Z|<2) = \Phi(2)-\Phi(-2)$를 계산하면 $0.9545$가
나온다. 정확히 $95\%$가 되는 경계는 $\pm 2$가 아니라 $\pm \mathbf{1.96}$
표준편차다(\Cref{sec:ci}에서 신뢰구간의 $1.96$과 같은 숫자). 슬라이드 7의
그림은 오히려 $1.96\sigma$ 쪽 눈금을 맞게 그려 두었으므로,
틀린 것은 본문의 ``$95\%$''라는 말뿐이다.
\end{pitfall}

\Cref{fig:dist-sigma-bands}가 패스트푸드 \texttt{calories}에 맞춘 정규
곡선 위에 $\pm 1\sigma$, $\pm 2\sigma$ 구간을 표시한다.

\begin{figure}[ht]\centering
  \includegraphics[width=0.85\linewidth]{06_02_normal_sigma_bands.png}
  \caption{패스트푸드 \texttt{calories}에 맞춘 정규 곡선과
  $\pm 1\sigma$($68.27\%$), $\pm 2\sigma$($95.45\%$) 구간. 실제 데이터
  중 $\pm 2\sigma$ 안에 든 비율은 $96.0\%$로, 이론값 $95.45\%$와
  가깝다.}
  \label{fig:dist-sigma-bands}
\end{figure}

\Cref{fig:normal-fit}(\Cref{ch:qq-normalization})은 같은 \texttt{calories}
데이터에 최대우도로 정규 곡선을 맞춘 그림이다. 히스토그램은 종 모양에
가까워 보이지만, 바로 아래 함정과 \Cref{sec:qq-apply}의 Q-Q 플롯은
정규성을 기각한다.

% source: deck 6 slide 8
분포를 이동하거나 척도를 바꾸는 조작도 자주 쓴다. $Z = (X-\mu)/\sigma$로
바꾸면 평균 0, 표준편차 1인 표준정규분포가 되는데, 이 변환의 자세한
성질과 다른 변환(범위 정규화 등)과의 비교는
\Cref{sec:normalization}(\Cref{ch:qq-normalization})에서 다룬다.

% source: deck 6 slide 9
정규분포가 아닌 모양도 있다. 오른쪽으로 긴 꼬리(양의 왜도), 왼쪽으로 긴
꼬리(음의 왜도), 대칭(왜도 0에 가까움), 이 세 모양의 정의와 판정 규칙은
\Cref{sec:skewness}와 \Cref{sec:skew-sign}이 다룬다(\Cref{sec:dist-background}에서
이미 배경으로 소개했다). 이 장은 링크만 걸고 새로 설명하지 않는다.

\begin{pitfall}[정규성 검정의 함정]
같은 표본에서 $\mu, \sigma$ 를 추정한 뒤 그 표본으로 KS 검정(Kolmogorov--Smirnov
Test)을 하면 p-value가 낙관적으로 나온다. 이 경우 Shapiro--Wilk 검정을 쓴다.
% TODO: calories 의 Shapiro p = 0.00038 -- 정규성이 기각된 사례로 정리.
\end{pitfall}

% =====================================================================
\section{지수분포}
\label{sec:exponential}

\begin{inbrief}
  \item[Why] 다음 사건까지 걸리는 시간(대기 시간, 수명)을 모수 하나
             $\lambda$로 표현하기 위해.
  \item[When] 사건이 평균적으로 일정한 비율(rate)로 무작위하게 일어날 때.
  \item[Where] 노트북 수명, 커피숍 손님 도착 간격, 패스트푸드
               \texttt{sugars}.
  \item[How] $\lambda$는 단위 시간당 사건 발생 비율. CDF
             $P(X<x) = 1 - e^{-\lambda x}$로 ``$x$ 안에 사건이 일어날
             확률''을 구한다.
\end{inbrief}

% source: deck 6 slide 15
지수분포(Exponential Distribution)는 사건과 사건 사이에 걸리는 시간을
모형화할 때 쓴다. 예를 들어 가게에서 다음 손님이 올 때까지 기다리는
시간이 그렇다.

\begin{pitfall}[모수 범위]
슬라이드 값은 모수 범위가 ``$0 < \lambda \le 1$''이고, 다시 계산하면
\textbf{$\lambda > 1$도 얼마든지 가능하다}. 올바른 범위는 $\lambda > 0$이다. 슬라이드 17
자신이 시간당 손님 10명, 즉 $\lambda = 10$을 예제로 쓰고 있어 슬라이드
15의 ``$\le 1$'' 제한과 슬라이드 17이 서로 어긋난다.
\end{pitfall}

% source: deck 6 slide 16
\begin{example}[노트북 수명]
\label{ex:dist-laptop}
델 노트북의 평균 수명이 5년일 때, 3년 안에 고장 날 확률을 구한다(슬라이드
16).
\begin{align*}
  \lambda &= 1/5 = 0.2      && \text{평균 5년의 역수}\\
  P(X<3) &= 1 - e^{-\lambda \cdot 3}  && \text{지수분포 CDF}\\
         &= 1 - e^{-0.2\times 3} && \text{$\lambda$ 대입}\\
         &= 1 - e^{-0.6}   && \text{정리}\\
         &= 1 - 0.548812   && \text{$e^{-0.6}$ 계산}\\
         &= \mathbf{0.4512} && \text{슬라이드 값과 확인}
\end{align*}
\end{example}

% source: deck 6 slide 17
\begin{example}[커피숍 손님 도착]
\label{ex:dist-coffee}
시간당 평균 10명이 오는 커피숍(즉 $\lambda = 10$)에서, 다음 손님까지
평균 대기 시간과 5분 안에 손님이 올 확률을 구한다.
\begin{align*}
  \text{평균 대기 시간} &= 1/\lambda           && \text{평균 대기 시간 공식}\\
    &= 1/10\,\text{시간}                        && \text{$\lambda=10$ 대입}\\
    &= 0.1\,\text{시간}                          && \text{소수로 정리}\\
    &= \mathbf{6}\,\text{분}                     && \text{시간을 분으로 환산}
\end{align*}
\end{example}

\begin{pitfall}[슬라이드 값은 0.0801, 다시 계산하면 0.5654]
$5$분을 시간 단위로 바꾸면 $5/60$시간이다.
\begin{align*}
  P(X<5/60) &= 1 - e^{-\lambda \cdot 5/60}   && \text{지수분포 CDF}\\
            &= 1 - e^{-10\times 5/60}        && \text{$\lambda=10$ 대입}\\
            &= 1 - e^{-0.8333}               && \text{지수 정리}\\
            &= 1 - 0.434598                  && \text{$e^{-0.8333}$ 계산}\\
            &= \mathbf{0.5654}               && \text{슬라이드 0.0801과 다르다}
\end{align*}
슬라이드의 $\mathbf{0.0801}$은 $\lambda$를 빠뜨리고
$1-e^{-5/60} = \mathbf{0.0800}$을 계산한 값과 거의 같다(넷째 자리에서만
차이). 상식으로도 검산된다. 평균 대기 시간이 $6$분인 가게에서 $5$분
안에 손님이 올 확률이 $8\%$뿐일 수는 없다. 절반을 넘는 것이 맞다.
\end{pitfall}

\subsection{기하분포로 근사하기}
\label{subsec:geometric}
% source: deck 6 slide 18

\paragraph{고치는 문제.} 지수분포는 연속 시간의 모형이라 ``정확히 몇 분째에
사건이 일어나는가''를 묻기는 어렵다. 슬라이드 18은 이를 1분 단위로 쪼갠
이산 질문으로 바꾼다: 첫 방문자가 1분째, 2분째, 3분째 각각에 올 확률은
얼마인가.

\paragraph{아이디어.} 1분마다 ``이번 1분 동안 방문자가 온다''는 베르누이
시행이 반복된다고 \textbf{근사}한다. 성공 확률은 슬라이드 17과 같은 방식으로
평균 대기 시간의 역수에서 얻는다(1분 평균 도착률 $\lambda = 1/20 = 0.05$,
\Cref{ex:dist-website} 참고). 이때 ``$k$번째 분에 처음 성공한다''는 확률은
기하분포(Geometric Distribution)
\[
  P(X=k) = (1-p)^{k-1} p
\]
를 따른다. 이는 근사이지 \textbf{오류가 아니다}. 실제 도착은 분 중간
어디서나 일어날 수 있는데, 그 안의 정확한 시점을 무시하고 ``그 분 안에서
왔는가/아닌가''만 보기 때문에 지수분포의 정확한 값과 조금 어긋난다.

\begin{example}[웹사이트 첫 방문자]
\label{ex:dist-website}
평균 20분에 한 번 방문자가 온다면 $\lambda = 1/20 = 0.05$(분당)다.
\begin{align*}
  P(0<x<1) &= p = \mathbf{0.05}          && \text{1분째 성공}\\
  P(1<x<2) &= (1-p)\,p                    && \text{1분 놓치고 2분째 성공}\\
           &= 0.95\times 0.05             && \text{$p=0.05$ 대입}\\
           &= \mathbf{0.0475}              && \text{정리}\\
  P(2<x<3) &= (1-p)^2 p                   && \text{2분 놓치고 3분째 성공}\\
           &= 0.95^2\times 0.05           && \text{$p=0.05$ 대입}\\
           &= \mathbf{0.045125}            && \text{정리}
\end{align*}
\end{example}

\Cref{tab:dist-geometric}가 이 기하분포 근사값을 지수분포의 정확한 값과
나란히 놓는다. 지수분포의 정확한 값은
$P(k-1<X<k) = e^{-\lambda(k-1)} - e^{-\lambda k}$로 구한다.

\begin{table}[ht]\centering
\begin{tabularx}{\linewidth}{@{}X c c c@{}}
\toprule
분(1분 구간) & 기하분포 근사 & 지수분포 정확값 & 차이 \\
\midrule
1분째 & 0.05000 & 0.04877 & 0.0012 \\
2분째 & 0.04750 & 0.04639 & 0.0011 \\
3분째 & 0.045125 & 0.04413 & 0.0010 \\
\bottomrule
\end{tabularx}
\caption{첫 방문자 도착 확률: 기하분포 근사값과 지수분포 정확값 비교.
근사와 정확값의 차이는 $0.001$ 안팎으로 작다.}
\label{tab:dist-geometric}
\end{table}

이 근사가 성립하는 이유는 각 1분 구간의 성공 여부가 \textbf{독립}이기
때문이다(\Cref{sec:ht-background}의 곱셈 규칙과 같은 전제). 독립이
아니라면(예: 손님이 몰려오는 시간대) 이 근사는 무너진다.

\Cref{fig:dist-exponential}가 \texttt{sugars} 데이터에 맞춘 지수분포의
PDF와 CDF를 함께 보여준다.

\begin{figure}[ht]\centering
  \includegraphics[width=0.95\linewidth]{06_04_exponential_fit.png}
  \caption{패스트푸드 \texttt{sugars}에 맞춘 지수분포. 평균
  $13.29\,\text{g}$에서 $\lambda = 1/\text{평균} = 0.0753$을 얻는다.}
  \label{fig:dist-exponential}
\end{figure}

% =====================================================================
\section{멱법칙}
\label{sec:powerlaw}
% source: deck 6 slide 19, 20, 21, 22

\begin{inbrief}
  \item[Why] 지수분포보다 훨씬 천천히 줄어드는 꼬리(극단적으로 큰 값이
             지수분포의 예측보다 훨씬 자주 나오는 경우)를 표현하기
             위해.
  \item[When] 소수가 대부분을 차지하는 현상, 즉 상위 몇 개가 전체 양의
              대부분을 차지할 때.
  \item[Where] 도시 인구, 단어 등장 빈도, 부의 분포, HW2의 공항 노선 수.
  \item[How] $f(x) \propto x^{-\alpha}$, $x \ge x_{\min}$. $\alpha$가
             작을수록 꼬리가 두껍다(큰 값이 더 자주 나온다).
\end{inbrief}

\term{멱법칙}(Power Law) 분포는 \term{파레토분포}(Pareto Distribution)라고도
부르며, 소수의 항목이 전체 발생 횟수나 값의 대부분을 차지하는 현상을
설명한다. 이른바 ``80/20 법칙''(파레토 원칙, Pareto Principle)이 대표적인
예다. 예를 들어 사회 전체 부의 약 $80\%$를 인구의 약 $20\%$가 가지고
있다는 관찰이 여기서 나온다.

\begin{definition}[멱법칙 분포]
\label{def:powerlaw}
멱법칙 분포는 $x \ge x_{\min}$에서
\[
  f(x) \propto x^{-\alpha}, \qquad \alpha > 1
\]
형태의 확률밀도함수를 갖는다. $\alpha$는 지수(Power Law Exponent),
$x_{\min}$은 멱법칙이 시작되는 하한이다.
\end{definition}

지수분포와 마찬가지로 값이 클수록 드물지만, 지수분포처럼 지수적으로
줄어들지 않고 \textbf{거듭제곱으로만} 줄어들기 때문에 훨씬 천천히
감소한다. 네트워크의 연결 수, 도시 인구, 단어 사용 빈도처럼 소수가
대부분을 차지하는 데이터에서 흔히 등장한다(슬라이드 20).

\term{여집합 누적분포함수}(Complementary CDF, $P(X \ge x)$)로 그리면
멱법칙을 더 쉽게 확인할 수 있다. 예를 들어 ``$x$번 이상 등장한 단어의
개수''를 세어 $x$에 대해 그리면 멱법칙은 로그-로그(log-log) 축에서
직선에 가깝다(슬라이드 21, University of Michigan의 Mark Newman
강의 노트를 슬라이드 자체가 인용한다).

\Cref{fig:dist-powerlaw}가 멱법칙과 지수분포를 로그-로그 축에서 비교한다
(슬라이드 22, Santa Fe Institute의 Aaron Clauset 강의 노트를 슬라이드가
인용한다). 멱법칙은 직선으로, 지수분포는 아래로 휘는 곡선으로
나타난다.

\begin{figure}[ht]\centering
  \includegraphics[width=0.95\linewidth]{06_06_powerlaw_vs_exponential.png}
  \caption{멱법칙(직선)과 지수분포(아래로 휘는 곡선)를 로그-로그 축에서
  비교. 오른쪽 패널은 패스트푸드 \texttt{sugars}의 여집합 CDF로,
  전체 $126$개 항목 중 로그 축에 표시할 수 없는 $0$ 이하 값을 뺀
  $109$개만 그렸다. 그 $109$개조차 꼬리가 짧아 멱법칙이라 부르기엔
  근거가 부족하다.}
  \label{fig:dist-powerlaw}
\end{figure}

이 데이터셋 안에는 멱법칙을 뚜렷이 보이는 변수가 없다. 실제 멱법칙
데이터와 모수 추정(최대우도, $\hat\alpha = 1 + n/\sum_i \ln(x_i/x_{\min})$)은
HW2의 공항 노선 수 예제(\Cref{sec:hw2-estimators},
\Cref{fig:hw2-airport-powerlaw})를 참고한다.

% =====================================================================
\section{균등분포}
\label{sec:uniform}

\begin{inbrief}
  \item[Why] 모든 결과가 똑같이 나올 가능성을 가질 때를 표현하기 위해.
  \item[When] 정해진 범위 안에서 특정 값을 더 선호할 이유가 없을 때.
  \item[Where] 주사위(이산), 난수 생성기(연속). 이 데이터셋에는 실제로
               균등분포를 따르는 변수가 없다.
  \item[How] 이산이면 각 결과에 $1/n$을, 연속이면 구간 폭의 역수를
             밀도로 준다.
\end{inbrief}

% source: deck 6 slide 11
균등분포(Uniform Distribution)는 결과가 정해진 범위 안에서 임의로
나오는 실험을 설명하며, 모든 결과가 똑같이 나올 가능성(Equally
Likely)을 갖는다. 공정한 6면체 주사위를 굴리는 것이 대표적인 이산
균등분포(Discrete Uniform Distribution)의 예로, $P(X=1) = \cdots =
P(X=6) = 1/6$이고 모든 확률의 합은 1이다.

% source: deck 6 slide 12
\begin{example}[연속 균등분포]
\label{ex:dist-uniform-continuous}
구간 $[0, 5]$ 위의 연속 균등분포를 생각한다. 밀도 함수는 구간
전체에서 상수 $h$이고, 넓이가 1이어야 하므로
\begin{align*}
  5 \times h &= 1  && \text{구간 폭 5, 넓이 1}\\
  h &= 1/5          && \text{밀도의 높이}
\end{align*}
슬라이드는 이 식을 ``$5\times X = 1 \to X = 1/5$''로 적어 구하는
대상을 확률변수 $X$처럼 보이게 썼지만, 실제로 구하는 것은 $X$가 아니라
밀도의 \textbf{높이} $h$다. 이 높이 $h=1/5$를 써서 구간 확률을 구한다.
\begin{align*}
  P(1<X<2) &= 1 \times \tfrac15 = \tfrac15
    && \text{폭 1 구간}\\
  P(4<X<4\tfrac13) &= \tfrac13 \times \tfrac15 = \mathbf{1/15}
    && \text{폭 $1/3$ 구간}\\
  P(X=4) &= 0
    && \text{점 하나의 넓이는 0}
\end{align*}
\end{example}

% source: deck 6 slide 14
일반적으로 구간 $[a, b]$ 위의 균등분포는(기호는 \Cref{sec:qq-background}의
함정에서 정한 것과 같이 $[a,b]$로 통일한다) 다음 평균과 표준편차를
갖는다.
\[
  \E[X] = \frac{a+b}{2}, \qquad
  \sd(X) = \frac{b-a}{\sqrt{12}}
\]
$[0,5]$ 구간에 대입하면
\begin{align*}
  \E[X] &= (0+5)/2 = 2.5            && \text{평균}\\
  \sd(X) &= 5/\sqrt{12} = 1.4434    && \text{표준편차}
\end{align*}

% source: deck 6 slide 13
슬라이드 13은 생후 8주 아기의 웃는 시간 55개를 난수로 만드는 예제를
든다. 균등분포에서 난수를 생성하는 방식(각 값이 똑같이 나올 가능성을
갖도록)을 보여주는 자리이며, 이 책은 슬라이드를 출처로만 인용하고 개별
값은 옮겨 적지 않는다.

\Cref{fig:dist-uniform}이 이산 균등분포(주사위)와 연속 균등분포
$[0,5]$를 함께 보여준다.

\begin{figure}[ht]\centering
  \includegraphics[width=0.95\linewidth]{06_05_uniform_reference.png}
  \caption{이산 균등분포(주사위, 왼쪽)와 연속 균등분포 $[0,5]$(오른쪽).
  이 데이터셋 안에는 균등분포를 따르는 변수가 없어 이론적 참고 그림으로
  그렸다.}
  \label{fig:dist-uniform}
\end{figure}

\begin{pitfall}[serving\_size는 균등분포가 아니다]
\texttt{serving\_size}가 균등분포인지
경험적 CDF와의 최대 세로 격차(\Cref{sec:cdf}의 $D$)로 확인하면
\textbf{0.152 대 0.101}이다. 균등분포로 맞췄을 때 최대격차가 $0.152$로,
정규분포로 맞췄을 때의 $0.101$보다 크다. 즉 \texttt{serving\_size}는
균등분포보다 정규분포에 더 가깝다.
\end{pitfall}

% =====================================================================
\section{베르누이와 이항분포}
\label{sec:binomial}

\begin{inbrief}
  \item[Why] 결과가 성공/실패 둘뿐인 시행 하나(베르누이), 또는 그런
             시행을 여러 번 반복했을 때의 성공 횟수(이항)를 표현하기
             위해.
  \item[When] 이진 결과(예/아니오, 성공/실패)를 셀 때.
  \item[Where] 동전 던지기, 지지/반대 설문, 패스트푸드
               \texttt{sodium}이 기준치를 넘는지 여부.
  \item[How] 베르누이는 성공확률 $p$ 하나. 이항은 여기에 시행 횟수
             $n$을 더해 $n$번 중 성공 횟수의 분포를 구한다.
\end{inbrief}

% source: deck 6 slide 23
베르누이분포(Bernoulli Distribution)는 각 데이터를 이진값(Binary Value)
하나로 매핑할 수 있을 때 쓴다. 실험 결과가 성공 또는 실패뿐이면, 성공(값
1)에는 확률 $p$를, 실패(값 0)에는 확률 $1-p$를 준다. 동전의 앞/뒤,
찬성/반대가 그 예다.

% source: deck 6 slide 24
\begin{example}[지지율 설문]
\label{ex:dist-favorability}
설문 응답자의 $60\%$가 찬성($1$), $40\%$가 반대($0$)라면 평균과 분산은
다음과 같다.
\begin{align*}
  \mu &= 0.4\times 0 + 0.6\times 1 = \mathbf{0.6}
    && \text{평균, 그러나 분포가 실제로 가질 수 있는 값은 0과 1뿐}\\
  \sigma^2 &= 0.4(0-0.6)^2 + 0.6(1-0.6)^2 && \text{분산: 편차제곱의 기댓값}\\
           &= 0.144 + 0.096 = \mathbf{0.24} && \text{정리}
\end{align*}
\end{example}

% source: deck 6 slide 25
일반적으로 성공확률 $p$인 베르누이분포의 평균과 표준편차는
\[
  \E[X] = p, \qquad \sd(X) = \sqrt{p(1-p)}
\]
이고, 확률질량함수는 $P(X=1)=p$, $P(X=0)=1-p$다(슬라이드 25의
표기와 같다).

% source: deck 6 slide 26
이항분포(Binomial Distribution)는 베르누이 시행을 서로 \textbf{독립}으로
$n$번 반복했을 때, 성공 횟수 $X$의 분포다. 예를 들어 동전을 5번 던져
``뒷면'' 2번이 나올 확률을 구하는 문제가 여기 속한다.

% source: deck 6 slide 27
\begin{example}[동전 5번 중 순서까지 정해진 경우]
\label{ex:dist-coin-order}
동전을 5번 던져 앞의 3번은 실패(성공=뒷면=1, 실패=앞면=0), 뒤의 2번은
성공인 \textbf{정해진 순서 하나}가 나올 확률은
\[
  0.5^5 = \mathbf{0.03125}
\]
다(각 시행이 독립이므로 곱한다, \Cref{sec:ht-background}).
\end{example}

% source: deck 6 slide 28
\begin{example}[순서를 가리지 않는 경우]
\label{ex:dist-coin-combination}
성공이 정확히 2번 나오는 \textbf{모든} 순서를 세려면 조합
$\binom{5}{2}$을 곱한다.
\begin{align*}
  \binom{5}{2} &= 10                      && \text{5번 중 2번을 고르는 방법 수}\\
  P(X=2) &= \binom{5}{2}\times 0.5^5      && \text{순서 무관 확률}\\
         &= 10/32                          && \text{값 대입}\\
         &= \mathbf{0.3125}                && \text{정리}
\end{align*}
\end{example}

% source: deck 6 slide 29
일반적으로 $n$번의 독립 베르누이 시행에서 성공확률 $p$인 이항분포의
평균과 표준편차는
\[
  \E[X] = np, \qquad \sd(X) = \sqrt{np(1-p)}
\]
다(슬라이드 29의 표기와 같다).

\Cref{fig:dist-bernoulli-binomial}이 패스트푸드 \texttt{sodium}이
$1000\,\text{mg}$을 넘는지 여부를 베르누이 시행으로, $n=10$번 뽑았을 때
고나트륨 항목 개수를 이항분포로 보여준다.

\begin{figure}[ht]\centering
  \includegraphics[width=0.95\linewidth]{06_07_bernoulli_binomial.png}
  \caption{패스트푸드 \texttt{sodium} $>1000\,\text{mg}$ 여부의
  베르누이 PMF($p=0.4524$)와, $n=10$번 뽑았을 때 고나트륨 개수의
  이항 PMF(평균 $np=4.52$, 표준편차 $1.57$).}
  \label{fig:dist-bernoulli-binomial}
\end{figure}

% =====================================================================
\section{누적분포함수로 비교하기}
\label{sec:cdf}
% source: deck 6 slide 16-17 (CDF as tool); figure from hw1 data

\begin{inbrief}
  \item[Why] 히스토그램 없이도, 분포 전체를 한 곡선으로 겹쳐 비교하기
             위해.
  \item[When] 데이터가 특정 분포를 따르는지 눈으로 먼저 확인하고 싶을 때.
  \item[Where] 지수분포 CDF로 대기 시간을 구했던 계산(슬라이드 16, 17)이
               이미 CDF를 도구로 쓴 예다.
  \item[How] 경험적 CDF(계단 모양)와 이론 CDF(매끄러운 곡선)를 같은
             축에 겹치고, 가장 크게 벌어지는 지점을 본다.
\end{inbrief}

값 $x$에서의 \term{누적분포함수}(Cumulative Distribution Function, CDF)는
한 줄로 요약하면 $F(x) = P(X \le x)$, 즉 ``$x$ 이하가 전체에서 차지하는
비율''이다. 데이터로 만든 계단 모양의 CDF를
\term{경험적 누적분포함수}(Empirical CDF, $\hat F$)라 부르며, 정식
정의는 \Cref{def:ecdf}(\Cref{ch:compare-dist})가 맡는다. 이 장은 CDF를
지수분포 예제(\Cref{ex:dist-laptop}, \Cref{ex:dist-coffee})에서 이미
계산 도구로 썼으므로, 정의를 다시 쓰지 않고 두 곡선을 겹쳐 보는 용도로만
쓴다.
모집단 CDF의 정식 정의(단조 증가, 우연속, 극한 0과 1)는 \Cref{def:cdf}에 있다.

\Cref{fig:dist-cdf}가 패스트푸드 \texttt{calories}의 경험적 CDF와
정규분포로 맞춘 CDF를 겹쳐 보여준다.

\begin{figure}[ht]\centering
  \includegraphics[width=0.85\linewidth]{06_03_cdf.png}
  \caption{패스트푸드 \texttt{calories}의 경험적 CDF(계단)와 최대우도로
  맞춘 정규 CDF($\hat\mu=532.49$, $\hat\sigma=249.85$). 두 곡선의 최대
  세로 격차 $D=0.1007$이 표시되어 있다.}
  \label{fig:dist-cdf}
\end{figure}

두 곡선이 가장 크게 벌어지는 세로 거리가
\term{콜모고로프-스미르노프 통계량}(Kolmogorov-Smirnov Statistic) $D$다.
정식 정의와 $D$를 이용한 검정 절차는 \Cref{def:ks-statistic}과
\Cref{sec:cd-ks}(\Cref{ch:compare-dist})가 다룬다. 같은 표본에서 모수를
추정한 뒤 이 표본으로 KS 검정을 하면 p값이 낙관적으로 나오므로
(\Cref{sec:normal}의 함정), 이 장에서는 $D$를 \textbf{눈으로 보는} 도구로만
쓰고 p값은 계산하지 않는다. 히스토그램 대신 CDF로 비교하면 구간 폭을
정할 필요가 없다는 장점은 \Cref{sec:dist-check-cdf}(\Cref{ch:qq-normalization})가
더 다룬다.

% =====================================================================
\section{요약}

\Cref{tab:dist-summary}가 이 장에서 다룬 분포를 모수, 평균, 분산과 함께
정리한다.

\begin{table}[ht]\centering
\begin{tabularx}{\linewidth}{@{}l l X c c c@{}}
\toprule
분포(Distribution) & 연속/이산 & 모수 & 평균 & 분산 & 절 \\
\midrule
정규(Normal)     & 연속 & $\mu, \sigma^2$ & $\mu$ & $\sigma^2$
  & \Cref{sec:normal} \\
지수(Exponential)     & 연속 & $\lambda$ & $1/\lambda$ & $1/\lambda^2$
  & \Cref{sec:exponential} \\
멱법칙(Power Law)   & 연속 & $\alpha, x_{\min}$ & (조건부 존재) & (조건부 존재)
  & \Cref{sec:powerlaw} \\
균등(Uniform)     & 연속 & $a, b$ & $(a+b)/2$ & $(b-a)^2/12$
  & \Cref{sec:uniform} \\
베르누이(Bernoulli) & 이산 & $p$ & $p$ & $p(1-p)$
  & \Cref{sec:binomial} \\
이항(Binomial)     & 이산 & $n, p$ & $np$ & $np(1-p)$
  & \Cref{sec:binomial} \\
\bottomrule
\end{tabularx}
\caption{이 장에서 다룬 분포의 모수와 평균, 분산.}
\label{tab:dist-summary}
\end{table}

% =====================================================================
\section{어떻게 적용하는가}
\label{sec:dist-apply}
% source: deck 6 slide 30 (내용 없음, Kahoot 퀴즈)

이 책 안에서 이 장의 분포가 실제로 쓰이는 자리는 다음과 같다.

\begin{itemize}
  \item 패스트푸드 \texttt{calories}: 히스토그램은 종 모양이지만 정규성이
        기각된다(\Cref{fig:normal-fit}, Shapiro-Wilk $p=0.00038$,
        \Cref{sec:normal}의 함정).
  \item 패스트푸드 \texttt{sugars}: 지수분포로 맞추면
        $\hat\lambda = 0.0753$(\Cref{fig:dist-exponential}).
  \item 패스트푸드 \texttt{sodium}: $1000\,\text{mg}$ 초과 여부를
        베르누이/이항으로 요약한다(\Cref{fig:dist-bernoulli-binomial}).
  \item HW2의 공항 노선 수, 영화 평점: 네 후보 분포(멱법칙, 지수, 균등,
        정규)의 모수를 추정하고 적합도를 판정한다
        (\Cref{sec:hw2-estimators}, \Cref{fig:hw2-airport-powerlaw}).
  \item 후보를 고른 뒤에는 Q-Q 플롯으로 한 번 더 확인한다
        (\Cref{sec:qq}, \Cref{fig:qq-calories}).
\end{itemize}

% =====================================================================
\section{앞으로}
\label{sec:dist-future}

\begin{itemize}
  \item \textbf{모수 추정과 Q-Q 플롯.} 이 장은 슬라이드의 공식으로
        모수를 구했다. 최대우도추정(Maximum Likelihood Estimation, MLE)의
        일반적인 절차와 Q-Q로 적합을
        확인하는 법은 \Cref{ch:qq-normalization}이 다룬다.
  \item \textbf{이산분포 위의 정확 검정.} 이항분포에서 관측값이 얼마나
        극단적인지 정확히 계산하는 법은 \Cref{sec:ht-exact}와
        \Cref{sec:pv-discrete}가 다룬다.
  \item \textbf{정규분포 위의 신뢰구간.} $\pm 1.96\sigma$가 왜 $95\%$의
        경계인지, 신뢰구간에 어떻게 쓰이는지는 \Cref{sec:ci}가 다룬다.
  \item \textbf{두 분포를 통계적으로 비교하기.} 이 장은 CDF를 눈으로만
        겹쳐 봤다. 숫자 하나(KS 통계량 $D$)로 ``얼마나 다른가''를
        검정하는 절차는 \Cref{sec:cd-ks}가 다룬다.
\end{itemize}

\chapter{가설검정: 증거가 귀무가설을 터무니없게 만드는가}
\label{ch:hypothesis}
% source: deck 9 slide 1

\begin{inbrief}
  \item[Why] \term{가설검정}(Hypothesis Testing)은 표본 하나에서 본 차이가
             진짜인지, 우연히 그렇게 나온 것인지를 가르는 규칙이다.
  \item[When] "차이가 없다"(또는 기존 값 그대로다)를 기본값으로 두고, 데이터가
              그 기본값을 뒤집을 만큼 강한 증거인지 물을 때.
  \item[Where] 두 약의 회복 시간 비교, 비율(수컷 강아지 비율, 반항적 인간
               비율), 성공 확률(복권)처럼 \term{모수}(Parameter)에 대한
               주장을 검증할 때.
  \item[How] \term{귀무가설}(Null Hypothesis) $H_0$ 와 \term{대립가설}
             (Alternative Hypothesis) $H_a$ 를 세운다 $\to$ 기준
             $\alpha$ 를 미리 정한다 $\to$ $H_0$ 가 참일 때 관측값 이상으로
             극단적인 결과가 나올 확률(\term{p값}, p-value)을 구한다 $\to$
             기준보다 작으면 \term{기각}(Reject), 아니면
             \term{기각 실패}(Fail to Reject).
\end{inbrief}

\begin{keyidea}
가설검정은 \analogy{현장 측량으로 지도 검증}이다(\Cref{tab:analogy}). 지도
(귀무가설)가 맞다면 이만큼 어긋난 측량이 나올 법한가를 묻는다. 너무 어긋나면
지도를 고치고(기각), 오차 안이면 고칠 근거가 없을 뿐 지도가 옳다는 증명은
아니다(기각 실패).
\end{keyidea}

이 장의 출처는 \texttt{notes/9\_Hypothesis Testing.pptx}(14 슬라이드) 하나다.
슬라이드가 외부 문헌을 밝히지 않으므로 슬라이드 자체를 출처로 적는다.

% =====================================================================
\section{배경}
\label{sec:ht-background}
% source: deck 9 slide 1

이 장을 읽는 데 필요한 개념을 먼저 깐다. 이미 다룬 것은 한 줄 요약과 링크로,
처음 나오는 것은 그 자리에서 정의한다.

\term{평균}(Mean)과 \term{퍼짐}(Spread, 표준편차)은
\Cref{sec:center,sec:spread}에서 이미 다뤘다. 이 장의 두 집단 비교
(\Cref{sec:ht-two-groups})에서 "평균에서 얼마나 떨어져 있는가"를 잴 때 그대로
쓴다.

\term{표본}(Sample)과 \term{모집단}(Population)은 이 장에서 처음 정의한다.
우리가 실제로 손에 쥐고 보는 것은 \textbf{표본}이고, 우리가 주장하려는 대상은
\textbf{모집단}이다. 강아지 15마리를 관찰한 것은 표본이고, "이 부모견이
낳는 강아지 전체의 성비"가 모집단이다. \term{모수}(Parameter)는 그 모집단을
숫자 하나로 나타낸 값이다(예: 모집단에서 수컷일 확률 $p$). 표본에서 계산한
값으로 모수를 추측하고, 그 추측이 미리 세운 주장과 맞는지를 검정한다.

\term{확률의 곱셈 규칙}(Multiplication Rule of Probability)도 이 장에서
처음 쓴다. \term{독립 시행}(Independent Trial, 사건 하나가 다른 사건에
영향을 주지 않는 시행)이 연달아 일어날 확률은 각 확률을
\textbf{곱한다}. 동전을 두 번 던져 둘 다 앞면이 나올 확률은
$0.5 \times 0.5 = 0.25$ 다. \Cref{sec:ht-exact}의 $0.4^3$ 계산이 이 규칙
그대로다. \term{이항분포}(Binomial Distribution)는 이 곱셈 규칙을 "$n$ 번
시행 중 성공이 정확히 $k$ 번" 이라는 경우의 수까지 센 분포다
(\Cref{sec:binomial}, 현재는 뼈대만 있는 절이며 \Cref{ch:distributions}에서
채워진다). \Cref{sec:ht-simulation}의 이항계수 계산이 이 분포를 그대로
쓴다.

\term{신뢰구간}(Confidence Interval)과 "5\% 미만이면 유의하다"는 기준은
\Cref{sec:ci,subsec:ci-compare}에서 이미 나왔다. \term{부트스트랩}
(Bootstrap, \Cref{subsec:bootstrap})의 "컴퓨터로 반복해서 분포를 만든다"는
발상은 \Cref{sec:ht-simulation}의 시뮬레이션과 같은 발상이다. 차이는
\Cref{sec:ht-apply}에서 다시 짚는다.

\begin{pitfall}[$p$ 라는 기호는 두 가지 뜻이 있다]
슬라이드 11에서 $p$ 는 "수컷일 확률"이라는 \textbf{모수}다
($p = 50\%$). 그런데 통계학 일반에서 $p$ 는 \textbf{p값}(p-value)도 가리킨다.
이 책에서는 슬라이드의 기호를 그대로 두고(모수는 $p$), \textbf{p값은 항상
"p값"이라는 한글 단어}로만 쓴다. 덱 9 자신은 "p-value"나 "significance
level"이라는 말을 쓰지 않고 "significance test"만 쓰지만, 이 장이 계산하는
양(슬라이드 11, 12의 "13 이상일 확률")은 \Cref{sec:ci}가 이미 쓴 p값과
같은 개념이다.
\end{pitfall}

\Cref{tab:ht-symbols}이 이 장에서 쓰는 기호를 모은다.

\begin{table}[ht]\centering
\begin{tabularx}{\linewidth}{@{}l X l@{}}
\toprule
기호 & 뜻 & 읽기 \\
\midrule
$H_0$ & \term{귀무가설}(Null Hypothesis): 차이 없음, 기존 값 그대로 & 에이치 제로 \\
$H_a$ & \term{대립가설}(Alternative Hypothesis): 찾으려는 주장 & 에이치 에이 \\
$p$ & 모수로서 성공 확률(예: $p = 50\%$). p값과는 다르다(위 함정 참고) & 피 \\
$\alpha$ & \term{유의수준}(Significance Level): 미리 정한 기각 기준(이 장은
1\%, \Cref{sec:ci}는 5\%) & 알파 \\
$n$, $X$ & 시행 횟수, 성공 횟수 & \\
\bottomrule
\end{tabularx}
\caption{이 장에서 쓰는 기호.}
\label{tab:ht-symbols}
\end{table}

% =====================================================================
\section{왜 필요한가}
\label{sec:ht-why}
% source: deck 9 slide 3-5, 7, 8

\begin{inbrief}
  \item[Why] 실험 한 번의 결과로 세운 가설은 다음 실험에서 뒤집히기 쉽다.
  \item[When] 재현되지 않는 결과를 어떻게 판정할지 정해야 할 때.
  \item[Where] 예비 자료(Preliminary Data)로 가설을 세우는 모든 실험.
  \item[How] 예비 자료가 필요 없는 유일한 값, 0(차이 없음)을 가설로 삼는다.
\end{inbrief}

같은 약을 먹어도 사람마다 회복 시간이 다르다. 예비 실험에서 약 A가 약 B보다
평균 15시간 더 빨리 회복시켰다면(슬라이드 3), "A는 B보다 15시간 빠르다"는
가설을 세우고 싶어진다. 문제는 \textbf{그 다음}이다.

\begin{enumerate}
  \item \textbf{재실험이 정반대로 나온다(슬라이드 4).} 다시 실험했더니 이번에는
        A가 오히려 35시간 \emph{더 느렸다}. 방향 자체가 뒤집혔으므로, "15시간
        빠르다"는 예비 가설은 \textbf{확신을 갖고 기각}한다. 이 기각은 아직
        귀무가설이 아니라 예비 자료로 세운 가설을 기각한 것이다(귀무가설은
        슬라이드 6에서 처음 정의된다).
  \item \textbf{재실험이 애매하게 나온다(슬라이드 5).} 약 C가 D보다 13시간
        빠르다는 가설은 첫 실험이 13시간이었다는 이유 하나로 정해졌다.
        재실험은 12시간, 13.5시간으로 나왔다. 기각할 만큼 다르지도, 맞다고
        확신시킬 만큼 같지도 않다. 결론은 \textbf{기각 실패}다.
  \item \textbf{방향이 실험마다 뒤집힌다(슬라이드 7).} 약 E와 F의 차이는
        한 번은 F가 0.25시간 높고, 다음번은 E가 0.5시간 높다. 차이의
        \emph{방향}조차 재현되지 않으므로 "차이가 있다"고 말할 근거가 없다.
        기각 실패다.
\end{enumerate}

\Cref{fig:ht-drug-experiments}가 첫 번째와 세 번째 경우를 나란히 보여준다.

\begin{figure}[ht]\centering
  \includegraphics[width=\linewidth]{09_01_drug_experiments.png}
  \caption{반복 실험이 가설을 지지하는가(슬라이드 3, 4, 7). 왼쪽: 예비
  실험이 세운 "A가 15시간 빠르다"는 가설을, 방향이 뒤집힌 재실험(35시간)이
  확신을 갖고 기각시킨다. 오른쪽: 약 E/F는 재실험마다 어느 쪽이 높은지가
  뒤집혀 기각 실패로 끝난다. 슬라이드 이미지의 점 위치가 축 눈금 라벨과
  어긋나 있어, 시간 값은 라벨의 차이만 재현한 책 자체의 개략(schematic) 값이다.}
  \label{fig:ht-drug-experiments}
\end{figure}

이 세 경우가 드러내는 문제는 하나다. \textbf{예비 자료로 세운 가설의 숫자
(15시간, 13시간)는 그 자체가 우연일 수 있고}, 그러면 무엇을 기준으로
"충분히 다르다"를 판정할지가 매번 달라진다.

처방은 슬라이드 8이 준다. 예비 자료가 필요 없는 유일한 값을 가설로 삼는다.
그 값은 \textbf{0}, 즉 "차이가 없다"다. 이것이 \term{귀무가설}(Null
Hypothesis)이다. 예비 실험에서 어떤 숫자가 나오든 귀무가설은 항상 같은
문장이므로, 판정 기준을 매번 새로 만들 필요가 없다.

% =====================================================================
\section{무엇인가}
\label{sec:ht-hypotheses}

\begin{inbrief}
  \item[Why] 증거를 잴 기준점(차이 0)이 필요하다.
  \item[When] 데이터를 보기 \textbf{전에} 세운다.
  \item[Where] 두 집단 비교, 비율, 확률에 대한 주장.
  \item[How] $H_0$ 는 등호("차이 없음", "모수 = 특정 값"), $H_a$ 는
             찾으려는 주장.
\end{inbrief}
% source: deck 9 slide 2, 6, 8, 9

\begin{definition}[귀무가설]
\label{def:ht-null}
\term{귀무가설}(Null Hypothesis) $H_0$ 는 "차이가 없다" 또는 "모수가 특정
값과 같다"는 기본 가정이다. 예비 자료 없이 세울 수 있는 유일한 값이 0이기
때문에 항상 등호($=$)로 쓴다.
\end{definition}

\begin{definition}[대립가설과 유의성 검정]
\label{def:ht-alt}
\term{대립가설}(Alternative Hypothesis) $H_a$ 는 우리가 증거로 뒷받침하려는
주장이다. \term{유의성 검정}(Significance Test)은 $H_0$ 가 참이라는 가정
아래에서, 관측한 증거가 $H_0$ 를 얼마나 터무니없게 만드는지를 재는 절차다.
\end{definition}

금성에 외계 생명체가 있는지 찾는 예(슬라이드 2)가 이 구조를 그대로 보여준다.
$H_0$: 이 행성에는 외계 생명체가 없다. $H_a$: 외계 생명체가 있다. 만약
증거를 전혀 찾지 못하면, $H_0$ 를 기각할 근거가 없으므로 \textbf{기각 실패}로
남는다. 증거를 찾으면 그때 $H_0$ 를 기각한다. "아무것도 찾지 못했다"는
"생명체가 없다는 증명"이 아니라 "기각할 증거가 없었다"는 뜻이다.

대부분의 경우 $H_a$ 는 $H_0$ 의 반대다(슬라이드 9). 약 C와 D의 예에서
$H_0$: 회복 시간에 차이가 없다, $H_a$: 차이가 있다. 이 대립가설을 "평균에서의
거리"로 판정하는 방법은 \Cref{sec:ht-two-groups}에서 다룬다.

\subsection{기각과 기각 실패}
\label{sec:ht-decision}

\begin{inbrief}
  \item[Why] 결론은 두 가지뿐이고, "채택"은 없다.
  \item[When] 증거의 확률을 기준 $\alpha$ 와 비교한 직후.
  \item[Where] 모든 검정의 마지막 문장.
  \item[How] 증거가 $H_0$ 를 터무니없게 만들면 기각, 아니면 "기각
             실패"라고만 쓴다.
\end{inbrief}
% source: deck 9 slide 2, 4, 5, 7

가설검정의 결론은 딱 두 가지다. \textbf{기각}(Reject $H_0$) 아니면
\textbf{기각 실패}(Fail to Reject $H_0$). "$H_0$ 를 채택한다"거나 "$H_0$ 가
증명됐다"는 결론은 없다.

\begin{pitfall}[기각 실패는 $H_0$ 가 참이라는 증명이 아니다]
\label{pit:ht-fail-not-proof}
금성에서 외계 생명체를 못 찾았다고 해서 생명체가 없다는 뜻은 아니다
(슬라이드 2). 약 C/D 예에서도 슬라이드는 "재실험이 가설이 맞다고 확신시키지도
않는다"(don't convince us the hypothesis is correct, 슬라이드 5)고 분명히
적는다. 이것은 \Cref{subsec:ci-compare}의 함정 "겹친다고 같은 것은
아니다"와 같은 논리다. \textbf{증거 부족은 무죄를 증명하지 않는다}, 단지
유죄를 선고할 근거가 없을 뿐이다.
\end{pitfall}

\Cref{tab:ht-wording}은 올바른 결론 문장과 틀린 결론 문장을 나란히 둔다.

\begin{table}[ht]\centering
\begin{tabularx}{\linewidth}{@{}X X@{}}
\toprule
올바른 표현 & 틀린 표현 \\
\midrule
"$H_0$ 를 기각한다" & "$H_0$ 가 틀렸다고 증명했다" \\
"$H_0$ 를 기각하지 못한다" & "$H_0$ 를 채택한다" \\
"기각할 증거가 부족하다" & "$H_0$ 가 옳다고 증명했다" \\
\bottomrule
\end{tabularx}
\caption{가설검정 결론의 올바른 표현과 틀린 표현.}
\label{tab:ht-wording}
\end{table}

\subsection{대립가설의 세 형태}
\label{sec:ht-forms}

\begin{inbrief}
  \item[Why] 무엇을 "극단적"이라 부를지는 $H_a$ 의 방향이 정한다.
  \item[When] $H_0$ 를 세울 때 함께, 데이터를 보기 전에.
  \item[Where] "더 크다"(슬라이드 11--13), "다르다"(슬라이드 9).
  \item[How] $>$, $<$ 는 한쪽 꼬리(\term{단측검정}, One-sided Test), $\neq$
             는 양쪽 꼬리(\term{양측검정}, Two-sided Test).
\end{inbrief}
% source: deck 9 slide 9, 10

$H_0$ 는 항상 모수가 특정 값과 같다는 형태다: $H_0: \text{모수} = \text{값}$.
$H_a$ 는 다음 셋 중 하나다.

\begin{center}
\begin{tabular}{@{}lll@{}}
\toprule
$H_a$ & 이름 & 극단값의 방향 \\
\midrule
모수 $>$ 값 & 단측검정(오른쪽 꼬리) & 관측값보다 큰 쪽 \\
모수 $<$ 값 & 단측검정(왼쪽 꼬리) & 관측값보다 작은 쪽 \\
모수 $\neq$ 값 & 양측검정 & 양쪽 다 \\
\bottomrule
\end{tabular}
\end{center}

\Cref{fig:ht-tails}가 세 형태의 꼬리를 보여준다.

\begin{figure}[ht]\centering
\begin{tikzpicture}[scale=1.0]
  \begin{scope}
    \draw[domain=-3:3,samples=60,smooth,draw=bookink] plot (\x,{exp(-\x*\x/2)});
    \fill[bookwarnbg] plot[domain=1.5:3,samples=30] (\x,{exp(-\x*\x/2)}) -- (3,0) -- (1.5,0) -- cycle;
    \draw[bookmuted] (-3,0) -- (3,0);
    \node[below] at (0,-0.15) {$>$: 오른쪽 꼬리};
  \end{scope}
  \begin{scope}[xshift=4.3cm]
    \draw[domain=-3:3,samples=60,smooth,draw=bookink] plot (\x,{exp(-\x*\x/2)});
    \fill[bookwarnbg] plot[domain=-3:-1.5,samples=30] (\x,{exp(-\x*\x/2)}) -- (-1.5,0) -- (-3,0) -- cycle;
    \draw[bookmuted] (-3,0) -- (3,0);
    \node[below] at (0,-0.15) {$<$: 왼쪽 꼬리};
  \end{scope}
  \begin{scope}[xshift=8.6cm]
    \draw[domain=-3:3,samples=60,smooth,draw=bookink] plot (\x,{exp(-\x*\x/2)});
    \fill[bookwarnbg] plot[domain=1.5:3,samples=30] (\x,{exp(-\x*\x/2)}) -- (3,0) -- (1.5,0) -- cycle;
    \fill[bookwarnbg] plot[domain=-3:-1.5,samples=30] (\x,{exp(-\x*\x/2)}) -- (-1.5,0) -- (-3,0) -- cycle;
    \draw[bookmuted] (-3,0) -- (3,0);
    \node[below] at (0,-0.15) {$\neq$: 양쪽 꼬리};
  \end{scope}
\end{tikzpicture}
\caption{대립가설의 방향이 정하는 꼬리(개략도, 축 값 없음). 색칠된 영역이
"극단적"으로 취급되는 쪽이다.}
\label{fig:ht-tails}
\end{figure}

\begin{pitfall}[방향은 데이터를 보고 고르지 않는다]
\label{pit:ht-direction}
방향($>$, $<$, $\neq$)은 데이터를 보기 \textbf{전에} 정한다. 관측값을 본
뒤에 "마침 이쪽으로 컸으니 $>$ 로 하자"고 고르면, 우연히 큰 쪽으로 나온
결과를 스스로 유리하게 골라 기각을 더 쉽게 만든 것이다. \Cref{sec:ht-simulation}
연습 문제(반항적 인간 14\%)에서 방향을 데이터를 본 뒤에 고른 사례
(\Cref{pit:ht-alien})가 이 함정의 구체적인 예다.
\end{pitfall}

\subsection{p값과 유의수준}
\label{sec:ht-pvalue}

\begin{inbrief}
  \item[Why] "터무니없다"를 숫자 하나로 바꿔야 판정이 재현된다.
  \item[When] 가설과 방향을 정한 뒤.
  \item[Where] 슬라이드 11--13이 모두 이 양을 계산한다.
  \item[How] $H_0$ 가 참이라 가정하고 "관측값 이상으로 극단적인 결과"의
             확률을 구해 $\alpha$ 와 비교한다.
\end{inbrief}
% source: deck 9 slide 11, 12, 13

\begin{definition}[p값]
\label{def:pvalue}
\term{p값}(p-value)은 $H_0$ 가 참이라고 가정했을 때, 관측한 값 또는 그보다
$H_a$ 방향으로 더 극단적인 값이 나올 확률이다.
\end{definition}

p값이 미리 정한 유의수준 $\alpha$ 보다 \textbf{작으면} $H_0$ 를 기각한다.
그렇지 않으면 기각 실패다.

\begin{pitfall}[기준 $\alpha$ 는 법칙이 아니라 약속이다]
\label{pit:ht-alpha}
이 장은 $\alpha = 1\%$ 를 쓰고(슬라이드 11--13), \Cref{sec:ci}는
$\alpha = 5\%$($z = 1.96$)를 쓴다. 둘 다 맞다. $\alpha$ 는 자연법칙이 아니라
\textbf{데이터를 보기 전에 정하는 약속}이다. \Cref{subsec:fence}의
"$1.5 \times \IQR$ 은 관례다"와 같은 구조의 함정이다. 기준이 낮을수록
기각하기 어려워진다.
\end{pitfall}

\subsection{시뮬레이션으로 p값 구하기}
\label{sec:ht-simulation}

\begin{inbrief}
  \item[Why] 공식을 몰라도 "$H_0$ 세상"을 컴퓨터로 여러 번 만들어 보면 된다.
  \item[When] 분포 공식이 없거나 복잡할 때.
  \item[Where] 강아지 성비, 반항적 인간 비율.
  \item[How] $H_0$ 로 1000회 모의 $\to$ 관측값 이상인 횟수를 센다 $\to$
             1000으로 나눈다.
\end{inbrief}
% source: deck 9 slide 11, 12

\begin{example}[강아지 15마리 중 수컷 13마리]
\label{ex:ht-puppy}
엘레나의 개 한 쌍이 강아지 15마리를 낳았고, 그중 13마리가 수컷이었다
(슬라이드 11).

\begin{align*}
  &H_0: p = 0.5 && \text{암수 확률이 같다}\\
  &H_a: p > 0.5 && \text{수컷일 확률이 더 크다}\\
  &\alpha = 1\%
\end{align*}

\textbf{1단계. 표에서 직접 센다.} 슬라이드는 $H_0$($p=0.5$) 아래에서 강아지
15마리를 1000번 시뮬레이션한 표를 준다(\Cref{tab:ht-puppy}, 빈도 열).
$k \ge 13$(관측값 이상)인 행을 더한다.

\begin{align*}
  &3 + 0 + 0 = 3 && \text{$k=13,14,15$ 의 빈도 합}\\
  &3 / 1000 = 0.003 && \text{시뮬레이션으로 구한 p값}
\end{align*}

\textbf{2단계. 이항계수로 정확히 구한다.} \term{이항계수}(Binomial
Coefficient) $\binom{n}{k}$ 는 "$n$ 개 중 $k$ 개를 고르는 경우의 수"라고
읽는다(순서는 따지지 않는다). 예를 들어 $\binom{15}{13}$ 은 강아지 15마리
중 수컷이 정확히 13마리인 조합이 몇 가지인지를 센 값이다. 이항분포
(\Cref{sec:binomial})를 쓰면 시뮬레이션 없이 정확한 값을 계산할 수 있다.

\begin{align*}
  &\binom{15}{13}+\binom{15}{14}+\binom{15}{15} = 105+15+1 = 121
    && \text{$k \ge 13$ 인 경우의 수}\\
  &121/32768 = 0.0037 && \text{$2^{15}=32768$ 로 나눈 정확한 p값}
\end{align*}

\textbf{3단계. 두 값이 다른 이유.} 시뮬레이션 값(0.003)과 정확한 값
(0.0037)이 완전히 같지 않다. \Cref{tab:ht-puppy}의 정확값 열을 보면 이유가
드러난다. 표의 모든 행이 $1000\binom{15}{k}/2^{15}$ 를 \textbf{반올림한}
값과 정확히 일치한다(예: $k{=}7$: $196.4 \to 196$, $k{=}14$:
$0.46 \to 0$). \term{기댓값}(Expected Value)은 "여러 번 반복했을 때 평균적으로
나올 것으로 예상되는 값"이다. 여기서는 $1000$ 번 시행했을 때 $k$ 명이 수컷일
경우의 수 평균, 즉 $1000\binom{15}{k}/2^{15}$ 를 뜻한다. 즉 이 표는
\textbf{실제 무작위 시뮬레이션 결과가 아니라 기댓값을 반올림해 적은 표}다.

\begin{pitfall}[슬라이드 표는 기댓값 표다]
\label{pit:ht-r4}
슬라이드 값은 "1000회 시뮬레이션 결과"라고 적혀 있지만, 다시 계산하면
정확히 $1000\binom{15}{k}/2^{15}$ 의 반올림값과 일치하는
\textbf{기댓값 표}다. 실제 무작위 시뮬레이션을 돌리면 이 값 근처에서
흔들리며 $k{=}14,15$ 행이 항상 0으로 반올림되지는 않는다. \Cref{fig:ht-puppy-sim}의
점이 정확값, 막대가 슬라이드 표다.
\end{pitfall}

$0.0037 < 1\%$ 이므로 \textbf{기각한다}. 결론: "이 강아지들에서 수컷이 더
잘 나온다는 증거가 있다."
\end{example}

\begin{table}[ht]\centering
\begin{tabularx}{\linewidth}{@{}c r r@{}}
\toprule
$k$ (수컷 수) & 슬라이드 빈도(1000회 중) & 정확값 $1000\binom{15}{k}/2^{15}$ \\
\midrule
6  & 153 & 152.7 \\
7  & 196 & 196.4 \\
8  & 196 & 196.4 \\
9  & 153 & 152.7 \\
10 & 92  & 91.6 \\
11 & 42  & 41.7 \\
12 & 14  & 13.9 \\
13 & 3   & 3.2 \\
14 & 0   & 0.46 \\
15 & 0   & 0.03 \\
\bottomrule
\end{tabularx}
\caption{강아지 예제(슬라이드 11)의 빈도표(일부, $k\ge 6$)와 정확값 비교.
슬라이드 값과 정확값이 반올림 오차 안에서 그대로 일치한다(\Cref{pit:ht-r4}).}
\label{tab:ht-puppy}
\end{table}

\begin{figure}[ht]\centering
  \includegraphics[width=\linewidth]{09_03_puppy_null.png}
  \caption{$H_0$: $p=0.5$ 아래 수컷 강아지 수의 분포(슬라이드 11). 막대는
  슬라이드의 빈도표, 점은 정확한 $1000\binom{15}{k}/2^{15}$ 다. 관측값
  13(세로 점선) 이상이 경고색이다.}
  \label{fig:ht-puppy-sim}
\end{figure}

\begin{pitfall}[연습: 반항적 인간 14\%는 경계가 애매하다]
\label{pit:ht-alien}
외계 제국은 반항적 인간 비율이 10\% 미만이어야 지구를 침공하지 않는다.
400명 표본에서 14\%가 반항적이었다(슬라이드 12). $H_0: p = 0.10$,
$\alpha = 1\%$.

슬라이드는 \Cref{tab:ht-alien}의 "14\% 이상" 행만 읽어
$7/1000 = 0.007 < 1\%$ 로 기각한다. 그런데 이 장은 세 가지를 다시 짚는다.

아래에서 \term{정규근사}(Normal Approximation), \term{연속성 보정}
(Continuity Correction), $z$ 값이 등장한다. \textbf{정규근사}는 이항분포를
매번 이항계수로 계산하는 대신, 표본이 충분히 크면(예: $n=400$) 종 모양의
정규분포로 어림잡아 빠르게 근사하는 방법이다. \textbf{연속성 보정}은
이항분포가 정수($0, 1, 2, \dots$)만 갖는 데 비해 정규분포는 모든 실수를
갖는다는 차이를 메우려고, 경계를 $0.5$ 만큼 늘려 잡는 보정이다(예: "$54$
이상"을 "$53.5$ 이상"으로 바꿔 근사한다). $z$ 는 관측값이 평균에서
표준편차 몇 개만큼 떨어져 있는지를 나타내는 수로, $z = (\text{관측값} -
\text{평균})/\text{표준편차}$ 로 구한다. 이항분포의 표준편차는
$\sqrt{np(1-p)}$ 다. 예를 들어 $n=400$, $p=0.10$ 이면
$\sqrt{400 \cdot 0.1 \cdot 0.9} = 6$ 이다.

\begin{enumerate}
  \item \textbf{경계가 애매하다.} 표본 400명에서 비율은 $1/400 = 0.25\%$
        단위로만 나온다. "14\% 행"이 $[13.5\%,\,14.5\%)$ 를 묶은 것이라면,
        "14\% 이상"은 실제로는 $X \ge 54$명(13.5\%)을 뜻할 수도 있다. 연속성
        보정을 쓰면 $z = (53.5-40)/6 = 2.25$ 로, 정규근사 p값은 약
        1.22\%다.
        \textbf{슬라이드 값은 $7/1000 = 0.007$(인원수 $X \ge 56$ 기준)이나,
        경계를 $X \ge 54$ 로 다시 잡으면 정확한 이항 꼬리 확률은
        $P(X \ge 54) = 0.0149$ 로 다시 계산된다}(\texttt{scipy.stats.binom.sf(53,\allowbreak
        400,\allowbreak 0.1)}). 이 값은 1\% 기준을 \textbf{넘는다}, 즉
        판정이 뒤집힐 수 있다. 그래서 꼬리는 퍼센트가 아니라
        \textbf{인원수}($X \ge 56$)로 정의해야 모호함이 없다.
  \item \textbf{정확값도 확인해 둔다.} $X \sim \mathrm{Bin}(400, 0.1)$,
        평균 $40$, 표준편차 $\sqrt{400 \cdot 0.1 \cdot 0.9} = 6$. 연속성
        보정을 쓰면 $z = (55.5-40)/6 = 2.58$. 정확한 이항 꼬리 확률을
        보면, \textbf{슬라이드 값은 $0.007$이나, 다시 계산하면
        \texttt{scipy.stats.binom.sf(55, 400, 0.1)} $= 0.00664$}다. 결론
        자체(기각, $X\ge56$ 기준)는 바뀌지 않는다.
  \item \textbf{방향을 데이터 뒤에 고른 흔적이 있다.} 외계인의 결정 규칙은
        "10\% \textbf{미만}이면 침공하지 않는다"이므로, 그 규칙에 맞는
        질문은 $H_a: p < 0.10$ 이다. \textbf{슬라이드 값은 $H_a: p > 10\%$
        이나, 결정 규칙에 맞게 다시 세우면 $H_a: p < 10\%$ 가 되어야
        한다.} 슬라이드는 관측값 14\%(10\%보다 큰 쪽)를 본 뒤 $>$ 방향을
        택한 것으로 읽힌다. 이것이 \Cref{pit:ht-direction}이 경고한
        상황이다.
\end{enumerate}

세 가지 모두 계산이 틀린 것은 아니지만, 경계와 방향을 정하는 방식에 따라
결론이 흔들릴 수 있다는 연습 사례로만 쓴다. \Cref{fig:ht-alien-null}이
13.5\%(54명)와 14\%(56명) 경계가 만드는 차이를 그림으로 보여준다.
\end{pitfall}

\begin{table}[ht]\centering
\begin{tabularx}{\linewidth}{@{}c r r@{}}
\toprule
반항적 비율(\%, 묶음) & 슬라이드 빈도(1000회 중) & 위에서부터 누적(꼬리 합) \\
\midrule
6  & 7   & 1000 \\
7  & 40  & 993 \\
8  & 93  & 953 \\
9  & 173 & 860 \\
10 & 327 & 687 \\
11 & 253 & 360 \\
12 & 73  & 107 \\
13 & 27  & 34 \\
14 & 7   & 7 \\
\bottomrule
\end{tabularx}
\caption{반항적 인간 비율(슬라이드 12)의 빈도표와 누적 꼬리 합. "14" 행은
$[13.5\%,14.5\%)$ 를 묶은 값일 수 있어 경계가 애매하다(\Cref{pit:ht-alien}).}
\label{tab:ht-alien}
\end{table}

\begin{figure}[ht]\centering
  \includegraphics[width=\linewidth]{09_04_alien_null.png}
  \caption{$H_0$: $p=0.10$ 아래 반항적 인원수의 분포(슬라이드 12). 슬라이드
  표(막대, 인원수 축으로 환산)와 $\mathrm{Bin}(400,0.1)$ 이론 곡선을 겹쳤다.
  13.5\%(54명)와 14\%(56명) 경계선이 \Cref{pit:ht-alien}의 판정 차이를
  보여준다.}
  \label{fig:ht-alien-null}
\end{figure}

\subsection{정확한 계산}
\label{sec:ht-exact}

\begin{inbrief}
  \item[Why] 경우가 단순하면 시뮬레이션 없이 곱셈만으로 정확한 값이 나온다.
  \item[When] 독립 시행이고 확률이 알려져 있을 때.
  \item[Where] 복권 3연승, \Cref{ex:ht-puppy}의 이항 계산.
  \item[How] 독립이면 곱한다. 여러 경우면 이항계수로 센다.
\end{inbrief}
% source: deck 9 slide 13

\begin{example}[빨간 셔츠를 입으면 복권에 당첨될까]
\label{ex:ht-juan}
후안은 40\% 확률로 당첨되는 복권을 매일 산다. 빨간 셔츠를 입을 때 자주
당첨된다고 느껴, 빨간 셔츠를 입은 날만 기록했더니 3번 중 3번 당첨됐다
(슬라이드 13).

\begin{align*}
  &H_0: p = 0.4 && \text{셔츠는 효과가 없다}\\
  &H_a: p > 0.4 && \text{셔츠를 입으면 확률이 더 높다}\\
  &\alpha = 1\%
\end{align*}

세 번 모두 당첨은 독립 시행이므로 \Cref{sec:ht-background}의 곱셈 규칙을
그대로 쓴다.

\begin{align*}
  &0.4 \times 0.4 \times 0.4 = 0.064 && \text{$P(X=3)$, 즉 3번 모두 당첨될 확률}\\
  &0.064 = 6.4\%
\end{align*}

$H_a$ 가 "더 크다"이므로 p값은 $P(X \ge 3) = P(X=3) = 0.064$ 다(가능한 최댓값이
3이라 다른 경우가 없다). $6.4\%$ 는 $1\%$ 보다 크므로 \textbf{기각하지
못한다}.

\begin{keyidea}
$n=3$ 에서 나올 수 있는 가장 작은 p값이 $0.4^3 = 0.064$ 다. 즉 \textbf{이
표본 크기로는 어떤 결과가 나와도 1\% 기준으로 기각할 수 없다.} 연승이 몇
번이어야 기각되는지 같은 방식으로 확인한다.
\begin{align*}
  &0.4^5 = 0.01024 && \text{5연승, 아직 $1\%$ 를 넘음}\\
  &0.4^6 = 0.004096 && \text{6연승, 비로소 $1\%$ 미만}
\end{align*}
\textbf{6연승}이어야 이 기준으로 기각할 수 있다.
\end{keyidea}

결론: 기각 실패는 "셔츠 효과가 없다"는 증명이 아니라, \textbf{"이 표본
크기(3번)로는 말할 수 없다"}는 뜻이다. \Cref{subsec:notch}의 "노치가
뒤집히면 표본이 부족한 것"과 같은 신호다.
\end{example}

\subsection{두 집단 비교: 평균에서의 거리}
\label{sec:ht-two-groups}

\begin{inbrief}
  \item[Why] "차이가 있다"는 두 집단 평균을 따로 두면 점들이 평균에 훨씬
             가까워진다는 뜻이다.
  \item[When] 두 집단의 연속값(회복 시간)을 비교할 때.
  \item[Where] 약 C 대 D.
  \item[How] 합친 평균 하나(귀무)와 집단별 평균 둘(대립)에서 거리를 재어
             비교한다. 줄어든 양이 작으면 기각 실패.
\end{inbrief}
% source: deck 9 slide 9

$H_0$: C와 D 사이에 회복 시간 차이가 없다는 말은, "C와 D를 하나로 합친 평균
하나로도 두 집단을 똑같이 잘 설명한다"는 말과 같다. $H_a$: 차이가 있다는
말은 "C와 D를 각각 따로 둔 평균 둘이 점들에 더 가깝다"는 말이다. 그래서
"차이가 있는가"는 "평균을 둘로 나누면 점들이 평균에서 얼마나 더 가까워지는가"로
바꿔 물을 수 있다.

\begin{example}[합친 평균 대 집단 평균]
\label{ex:ht-distance}
아래 값은 슬라이드 9(\texttt{image22}, \texttt{image23})의 그림 모양을
흉내 낸 \textbf{이 책 자체의 예시 값}이다. 슬라이드는 숫자를 밝히지 않는다.
$C = \{20, 28, 34\}$, $D = \{22, 26, 37\}$.

\begin{align*}
  &\bar{x}_{\text{pooled}} = \frac{20+28+34+22+26+37}{6} = \frac{167}{6} = 27.83
    && \text{합친 평균}\\
  &\bar{x}_C = \frac{20+28+34}{3} = \frac{82}{3} = 27.33 && \text{C의 평균}\\
  &\bar{x}_D = \frac{22+26+37}{3} = \frac{85}{3} = 28.33 && \text{D의 평균}
\end{align*}

\Cref{sec:spread}의 분산처럼 거리는 \textbf{제곱거리 합}으로 잰다.

\begin{align*}
  &(20-27.83)^2+(28-27.83)^2+(34-27.83)^2 \\
  &\quad +(22-27.83)^2+(26-27.83)^2+(37-27.83)^2 = 220.83
    && \text{합친 평균에서의 제곱거리 합}\\
  &(20-27.33)^2+(28-27.33)^2+(34-27.33)^2 = 98.67 && \text{C, 집단 평균에서}\\
  &(22-28.33)^2+(26-28.33)^2+(37-28.33)^2 = 120.67 && \text{D, 집단 평균에서}\\
  &98.67 + 120.67 = 219.33 && \text{집단 평균에서의 제곱거리 합}\\
  &220.83 - 219.33 = 1.50 && \text{줄어든 양}
\end{align*}

검산: $3(27.33-27.83)^2 + 3(28.33-27.83)^2 = 0.75+0.75 = 1.50$, 일치한다.

전체 220.83 중 겨우 1.50, 즉 \textbf{0.7\%}만 줄어든다. \Cref{fig:ht-pooled-vs-group}이
이 두 제곱거리 합을 나란히 보여준다. 평균을 둘로 나누어도 점들이 평균에
거의 더 가까워지지 않으므로 "극적인 차이가 없다"(슬라이드 9). 결론은
기각 실패다. 이 거리를 정식 검정 통계량(Test Statistic)으로 만드는 방법은
이 덱에 없다. \Cref{sec:ht-future}에서 열린 질문으로 남긴다.
\end{example}

\begin{figure}[ht]\centering
  \includegraphics[width=\linewidth]{09_02_pooled_vs_group.png}
  \caption{\Cref{ex:ht-distance}의 값(책 자체의 예시, 슬라이드 데이터
  아님). 왼쪽: 합친 평균 하나에서 잰 제곱거리 합 220.83. 오른쪽: C, D
  각각의 평균에서 잰 제곱거리 합 219.33. 줄어든 양 1.50이 크지 않아
  기각 실패로 이어진다.}
  \label{fig:ht-pooled-vs-group}
\end{figure}

\Cref{tab:ht-cases}가 이 장에서 다룬 네 사례를 한눈에 정리한다.

\begin{table}[ht]\centering
\begin{tabularx}{\linewidth}{@{}l X X l l l@{}}
\toprule
사례 & $H_0$ & $H_a$ & p값 & $\alpha$ & 판정 \\
\midrule
금성 생명체 & 생명체 없음 & 생명체 있음 & (증거 없음) & -- & 기각 실패
  (증거가 나오면 기각) \\
강아지(\Cref{ex:ht-puppy}) & $p=0.5$ & $p>0.5$ & 0.0037 & 1\% & 기각 \\
외계인 침공(\Cref{pit:ht-alien}) & $p=0.10$ & $p>0.10$ & 0.00664 &
  1\% & 기각(방향 문제는 \Cref{pit:ht-alien}) \\
후안의 복권(\Cref{ex:ht-juan}) & $p=0.4$ & $p>0.4$ & 0.064 & 1\% & 기각 실패 \\
\bottomrule
\end{tabularx}
\caption{이 장에서 다룬 네 가설검정 사례의 요약. 금성 사례는 슬라이드 2가
증거 유무 두 갈래를 모두 보여주므로 하나의 확정된 결과로 단정하지 않는다.}
\label{tab:ht-cases}
\end{table}

% =====================================================================
\section{어떻게 적용하는가}
\label{sec:ht-apply}
% source: deck 9 slide 14

\begin{inbrief}
  \item[Why] 지금까지의 절차를 하나의 흐름으로 정리하고, 이 책 안의 다른
             자리와 연결한다.
  \item[When] 새 데이터에서 "차이가 있다"는 문장을 쓰기 전에.
  \item[Where] 과제와 프로젝트의 결론 문장 전부.
  \item[How] 아래 다섯 단계를 순서대로 밟는다.
\end{inbrief}

\begin{figure}[ht]\centering
\begin{tikzpicture}[
  node distance = 7mm,
  flowbox/.style = {draw = bookrule, fill = booksurface, rounded corners = 2pt,
                   minimum width = 62mm, minimum height = 9mm, align = center,
                   font = \small},
  arrow/.style  = {-{Stealth[length=2.2mm]}, draw = bookmuted}]
  \node[flowbox] (h)      {$H_0$ 와 $H_a$ 를 세운다};
  \node[flowbox, below = of h]   (alpha)  {유의수준 $\alpha$ 를 미리 정한다};
  \node[flowbox, below = of alpha] (dir)  {극단의 방향($>$, $<$, $\neq$)을 정한다};
  \node[flowbox, below = of dir]  (pval)  {p값을 구한다\\[-1pt]
        \footnotesize 시뮬레이션(\Cref{sec:ht-simulation}) 또는 정확 계산
        (\Cref{sec:ht-exact})};
  \node[flowbox, below = of pval, fill = bookaccentbg, draw = bookaccent] (dec)
        {p값 $<\alpha$ ?\\[-1pt] \footnotesize 예: 기각 / 아니오: 기각 실패};
  \draw[arrow] (h) -- (alpha);
  \draw[arrow] (alpha) -- (dir);
  \draw[arrow] (dir) -- (pval);
  \draw[arrow] (pval) -- (dec);
\end{tikzpicture}
\caption{가설검정 다섯 단계. 방향(3단계)은 데이터를 보기 전에 정한다
(\Cref{pit:ht-direction}).}
\label{fig:ht-flow}
\end{figure}

이 책 안에서 같은 흐름이 이미 쓰인 자리, 그리고 앞으로 쓸 자리를 짚는다.

\begin{itemize}
  \item \Cref{subsec:ci-compare}: 버거(평균 620, 신뢰구간 $[557, 686]$)와
        치킨너겟(평균 275, 신뢰구간 $[211, 345]$)을 비교했다. 같은 질문을
        $H_0$: "두 평균 칼로리 차이 = 0"으로 다시 쓰면, 신뢰구간이 겹치지
        않는다는 그림이 곧 가설검정 하나를 그림으로 그린 것임을 알 수 있다.
        새 수치는 계산하지 않는다.
  \item \Cref{subsec:bootstrap}: 재표집으로 분포를 만든다는 발상이
        \Cref{sec:ht-simulation}과 같다. 차이는 하나다. 부트스트랩은
        \textbf{표본}에서 다시 뽑고, 가설검정 시뮬레이션은 \textbf{$H_0$
        가 정한 세상}(예: $p=0.5$ 인 동전)에서 뽑는다.
  \item \Cref{ch:hw1}: 과제 데이터에서 세울 수 있는 $H_0$ 의 예, "버거와
        치킨너겟의 평균 칼로리 차이는 0이다" 같은 가설 문장을 계산 없이
        세워 본다.
  \item \Cref{ch:project}: 프로젝트에서 "차이가 있다"고 쓰기 전에
        \Cref{fig:ht-flow}의 다섯 단계를 거쳤는지 점검한다.
\end{itemize}

확인 문제(슬라이드 14의 Kahoot을 대신한다).

\begin{enumerate}
  \item p값이 $\alpha$ 보다 작을 때 올바른 결론 문장을 \Cref{tab:ht-wording}
        에서 고른다.
  \item 외계인이 "10\% 미만이면 침공하지 않는다"는 규칙을 쓴다면 $H_a$ 의
        방향은 $>$ 인가 $<$ 인가.
  \item 성공 확률 0.4인 독립 시행에서 연속 성공 확률이 처음으로 1\% 밑으로
        내려가는 최소 연승 횟수는 몇 번인가.
\end{enumerate}

% =====================================================================
\section{앞으로}
\label{sec:ht-future}
% source: deck 9 slide 14

\begin{itemize}
  \item \term{순열검정}(Permutation Test), \term{통계적 유의성}
        (Statistical Significance), \term{차등 분석}(Differential
        Analysis)은 강의계획서가 예고한 주제다. 이 덱에는 내용이 없으므로
        다음 강의로 이어지는 갈래로만 적는다.
  \item \Cref{ex:ht-distance}의 "평균에서의 거리 비교"를 정식 검정
        통계량으로 만드는 방법은 이 덱에 없다. 열린 질문으로 남긴다.
  \item \Cref{ex:ht-juan}이 드러낸 질문: $H_0$ 가 실제로 틀렸을 때, 표본
        크기가 작으면 기각할 수 없다. "얼마나 큰 표본이어야 진짜 차이를
        잡아낼 수 있는가"(검정력, Power)는 이 덱에 없는 질문이다.
  \item 정규성 가정이 필요한 검정과의 연결은 아직 없다. 분포 확인 방법이
        다뤄질 때 다시 연결한다.
\end{itemize}

\chapter{p값: 우연이라면 이 정도 결과가 얼마나 드문가}
\label{ch:pvalue}
% source: deck 10 slide 1

\begin{inbrief}
  \item[Why] \term{p값}(p-value)은 "우연만으로 이 정도 결과가 나올 확률"을
             숫자 하나로 준다. 그 숫자를 제대로 읽지 못하면 검정 결론을
             거꾸로 쓴다.
  \item[When] 귀무가설(Null Hypothesis)을 세운 뒤, 관측한 데이터가 그
              가설 아래에서 얼마나 드문지 판정할 때.
  \item[Where] 두 약의 완치율 비교, 동전의 공정성, 한 사람의 키가
               모집단에서 튀는지, 평균 비타민 D가 기준값과 다른지.
  \item[How] 관측값과 똑같이 드문 결과, 더 드문 결과의 확률을 모두 더한다
             $\to$ 미리 정한 기준(보통 0.05)과 비교한다 $\to$ p값은 차이의
             \textbf{크기}가 아니라 차이의 \textbf{확실성}만 말한다는 점을
             결론 문장에 반영한다.
\end{inbrief}

\begin{keyidea}
p값도 \analogy{현장 측량으로 지도 검증}이다(\Cref{tab:analogy}). 다른 점은
"어긋났는가 아닌가"만 보는 \Cref{ch:hypothesis}의 기각/기각 실패에서 한 걸음
더 나아가, \textbf{얼마나 어긋났는지를 확률 하나로 재는 자}를 준다는
것이다. 개가 무죄인 세상을 상상하고, 그 세상에서 지금 증거만큼(또는 더)
나쁜 일이 일어날 확률을 재판 예시로 쓴다(슬라이드 3). 이것은 비유가 아니라
p값의 정의를 그대로 보여주는 \textbf{예시}다. \Cref{pit:ht-fail-not-proof}의
"증거 부족은 무죄를 증명하지 않는다"가 이 재판 예시와 같은 문장이다.
\end{keyidea}

이 장의 출처는 \texttt{notes/10\_P-value.pptx}(25 슬라이드) 하나다. 슬라이드
속 예제(약 A/B, 동전, 브라질 여성 키, 비타민 D, 배 칼로리)는 슬라이드 자체를
출처로 적는다.

% =====================================================================
\section{배경}
\label{sec:pv-background}
% source: deck 10 slide 2

이 장은 p값을 \textbf{처음 정의하지 않는다.} 정의는 \Cref{def:pvalue}
하나이고, 이 장은 그 정의를 이산 데이터의 세 부분 분해, 연속 데이터, 기준의
의미, 효과 크기와의 구분, 시뮬레이션과 모수적 계산, 신뢰구간과의 관계로
\textbf{넓힌다.}

\term{귀무가설}(Null Hypothesis) $H_0$ 와 \term{대립가설}(Alternative
Hypothesis) $H_a$ 는 \Cref{def:ht-null,def:ht-alt}에서 이미 정의했다. 약 C와
D의 회복 시간을 비교하는 슬라이드 2의 예(\textbf{덱 9 슬라이드 5, 6의
반복})는 $H_0$: 두 약 사이에 차이가 없다, $H_a$: 차이가 있다는 문장 그대로
\Cref{def:ht-null,def:ht-alt}의 구조다. 새로 덧붙일 내용이 없으므로 여기서
끝낸다.

기각과 기각 실패의 뜻은 \Cref{sec:ht-decision}, 단측/양측 대립가설의 구분은
\Cref{sec:ht-forms,fig:ht-tails}에서 다뤘다. 세 부분 p값(\Cref{sec:pv-discrete})은
이 구분이 왜 필요한지를 다시 보여준다. 곱셈 규칙과 이항계수는
\Cref{sec:ht-background,ex:ht-puppy}에서 다뤘고, 동전 5번의
$\binom{5}{k}/32$(\Cref{sec:pv-discrete})가 그대로 쓴다.

\term{정규분포}(Normal Distribution)는 \Cref{sec:normal,def:normal}이
아직 뼈대만 갖춘 절이라, 여기서 필요한 두 문장만 직접 설명한다. 정규분포는
가운데가 높고 양옆으로 갈수록 낮아지는 종 모양 곡선이고, \textbf{곡선 아래
넓이가 곧 확률}이다. 예를 들어 곡선의 왼쪽 절반 넓이는 $0.5$, 즉
"평균보다 작을 확률 50\%"를 뜻한다. \Cref{sec:pv-continuous}의 키 예제가 이
넓이를 직접 계산한다.

\term{표준오차}(Standard Error)는 \Cref{pit:sd-se-ci}에서 표준편차와
구분됐다. \Cref{sec:pv-sim-vs-param}의 비타민 D 예제에서 3.3이라는 값이
표준편차가 아니라 표준오차라는 점이 판정을 좌우한다. \term{신뢰구간}
(Confidence Interval)은 \Cref{sec:ci,subsec:ci-compare}에서 다뤘고,
\Cref{sec:pv-ci-duality}에서 p값과의 관계를 짚는다.

\Cref{tab:ht-symbols}가 이미 $H_0$, $H_a$, $\alpha$, $n$, $X$ 를 정리했다.
\Cref{tab:pv-symbols}는 이 장에서 새로 쓰는 기호만 더한다.

\begin{table}[ht]\centering
\begin{tabularx}{\linewidth}{@{}l X l@{}}
\toprule
기호 & 뜻 & 읽기 \\
\midrule
$z$ & \term{표준화 점수}(z-score): 관측값이 평균에서 표준오차 몇 개만큼
      떨어져 있는가, $z=(\bar x-\mu_0)/\mathrm{SE}$ & 지 \\
$\mathrm{SE}$ & \term{표준오차}(Standard Error): 표본 평균의 흩어짐. 표준편차와
                다르다(\Cref{pit:sd-se-ci}) & 에스이 \\
$\mu_0$ & 귀무가설이 정한 값(예: 비타민 D 기준 100) & 뮤 제로 \\
\bottomrule
\end{tabularx}
\caption{이 장에서 새로 쓰는 기호. $H_0$, $H_a$, $\alpha$, $n$, $X$ 는
\Cref{tab:ht-symbols}를 그대로 쓴다.}
\label{tab:pv-symbols}
\end{table}

% =====================================================================
\section{왜 필요한가}
\label{sec:pv-why}
% source: deck 10 slide 3, 4

\begin{inbrief}
  \item[Why] "다르다"는 판단에는 근거의 강도를 나타내는 숫자가 필요하다.
             느낌만으로는 사람마다 결론이 갈린다.
  \item[When] 두 집단(또는 값)을 비교했는데 그 차이가 우연인지 아닌지
              애매할 때.
  \item[Where] 신약 A와 B의 완치율 비교(슬라이드 4).
  \item[How] "차이가 없는 세상"을 상상하고, 그 세상에서 지금 본 결과(또는
             더 극단적인 결과)가 나올 확률을 잰다. 그 확률이 p값이다.
\end{inbrief}

이 절의 질문은 이 장 도입부의 개 재판 비유(슬라이드 3)를 약 A/B 비교로
구체화한 것이다. "차이가 없는 세상"을 상상하고 그 세상에서 지금 본
증거가 얼마나 드문지 재는 절차는 똑같다.

한 사람에게 약 A, 한 사람에게 약 B를 주고 A만 나았다고 "A가 낫다"고 말할 수
없다. 우연히 그럴 수 있기 때문이다. 두 명씩 줘도 마찬가지다. 그렇다면
얼마나 많은 사람에게 줘야 확신할 수 있는가.

반대 극단에서는 답이 분명하다. 1,046명 중 1,043명이 완치된 약과 1,434명 중
2명만 완치된 약을 비교하면, 우연이라고 믿기 어렵다.

\begin{align*}
  &1043/1046 = 0.99713 && \text{99.7\%, 거의 전원 완치}\\
  &2/1434 = 0.001395 && \text{0.14\%, 거의 아무도 완치되지 않음}
\end{align*}

\begin{pitfall}[슬라이드 값은 0.001\%, 다시 계산하면 0.14\%]
슬라이드 4는 $2/1434$ 를 \textbf{0.001\%}로 적었다. $2/1434 = 0.0014$ 는
비율이지 퍼센트가 아니다. 이를 퍼센트로 바꾸면 $0.14\%$ 다(비율 자체를
퍼센트로 착각한 것으로 보인다). 99.7\%는 맞다.
\end{pitfall}

그 사이, 198명 중 73명 대 190명 중 59명은 판단이 서지 않는다.

\begin{align*}
  &73/198 = 0.3687 && \text{36.9\%}\\
  &59/190 = 0.3105 && \text{31.1\%}
\end{align*}

\begin{pitfall}[슬라이드 값은 29\%, 다시 계산하면 31.1\%]
\label{pit:pv-drug-2931}
슬라이드 4는 $59/190$ 을 \textbf{29\%}로 적었다. $59/190=0.3105$ 는
\textbf{31.1\%}다(이 장은 이 값을 그대로 쓴다). 36.9\%는 맞다.
\end{pitfall}

\textbf{"얼마나 확신하는가"를 숫자로 줄 도구가 필요하다.} 그것이 p값이고,
슬라이드는 이 마지막 경우의 p값을 $0.24$로 적는다.

\begin{example}[약 A 대 약 B]
\label{ex:pv-drug}
\Cref{tab:pv-drug-trials}가 표본 크기에 따라 같은 질문("A가 B보다 낫다고
말할 수 있는가")의 확신이 어떻게 달라지는지 정리한다.
\end{example}

\begin{table}[ht]\centering
\begin{tabularx}{\linewidth}{@{}X l l l l X@{}}
\toprule
사례 & 약 & 완치 & 비완치 & 완치율(재계산) & 판정(재계산한 p값, p-value) \\
\midrule
1명씩, 2명씩 & 없음 & 없음 & 없음 & 없음 & 표본 너무 작음, 판단 불가 \\
대규모 & A & 1043 & 3 & 99.7\% & 우연으로 보기 어려움(계산 없음) \\
       & B & 2 & 1432 & 0.14\% & \\
약 200명씩 & A & 73 & 125 & 36.9\% & 확신하기 어려움($p\approx0.24$) \\
           & B & 59 & 131 & 31.1\% & \\
약 A 대 약 A(1) & A & 71 & 127 & 35.9\% & 진짜 차이 없음($p=1.0$) \\
                & A' & 72 & 126 & 36.4\% & \\
약 A 대 약 A(2) & A & 73 & 125 & 36.9\% & 진짜 차이 없음($p=0.917$) \\
                & A' & 71 & 127 & 35.9\% & \\
약 A 대 약 A(3) & A & 60 & 138 & 30.3\% & 진짜 차이 없어도 $p=0.012$
                                          로 거짓 양성 가능 \\
                & A' & 84 & 114 & 42.4\% & \\
대표본(효과 크기) & A & 5005 & 9868 & 33.65\% & 차이 작지만 $p=0.044$
                                                (\Cref{sec:pv-effect-size}) \\
                  & B & 4800 & 9000 & 34.78\% & \\
\bottomrule
\end{tabularx}
\caption{약 A 대 약 B(그리고 약 A 대 약 A) 비교, 표본 크기와 진짜 차이
유무에 따라 확신·p값이 어떻게 달라지는지 정리(슬라이드 4, 5, 6, 7).
완치율은 재계산 값이다. 슬라이드는 약 200명씩 사례를 각각 37\%,
29\%(오류, \Cref{pit:pv-drug-2931} 함정 참고)로 적는다.}
\label{tab:pv-drug-trials}
\end{table}

% =====================================================================
\section{기준 0.05와 거짓 양성}
\label{sec:pv-threshold}
% source: deck 10 slide 5, 6

\begin{inbrief}
  \item[Why] "얼마나 작아야 작은가"에 답이 없으면 판정이 사람마다 달라진다.
  \item[When] 데이터를 보기 전에 기준을 정할 때.
  \item[Where] 약 A 대 약 A처럼 실제 차이가 없는 비교.
  \item[How] 기준 0.05는 "차이가 없을 때 실험 100번 중 약 5번은 잘못
             기각한다"는 약속이다. 이 잘못을 \term{거짓 양성}(False
             Positive), \term{1종 오류}(Type I Error)라 부른다.
\end{inbrief}

슬라이드 5, 6은 약 A를 \textbf{약 A 자신}과 비교한다. 진짜 차이는 없다.
그런데도 두 쌍을 뽑으면 p값이 달라진다.

여기서 쓰는 $[[a,b],[c,d]]$ 표기는 $2\times2$ 표다. 첫 행이 첫 번째 표본의
완치·비완치 인원, 둘째 행이 두 번째 표본의 완치·비완치 인원이다(예:
$[[71,127],[72,126]]$은 첫 번째 표본에서 완치 71명·비완치 127명, 두 번째
표본에서 완치 72명·비완치 126명). \term{카이제곱 검정}(Chi-squared Test)은
"두 표본의 완치율이 같다"는 귀무가설 아래 이 표가 얼마나 드문지를 재는
검정이고, 이 절의 p값은 모두 이 검정에서 나온다. \term{Yates 보정}(Yates'
Continuity Correction)은 표본이 작을 때 카이제곱 값을 살짝 줄여 더 보수적인
p값을 주는 보정이며, \term{피셔 정확검정}(Fisher's Exact Test)은 근사 없이
정확한 확률을 계산하는 대안이다. 이 검정들의 계산식은 이 장에서 다루지
않고, 슬라이드와 재계산이 주는 결과값만 그대로 옮긴다.

\begin{align*}
  &[[71,127],[72,126]] && \text{첫째 쌍}\\
  &\chi^2_{\text{Yates}} = 0 && \text{Yates 보정 카이제곱}\\
  &p = 1.0 && \text{완전히 같은 비율}
\end{align*}
\begin{align*}
  &[[73,125],[71,127]] && \text{둘째 쌍}\\
  &\chi^2_{\text{Yates}} = 0.011 && \text{Yates 보정 카이제곱}\\
  &p = 0.917 && \text{역시 큰 p값}
\end{align*}

두 쌍 모두 큰 p값을 주지만, 표본을 다르게 뽑으면 작은 p값도 나올 수 있다.

\begin{align*}
  &[[60,138],[84,114]] && \text{역시 약 A 대 약 A}\\
  &60/198 = 0.303 && \text{완치율(첫 행)}\\
  &84/198 = 0.424 && \text{완치율(둘째 행)}\\
  &\chi^2 = 6.29 && \text{보정 없음}\\
  &p = 0.012 && \text{보정 없음 p값}\\
  &p_{\text{Yates}} = 0.016 && \text{Yates 보정 p값}
\end{align*}

\textbf{진짜 차이가 없는데도} p값이 작게 나올 수 있다. 이것이 슬라이드 6의
문장이다. 기준 0.05를 쓰면, 차이가 없는 실험 100번 중 약 5번은 이렇게
잘못 기각한다. 이 잘못을 \term{거짓 양성}(False Positive), \term{1종 오류}
(Type I Error)라 부른다. \Cref{fig:pv-false-positive}가 이것을 1,000번
모의로 직접 보여준다.

\begin{figure}[ht]\centering
  \includegraphics[width=0.75\linewidth]{10_05_false_positive.png}
  \caption{완치율이 똑같은(36.9\%) 두 집단을 1,000번 비교한 p값의 분포.
  카이제곱 검정(보정 없음)은 5\%에 가장 가깝게(약 4.6\%) 0.05 아래로
  떨어진다. Yates 보정이나 피셔 정확검정을 쓰면 약 3.3\%로 더 보수적이게
  나와, 슬라이드 6의 "5\%"를 눈으로 확인하기 어려워진다.}
  \label{fig:pv-false-positive}
\end{figure}

기준 0.05가 법칙이 아니라 약속이라는 점은 \Cref{pit:ht-alpha}에서 이미
다뤘다. 슬라이드 6도 "The threshold is arbitrary"라고 직접 적어 같은
점을 확인해 준다.

\begin{pitfall}[p값 0.24는 "24\% 확률로 차이가 없다"는 뜻이 아니다]
p값은 \textbf{$H_0$ 가 참이라는 가정 아래} 지금 증거(또는 더 극단적인
증거)가 나올 확률이다(\Cref{def:pvalue}). "차이가 없을 확률이 24\%"는
$H_0$ 자체의 확률을 말하는 문장인데, p값은 그런 확률을 계산하지 않는다.
p값 $0.24$는 "차이가 없다는 가정 아래에서도 이 정도 결과는 꽤 흔하다"는
뜻일 뿐이다.
\end{pitfall}

% =====================================================================
\section{p값은 효과 크기가 아니다}
\label{sec:pv-effect-size}
% source: deck 10 slide 7

\begin{inbrief}
  \item[Why] 작은 p값을 "큰 차이"로 읽는 실수가 가장 흔하다.
  \item[When] 결론 문장을 쓸 때.
  \item[Where] 약 A 34\% 대 약 B 35\%, 표본 약 29,000명.
  \item[How] p값과 함께 차이 자체(\term{효과 크기}, Effect Size)를 따로
             적는다.
\end{inbrief}

$[[5005,9868],[4800,9000]]$(슬라이드 7, \texttt{image15})은 완치율이
겨우 다르다.

\begin{align*}
  &5005/14873 = 0.3365 && \text{33.65\%}\\
  &4800/13800 = 0.3478 && \text{34.78\%}\\
  &34.78 - 33.65 = 1.1\ \%\text{포인트} && \text{차이}\\
  &\chi^2 = 4.07 && \text{보정 없음}\\
  &z = 2.02 && \text{표준화 점수}\\
  &p = 0.044 && \text{양측 p값}
\end{align*}

\begin{pitfall}[p값이 작다고 차이가 크다는 뜻은 아니다]
표본이 \textbf{크면}(여기서는 약 14,000명씩) 아주 작은 차이(1.1\%포인트)도
p값을 0.05 아래로 밀어낸다. p값 0.044는 "차이가 우연이 아닐 가능성이
크다"는 뜻이지 "차이가 크다"는 뜻이 아니다. 결론 문장에는 항상 p값과
\term{효과 크기}(Effect Size, 여기서는 1.1\%포인트)를 \textbf{함께} 적는다.
\end{pitfall}

% =====================================================================
\section{이산 데이터의 p값: 세 부분}
\label{sec:pv-discrete}
% source: deck 10 slide 8-12

\begin{inbrief}
  \item[Why] "관측값이 나올 확률"만 보면 모든 결과가 드물어 보인다.
  \item[When] 결과를 셀 수 있을 때(동전).
  \item[Where] 동전 2번, 5번.
  \item[How] 관측값의 확률 $+$ 똑같이 드문 다른 결과의 확률 $+$ 더 드문
             결과의 확률.
\end{inbrief}

동전을 2번 던져 2번 다 앞면(HH)이 나왔다. "내 동전은 특별한가"라는 질문에
"HH가 나올 확률은 $0.25$뿐이니 특별하다"고 답하면 안 된다. HH, HT, TH, TT
\textbf{네 결과 모두} 확률이 $0.25$ 이하다. 모든 결과가 드물다면 "드묾"만으로는
아무것도 가려낼 수 없다. \Cref{fig:pv-coin}가 동전 2번과 5번(뒤에 나올
\Cref{ex:pv-coin5})의 결과 분포를 막대로 보여준다.

\begin{align*}
  &P(\text{HH}) = 0.5\times0.5 = 0.25 && \text{관측값}\\
  &P(\text{TT}) = 0.25 && \text{똑같이 드묾(2/2 반대 방향)}\\
  &P(\text{HT})+P(\text{TH}) = 0.25+0.25 = 0.5 && \text{HH보다 흔함, 더 드물지 않음}
\end{align*}

p값은 관측값의 확률에 \textbf{똑같이 드문 결과}, \textbf{더 드문 결과}의
확률까지 더한다. \Cref{fig:pv-three-parts}가 동전 2번의 수형도와 세
부분을 그림으로 정리한다.

\begin{align*}
  &0.25\ (\text{관측 HH}) + 0.25\ (\text{똑같이 드묾, TT}) + 0\ (\text{더 드문 결과 없음}) \\
  &= 0.5
\end{align*}

$0.5$ 는 $0.05$ 보다 훨씬 크므로 기각 실패다. "HH가 나온 건 특별해 보이지만,
TT도 똑같이 드물다는 걸 알면 그렇게 특별하지 않다."

꽃에 비유하면 이해가 쉽다(슬라이드 11). 지금 보는 꽃이 "가장 희귀한 꽃"
같아 보여도, 그것과 똑같이 희귀한 꽃이 여럿 있고 그보다 \textbf{더} 희귀한
꽃도 많다는 걸 알면 더 이상 특별해 보이지 않는다. "더 드문 결과가 많을수록
지금 본 결과는 덜 특별하다"는 것이 세 번째 부분(더 드묾)이 하는 일이다.

\begin{example}[동전 5번 중 앞면 4번]
\label{ex:pv-coin5}
동전을 5번 던져 앞면이 4번 나왔다(슬라이드 12). $2^5=32$가지 결과가 있고,
앞면 수 $k$ 별 경우의 수는 $\binom{5}{k}$ 로 $1,5,10,10,5,1$ 이다.

\begin{align*}
  &\binom{5}{4} = 5 && \text{관측값(앞4 뒤1)의 경우의 수}\\
  &\binom{5}{1} = 5 && \text{똑같이 드묾(앞1 뒤4)의 경우의 수}\\
  &\binom{5}{5}+\binom{5}{0} = 1+1 = 2 && \text{더 드묾(앞5 또는 앞0)의 경우의 수}\\
  &(5+5+2)/32 = 12/32 = 0.375 && \text{양측 p값}
\end{align*}

$0.375 > 0.05$ 이므로 기각 실패다. "내 동전이 특별하다는 증거가 없다."
\end{example}

\begin{pitfall}[세 부분 = 양측 p값이다]
\Cref{ex:ht-puppy}(강아지)와 \Cref{ex:ht-juan}(복권)은 단측
($H_a$: 수컷이 더 많다, 당첨이 더 잦다)이라 "둘째 부분"(똑같이 드문
반대쪽)이 애초에 없다. 세 부분 분해는 정의 하나(\Cref{def:pvalue})를
\textbf{양측} 대립가설에 적용한 모양일 뿐이다. 같은 관측(앞면 4번)을
단측($H_a$: 앞면이 더 잘 나온다)으로 물으면 둘째 부분이 빠진다.

\begin{align*}
  &(5+1)/32 = 6/32 = 0.1875 && \text{단측: 관측값(앞4) + 더 드묾(앞5)만}
\end{align*}

0번(앞면 0번)은 "앞면이 \textbf{더} 잘 나온다"는 방향의 반대쪽이므로 이
단측 p값에는 들어가지 않는다.

\Cref{tab:pv-three-parts-compare}가 양측과 단측을 나란히 둔다.
\end{pitfall}

\begin{table}[ht]\centering
\begin{tabularx}{\linewidth}{@{}X l l@{}}
\toprule
$H_a$ & 포함하는 부분 & p값 \\
\midrule
$\neq$(양측) & 관측 + 똑같이 드묾 + 더 드묾 & $12/32=0.375$ \\
$>$(단측, 앞면이 더 잘 나온다) & 관측 + 더 드묾만 & $6/32=0.1875$ \\
\bottomrule
\end{tabularx}
\caption{동전 5번 중 앞면 4번, 양측과 단측 p값 비교.}
\label{tab:pv-three-parts-compare}
\end{table}

\begin{figure}[ht]\centering
  \includegraphics[width=\linewidth]{10_01_coin_pvalue.png}
  \caption{동전 2번(왼쪽)과 5번(오른쪽)의 결과 분포. 오른쪽 막대 색이
  관측(주황, 앞면 4번), 똑같이 드묾(파랑, 앞면 1번), 더 드묾(빨강, 앞면
  0 또는 5번)을 구분한다. 제목의 $12/32=0.375$가 세 색을 모두 더한 값이다.}
  \label{fig:pv-coin}
\end{figure}

\begin{figure}[ht]\centering
\begin{tikzpicture}[
  level distance = 11mm,
  edge from parent/.style = {draw = bookmuted},
  every node/.style = {font=\small},
  box/.style = {draw = bookrule, rounded corners = 2pt, minimum width = 30mm,
                minimum height = 8mm, align = center}]
  \node {시작}
    child { node {앞} edge from parent node[left] {0.5}
      child { node {HH, $0.25$} edge from parent node[left] {0.5} }
      child { node {HT, $0.25$} edge from parent node[right] {0.5} } }
    child { node {뒤} edge from parent node[right] {0.5}
      child { node {TH, $0.25$} edge from parent node[left] {0.5} }
      child { node {TT, $0.25$} edge from parent node[right] {0.5} } };
\end{tikzpicture}
\hspace{4mm}
\begin{tikzpicture}[node distance = 3mm]
  \node[draw=bookaccent, fill=bookaccentbg, rounded corners=2pt,
        minimum width=30mm, minimum height=8mm, align=center, font=\small]
        (obs) {관측\\$0.25$};
  \node[draw=bookmuted, fill=booksurface, rounded corners=2pt,
        minimum width=30mm, minimum height=8mm, align=center, font=\small,
        below = of obs] (eq) {똑같이 드묾\\$0.25$};
  \node[draw=bookwarn, fill=bookwarnbg, rounded corners=2pt,
        minimum width=30mm, minimum height=8mm, align=center, font=\small,
        below = of eq] (more) {더 드묾\\$0$};
\end{tikzpicture}
\caption{동전 2번의 수형도(왼쪽, 슬라이드 9 구조)와 관측이 HH일 때 세 부분
(오른쪽, 값 없는 흐름도). 오른쪽 상자를 모두 더하면 $0.5$다.}
\label{fig:pv-three-parts}
\end{figure}

% =====================================================================
\section{연속 데이터의 p값: 넓이}
\label{sec:pv-continuous}
% source: deck 10 slide 13, 14

\begin{inbrief}
  \item[Why] 연속값에서는 한 점의 확률이 0이라 "관측값의 확률"을 셀 수
             없다.
  \item[When] 키, 농도 같은 측정값.
  \item[Where] 브라질 여성 키.
  \item[How] 관측값보다 극단적인 쪽의 곡선 아래 넓이를 더한다(양측이면
             양쪽).
\end{inbrief}

동전은 결과를 셀 수 있지만, 키는 셀 수 없다. 키가 정확히 $142.000\ldots$
cm일 확률은 $0$이다. 대신 "142 cm \textbf{보다} 작은 쪽"처럼 범위의 넓이를
잰다.

\begin{example}[브라질 여성 키]
\label{ex:pv-height}
1996년 브라질 15--49세 여성 키의 95\%가 142--169 cm 사이였다(슬라이드
13, \Cref{fig:pv-height}). 이 구간의 중점과 폭에서 평균과 표준편차를
되짚는다.

\begin{align*}
  &(142+169)/2 = 155.5 && \text{구간 중점, 평균 155.7과 근사(반올림 차)}\\
  &169-142 = 27 && \text{구간 폭}\\
  &2\times1.96 = 3.92 && \text{표준정규 95\% 구간 폭의 배수}\\
  &27/3.92 = 6.89 && \text{95\% 구간에서 되짚은 표준편차}
\end{align*}

$\sigma=6.89$ cm는 \textbf{슬라이드에 없는 이 책의 유도값}이다. 키가
$N(155.7,\,6.89^2)$을 따른다고 두고 142 cm의 양측 p값을 계산한다.

\begin{align*}
  &P(h<142) = 0.0234 && \text{왼쪽 꼬리}\\
  &P(h>169) = 0.0268 && \text{오른쪽 꼬리(중점이 155.7이라 완전히 대칭은 아님)}\\
  &0.0234+0.0268 = 0.0502 && \text{양측 p값}
\end{align*}

$0.0502$ 는 $0.05$ 보다 \textbf{약간 크다.}
\end{example}

\begin{pitfall}[슬라이드 값은 "귀무가설: 이 사람이 특별하다", 바르게는 반대다]
\label{pit:pv-null-reversed}
슬라이드 13은 "Null Hypothesis: This person is special from this
distribution"이라고 적는다. 귀무가설은 항상 "차이가 없다"는 쪽이어야
한다(\Cref{def:ht-null}). 올바른 문장은 "$H_0$: 이 사람은 이 분포에서
나왔다(특별하지 않다)"다. 기각하면 그때 "특별하다"고 말한다.
\end{pitfall}

\begin{pitfall}[슬라이드는 p = 0.05를 유의하다고 적지만, 바르게는 경계에서
기각 실패다]
슬라이드 13은 $0.025+0.025=0.05$ 를 "statistically significant"라고
적는다. \Cref{def:pvalue}와 \Cref{sec:ht-pvalue}의 규칙은 "p값이 $\alpha$
\textbf{보다 작으면} 기각"이다. $0.05$는 $0.05$보다 작지 않으므로
엄밀히는 \textbf{기각 실패}다. 142 cm는 경계선 바로 위에 있는 값이다. 이
장의 재계산 $0.0502$도 마찬가지로 경계를 살짝 넘는다.
\end{pitfall}

두 번째 예에서는 평균 근처를 묻는다. 155.4--156 cm 사이에서 사람을 볼
확률의 p값은 얼마인가.

\begin{align*}
  &P(155.4<h<156) = 0.035 && \text{가운데 띠}\\
  &P(h<155.4) = 0.483 && \text{왼쪽 나머지}\\
  &P(h>156) = 0.483 && \text{오른쪽 나머지}\\
  &0.035+0.483+0.483 = 1.001\approx1 && \text{합계(반올림 오차)}
\end{align*}

\begin{keyidea}
평균을 포함하는 구간의 양측 p값은 \textbf{언제나 1에 가깝다.} 평균 근처는
전혀 드물지 않기 때문이다. p값이 크다고 "특별하다"는 뜻이 되지 않는 것과
정반대로, p값이 작아야 비로소 "드물다"는 신호다.
\end{keyidea}

\begin{figure}[ht]\centering
  \includegraphics[width=\linewidth]{10_02_height_pvalue.png}
  \caption{브라질 여성 키 $N(155.7,\,6.89^2)$($\sigma$는 이 책의 유도값,
  \Cref{ex:pv-height}). 왼쪽: 142/169 cm 꼬리(각 0.0234, 0.0268), 양측
  p값 0.0502(슬라이드는 각 0.025, 합 0.05로 적는다). 오른쪽:
  [155.4, 156] 띠 0.035와 양쪽 나머지 각 0.483.}
  \label{fig:pv-height}
\end{figure}

% =====================================================================
\section{시뮬레이션과 모수적 계산}
\label{sec:pv-sim-vs-param}
% source: deck 10 slide 16, 17, 18

\begin{inbrief}
  \item[Why] 귀무가설 아래 분포를 알아야 넓이를 잰다. 그 분포를 만드는
             길이 두 가지다.
  \item[When] 표준오차를 알 때.
  \item[Where] 비타민 D 평균 대 기준 100.
  \item[How] (a) 귀무가설 세상에서 연구 10,000번을 모의해 관측값 이상을
             센다, (b) 정규분포를 알면 $z=(\bar x-\mu_0)/\mathrm{SE}$ 로
             바로 계산한다.
\end{inbrief}

\begin{example}[비타민 D 평균이 100보다 낮은가]
\label{ex:pv-vitd}
연구 질문(슬라이드 16): 중년/노년 유럽 남성의 평균 비타민 D가 "적정"
수준인 100 nmol/L보다 낮은가. $H_0$: 평균 $=100$. 관측한 표본 평균은
63(슬라이드 17, 18), 표준오차 $\mathrm{SE}=3.3$.

\textbf{(a) 시뮬레이션.} $H_0$ 세상(평균 100, SE 3.3인 정규분포)에서 연구를
10,000번 모의하면, 관측값 63 아래로 내려가는 경우가 \textbf{0번}이다
(\Cref{fig:pv-vitd} 왼쪽).

\begin{align*}
  &0/10{,}000 && \text{시뮬레이션에서 63 이하가 나온 횟수}\\
  &p < 1/10{,}000 && \text{p값은 이 비율보다 작다고만 말할 수 있다}
\end{align*}

\textbf{(b) 모수적 계산.} 정규분포임을 알면 $z$ 점수로 바로 계산한다.

\begin{align*}
  &63-100 = -37 && \text{분자(관측 평균 $-$ 귀무값)}\\
  &-37/3.3 = -11.21 && \text{표준화 점수 } z\\
  &p = 2\times P(Z<-11.21) \approx 3.6\times10^{-29} && \text{양측 p값, 사실상 0}
\end{align*}

두 방법 모두 $H_0$ 를 기각한다. 표본 평균의 분포가 정규분포에 가깝기
때문에, 10,000번 모의(a)와 정규분포 공식(b)이 같은 결론에 이른다. 이는
\Cref{sec:ht-simulation,sec:ht-exact}가 강아지 예제에서 보인 "시뮬레이션
대 정확 계산" 짝과 같은 구조다.
\end{example}

\begin{pitfall}[슬라이드는 관측 평균으로 62와 63을 섞어 쓴다, 이 책은 63으로
통일한다]
\label{pit:pv-vitd-62-63}
슬라이드 15(\texttt{image36})는 신뢰구간 계산에서 관측 평균 \textbf{62}를
쓰지만, 슬라이드 17, 18(\texttt{image39}, \texttt{image41})의 시뮬레이션과
$z$ 계산은 \textbf{63}을 쓴다. 이 장은 본문·그림 전부 \textbf{63}으로
통일한다. 62를 써도 결론은 바뀌지 않는다($z=(62-100)/3.3=-11.52$, 여전히 사실상 0).
\end{pitfall}

\begin{pitfall}[슬라이드는 부호 없이 $z=11.2$, 다시 계산하면 $z=-11.21$이다]
슬라이드 18은 "$z=(63-100)/3.3=11.2$"라고 부호 없이 적는다. $63-100=-37$
이므로 $z=-11.21$ 이 정확한 값이다(크기만 보면 11.2가 맞다). 같은 슬라이드의
"true average weight"라는 표현은 "비타민 D 평균"의 오타로 보인다.
\end{pitfall}

\begin{pitfall}[슬라이드는 연구 질문을 단측으로, 대립가설을 양측으로
적었다. 바르게는 둘 다 단측($H_a: \mu<100$)이다]
슬라이드 16의 연구 질문은 "100보다 \textbf{낮은가}"($H_a: \mu<100$, 단측)인데,
같은 슬라이드가 적은 대립가설은 "평균이 100이 \textbf{아니다}"($H_a:\mu\neq
100$, 양측)다. 바르게는 연구 질문과 대립가설의 방향을 하나로 맞춰야 한다.
\Cref{pit:ht-direction}이 경고한 것과 같은 종류의 불일치다.
어느 쪽으로 계산해도 이 예에서는 기각으로 결론이 같다.
\end{pitfall}

\begin{figure}[ht]\centering
  \includegraphics[width=\linewidth]{10_03_vitd_null.png}
  \caption{비타민 D: $H_0$ 아래 표본 평균의 시뮬레이션 10,000개(왼쪽, 평균
  100, SE 3.3)와 관측값 63(점선). 10,000개 중 63 이하는 0개다. 오른쪽:
  관측 평균 63을 중심으로 한 표본 평균의 분포와 95\% 구간
  $[56.4,\,69.6]$(\Cref{ex:pv-vitd}, \Cref{pit:pv-vitd-62-63}).}
  \label{fig:pv-vitd}
\end{figure}

% =====================================================================
\section{신뢰구간과 가설검정}
\label{sec:pv-ci-duality}
% source: deck 10 slide 15, 19, 20

\begin{inbrief}
  \item[Why] 같은 데이터에서 두 도구가 서로 다른 말을 하면 안 된다.
  \item[When] 신뢰구간이 이미 있을 때.
  \item[Where] 비타민 D 95\%, 99\% 구간, 배 두 품종.
  \item[How] \term{유의수준}(Significance Level) $=1-$ \term{신뢰수준}
             (Confidence Level). 귀무가설의 값이 구간 밖이면 기각,
             안이면 기각 실패. 구간은 p값보다 정보가 많다(크기와
             불확실성을 함께 준다).
\end{inbrief}

비타민 D 예제(관측 평균 63, SE 3.3)로 이어간다. 신뢰구간은
$\text{관측값}\pm z\cdot\mathrm{SE}$ 형태로 만든다(슬라이드 15).

\begin{align*}
  &63 - 1.96\times3.3 = 56.53 && \text{95\% 구간 하한(정확한 1.96)}\\
  &63 + 1.96\times3.3 = 69.47 && \text{95\% 구간 상한}\\
  &63 \pm 2\times3.3 = [56.4,\ 69.6] && \text{슬라이드가 쓰는 반올림 2}
\end{align*}

\begin{pitfall}[슬라이드 값은 62 기준 95\% 구간 (55.4, 68.6), 다시
계산하면 63 기준 (56.4, 69.6)이다]
\label{pit:pv-vitd-ci}
슬라이드 15는 관측 평균 62로 $62\pm2\times3.3=[55.4,\,68.6]$을 계산한다.
\Cref{pit:pv-vitd-62-63}에서 밝힌 대로 이 장은 63으로 통일하므로, 95\%
구간은 $[56.4,\,69.6]$, 99\% 구간은 $63\pm2.576\times3.3=[54.5,\,71.5]$다
(슬라이드의 62 기준 99\% 구간은 $[53.5,\,70.5]$). 어느 쪽을 쓰든 귀무값
100은 구간 \textbf{밖}이라 결론(기각)은 같다.
\end{pitfall}

두 구간 모두 귀무값 $100$을 포함하지 않는다(\Cref{fig:pv-ci-test} 위).
\textbf{같은 데이터에서 신뢰구간과 가설검정이 같은 결론에 이른다.}
\Cref{tab:pv-ci-vs-p}가 이 규칙을 정리한다.

\begin{table}[ht]\centering
\begin{tabularx}{\linewidth}{@{}l l X X@{}}
\toprule
신뢰수준(Confidence Level) & 유의수준(Significance Level, $\alpha$) & 귀무값이 구간 밖이면 & 귀무값이 구간 안이면 \\
\midrule
95\% & 0.05 & $p<0.05$, 기각 & $p\ge0.05$, 기각 실패 \\
99\% & 0.01 & $p<0.01$, 기각 & $p\ge0.01$, 기각 실패 \\
\bottomrule
\end{tabularx}
\caption{유의수준 $=1-$ 신뢰수준(슬라이드 19). 비타민 D의 두 구간 모두
귀무값 100을 포함하지 않으므로 기각이다.}
\label{tab:pv-ci-vs-p}
\end{table}

슬라이드 19는 "신뢰구간에서 p값을 항상 이끌어낼 수 있다"고 적는다. 이
문장이 성립하려면 구간이 \textbf{관측값 $\pm z\cdot\mathrm{SE}$} 형태로
만들어져야 한다. 이 장의 모든 구간(비타민 D, 배)이 그 형태다. 구간이 다른
방식(예: 비대칭 구간, 부트스트랩 구간)으로 만들어졌다면 이 변환이 그대로
성립하지 않는다.

\begin{example}[배 두 품종의 평균 칼로리 차이]
\label{ex:pv-pears}
Bosc 배 65개의 평균 120 kcal, 표준편차 15. Anjou 배 65개의 평균 116
kcal, 표준편차 13(슬라이드 20). $H_0$: 두 품종의 평균 칼로리 차이는
$0$이다.

\begin{align*}
  &15/\sqrt{65} = 1.860 && \text{Bosc의 표준오차}\\
  &13/\sqrt{65} = 1.612 && \text{Anjou의 표준오차}\\
  &\mathrm{SE}_{\text{diff}} = \sqrt{1.860^2+1.612^2} = 2.462 && \text{차이의 표준오차}\\
  &120-116 = 4 && \text{관측된 평균 차이}\\
  &6.44/2.462 = 2.616 && \text{자유도 약 125인 t분포의 99\% 임계값과 일치}\\
  &4 \pm 6.44 = [-2.44,\ 10.44] && \text{99\% 구간}
\end{align*}

구간 $[-2.44,\,10.44]$가 귀무값 $0$을 \textbf{포함한다.} 따라서 $p>0.01$,
\textbf{기각 실패}다. 올바른 결론 문장은 "두 품종의 평균 칼로리가
다르다고 말할 근거가 없다"다.
\end{example}

\begin{pitfall}[슬라이드는 가정문과 실제 결론을 한 화면에 섞어 적는다,
바르게는 처음부터 끝까지 기각 실패다]
슬라이드 20의 텍스트를 그대로 옮기면 다음과 같다.
\begin{enumerate}
  \item "이 값을 100번 계산하면 99번이 이 구간 안에 있다"는 해석이 틀렸다.
        올바른 해석은 \textbf{같은 절차로 구간을 100번 만들면 약 99개가
        참값을 포함한다}는 뜻이다(구간 하나를 100번 다시 \textbf{계산}하는
        것이 아니라, 표본을 100번 다시 \textbf{뽑아} 구간을 100번 만든다).
  \item "P-value $<$ 0.01 $\to$ reject $\to$ the pears \textbf{are}
        different" 줄은 \textbf{"구간이 겹치지 않았다면"이라는 가정 아래의
        문장}이다("If it did not overlap, we could say this is that 1 in
        100 times"). 실제로는 구간이 $0$을 포함하므로 이 가정은 성립하지
        않는다. 슬라이드의 실제 결론은 그 다음 줄, "$0$ is in the
        interval $\to$ nothing surprising $\to$ P-value $>$ 0.01 $\to$ fail
        to reject"다. 두 갈래를 나란히 적어 놓아 가정문을 실제 결론으로
        잘못 읽기 쉽다.
  \item 실제 결론 줄의 "Apples are the same in their average caloric
        content"의 "Apples"는 "\textbf{Pears}"의 오타로 보인다(이 예는
        배(Bosc, Anjou) 이야기이고 사과는 등장하지 않는다).
  \item "the pears are different", "Apples[Pears] are the same"은 그
        자체로도 \Cref{tab:ht-wording}의 틀린 표현이다. 기각했을 때도
        "다르다고 증명했다"가 아니라 "다르다고 말할 근거가 있다"이고,
        기각 실패일 때도 "같다"가 아니라 "다르다고 말할 근거가
        \textbf{없다}"다.
\end{enumerate}
이 장의 실제 계산(위)과 슬라이드의 실제 결론은 \textbf{같다.} 둘 다 구간이
$0$을 포함하므로 기각 실패다. 슬라이드에는 "반대 결론"이 없다. 가정문
줄만 따로 떼어 읽으면 기각처럼 보일 뿐이다.
\end{pitfall}

\begin{table}[ht]\centering
\begin{tabularx}{\linewidth}{@{}l X X@{}}
\toprule
& 판정: 기각 & 판정: 기각 실패 \\
\midrule
진실: 차이 없음 & \term{거짓 양성}(False Positive, 1종 오류/Type I Error,
  \Cref{sec:pv-threshold}) & 올바른 판정 \\
진실: 차이 있음 & 올바른 판정 & (\term{2종 오류}(Type II Error),
  \Cref{sec:pv-future}) \\
\bottomrule
\end{tabularx}
\caption{진실과 판정의 $2\times2$ 표. 이 장은 거짓 양성(1종 오류) 칸만
채운다. 나머지 오류 칸은 \Cref{sec:pv-future}에서 열린 질문으로 남긴다.}
\label{tab:pv-error-types}
\end{table}

\begin{figure}[ht]\centering
  \includegraphics[width=0.8\linewidth]{10_04_ci_vs_test.png}
  \caption{위: 비타민 D 95\% 구간 $[56.4,\,69.6]$과 99\% 구간
  $[54.5,\,71.5]$(둘 다 63 기준, \Cref{pit:pv-vitd-ci}) 모두 귀무값
  100을 벗어난다. 아래: 배 두 품종의 칼로리 차이 99\% 구간
  $[-2.44,\,10.44]$는 귀무값 0을 포함해 기각 실패다.}
  \label{fig:pv-ci-test}
\end{figure}

% =====================================================================
\section{어떻게 적용하는가}
\label{sec:pv-apply}
% source: deck 10 slide 21

\begin{inbrief}
  \item[Why] p값 하나만으로는 결론을 온전히 전달하지 못한다.
  \item[When] 검정 결과를 문장으로 옮길 때.
  \item[Where] 이 책 안의 모든 결론 문장.
  \item[How] "[p값]이므로 [기각/기각 실패]다. 차이는 [효과 크기, 신뢰구간]
             이다"라는 틀을 쓴다.
\end{inbrief}

\Cref{fig:ht-flow}의 다섯 단계 마지막에 이 문장 틀을 붙인다. p값만 적고
끝내지 않는다.

\begin{itemize}
  \item \Cref{subsec:ci-compare}, \Cref{fig:ci-compare}: 버거(평균 620,
        95\% 구간 $[557,686]$)와 치킨너겟(평균 275, 구간 $[211,345]$)의
        캡션이 쓴 $p<0.05$가 바로 이 장의 신뢰구간-검정 쌍대성
        (\Cref{sec:pv-ci-duality})의 한 사례다. 새 수치는 계산하지
        않는다.
  \item \Cref{pit:sd-se-ci}: 표준오차와 표준편차의 구분이
        \Cref{ex:pv-vitd}에서 $z=-11.21$이라는 큰 값을 만드는 이유다.
  \item \Cref{ch:hw1}: 과제 데이터로 결론 문장 틀을 채워 보는 연습(계산
        없이 문장만).
  \item \Cref{ch:project}: 프로젝트 결과 절에서 p값과 효과 크기를 함께
        적었는지 점검한다.
\end{itemize}

확인 문제(슬라이드 21 Kahoot을 대신한다).

\begin{enumerate}
  \item 동전을 3번 던져 3번 모두 앞면이 나왔다. 양측 p값은
        $2/8=0.25$다. 이 값이 $\alpha=0.05$ 보다 큰가 작은가, 결론은
        기각인가 기각 실패인가.
  \item p값이 $0.04$이고 두 집단의 차이가 1\%포인트일 때, "차이가 크다"고
        말할 수 있는가.
  \item 95\% 신뢰구간이 $0$을 포함할 때, $\alpha=0.05$ 기준의 결론은
        무엇인가.
\end{enumerate}

% =====================================================================
\section{앞으로}
\label{sec:pv-future}
% source: deck 10 slide 22-25

이 덱의 마지막 네 슬라이드(22--25)는 다음 강의를 미리 보여준다. 두 집단의
평균을 정식으로 비교하는 \term{t검정}(t-test), 분포 모양 전체를 비교하는
\term{콜모고로프-스미르노프 검정}(Kolmogorov-Smirnov Test)이 그것이다.
슬라이드 24, 25의 두 표본 t검정 공식과 토마토 예제는 이 장에서 계산하지
않는다. 분포 비교를 다루는 다음 장이 이어받는다.

\begin{itemize}
  \item \Cref{ex:ht-distance}가 열어 둔 질문, "평균에서의 거리 비교를
        정식 검정 통계량으로 만드는 방법"이 다음 장에서 이어진다.
  \item \term{2종 오류}(Type II Error)와 \term{검정력}(Power): 이 장의
        \Cref{tab:pv-error-types}가 채우지 못한 칸이다. \Cref{sec:ht-future}
        의 검정력 질문과 같은 갈래다. 이 덱에는 계산 방법이 없다.
  \item 검정을 여러 번 반복하면(유전자 수천 개를 한 번에 검정하는 경우처럼)
        기준 $0.05$가 만드는 거짓 양성이 쌓인다. 이 덱에는 다루지 않으므로
        열린 질문으로만 남긴다.
  \item 평균, 분산, 분포 전체를 정식 통계량으로 비교하는 방법은 \Cref{ch:compare-dist}가 이어받는다.
\end{itemize}


% 2부 - 실전. 결과가 아니라 판단의 과정을 남긴다.
\part{실전}

\input{chapters/05-hw1-casebook}
\chapter{과제 2: 데이터와 분포}
\label{ch:hw2}

\begin{aside}
결과 자체는 \texttt{hw2/report/}에 있다. (\texttt{hw2/report\_submit/}은 채점 제출용
빈 템플릿이라 답이 들어 있지 않으므로, 이 장은 \texttt{hw2/report/}를 유일한 출처로
쓴다. 두 파일의 차이는 \Cref{sec:hw2-learned}에 적는다.) 이 장에는 \textbf{판단의 과정}을
남긴다. 무엇을 망설였고, 왜 그렇게 정했는지를 적는다. 모든 수치는
\texttt{hw2/code/p1.py}, \texttt{hw2/code/p2.py}를 다시 실행해 재확인했다.
\end{aside}

\section{데이터 세 가지}
\label{sec:hw2-data}

세 데이터셋의 크기와 성질이 이미 어떤 모델이 맞을지 힌트를 준다.

\begin{itemize}
  \item \textbf{공항 노선 수(Airport Routes)}: 공항 3409곳, 노선 수는 정수, 최소 1에서
    최대 915. 평균 19.845, 중앙값 4로 평균이 중앙값의 약 5배다. 평균과 중앙값이 크게
    벌어지는 것은 오른쪽 꼬리가 길다는 신호다(\Cref{sec:skew-sign}).
  \item \textbf{영화 평점(Movie Votes)}: 영화 4392편, 평점은 0.1 단위로 반올림된 실수,
    범위 1.9에서 8.5. 평균 6.227, 중앙값 6.3으로 둘이 거의 같다. 대칭에 가깝다는 신호다.
  \item \textbf{배우 나이(Actor Age)}: 원본 460{,}718개 생년월일 중 정제 후 458{,}015명
    (\Cref{sec:hw2-age}에서 정제 과정을 설명한다). 연 단위 정수, 15에서 99.
\end{itemize}

\section{모수 추정: 네 가지 분포, 하나의 원리}
\label{sec:hw2-estimators}

\begin{inbrief}
  \item[Why] 데이터가 어떤 분포를 따르는지 주장하려면 먼저 후보 분포마다 모수(Parameter)를
    데이터에서 뽑아내야 비교할 대상이 생긴다.
  \item[When] 후보 분포의 형태(멱함수, 지수, 균등, 정규)를 고정하고 그 안의 숫자만
    추정할 때.
  \item[Where] 뒤에 나오는 시각적 비교(\Cref{sec:hw2-compare-how})와 적합도 판정
    전 단계. 추정이 틀리면 비교 자체가 무의미하다.
  \item[How] 균등·지수·멱법칙 세 분포는 \term{최대우도추정}(Maximum Likelihood
    Estimation, MLE)을 그대로 쓴다. 정규분포만 $\sigma$를 $n-1$로 나눈 표본표준편차로
    바꿔 쓴다(자유도 보정). 이 $\sigma$를 뺀 나머지 모수는 모두 데이터가 관측될
    확률(우도, Likelihood)을 가장 크게 만드는 값이다.
\end{inbrief}

문제는 모수 하나(예: 정규분포의 $\mu$)가 아니라 \textbf{모형 자체가 맞는지}를 물었다.
그런데 모형이 맞는지 보려면 먼저 각 모형이 이 데이터에 가장 잘 맞을 때 모수가 얼마인지
정해야 공정한 비교가 된다. 임의의 모수로 비교하면 모형이 아니라 모수를 잘못 고른
탓에 진 것처럼 보일 수 있다. 그래서 네 모형 모두 모수를 먼저 고정한 뒤에만
\Cref{sec:hw2-compare-how}에서 비교한다.

\paragraph{균등분포(Uniform Distribution)와 정규분포(Normal Distribution).}
가장 단순한 둘부터 본다. 균등분포 $[a, b]$의 MLE는 관측값의 최솟값과 최댓값이다.
정규분포(\Cref{sec:normal})의 $\mu$는 MLE(표본평균, Sample Mean) $\hat\mu = \bar x$를
쓴다. $\sigma$는 MLE(즉 $n$으로 나눈 값) 대신 $n-1$로 나눈 \term{표본표준편차}(Sample
Standard Deviation)를 썼다. $n$으로 나누면 분산을 과소추정하기 때문에(자유도 보정,
Bessel's Correction) \textbf{분산}의 불편추정량(Unbiased Estimator) $s^2$를 택한 것이다.
표준편차 $s$ 자체는 여전히 약간 편향되어 있다. $n=3409, 4392$처럼 $n$이 큰 경우
두 값의 차이는 소수 셋째 자리 이하다.

\paragraph{지수분포(Exponential Distribution).}
지수분포(\Cref{sec:exponential})는 평균 대기 시간이 $1/\lambda$인 분포다. 로그우도
(Log-Likelihood)를 $\lambda$에 대해 미분해 0으로 놓으면 $\hat\lambda$가 나온다.
\begin{align*}
  \ell(\lambda) &= n\ln\lambda - \lambda\sum_i x_i   && \text{로그우도}\\
  \frac{d\ell}{d\lambda} &= \frac{n}{\lambda} - \sum_i x_i  && \text{$\lambda$로 미분}\\
  0 &= \frac{n}{\hat\lambda} - \sum_i x_i  && \text{기울기가 0인 점}\\
  \hat\lambda &= \frac{n}{\sum_i x_i} = \frac{1}{\bar x}   && \text{정리}
\end{align*}
즉 표본평균의 역수 하나면 끝난다.

\begin{example}[공항 데이터에 지수분포 MLE 적용]
\label{ex:hw2-exp-mle}
공항 노선 수의 평균은 $\bar x = 19.845$다.
\begin{align*}
  \hat\lambda &= 1/\bar x                 && \text{MLE 공식}\\
              &= 1/19.845                 && \text{평균 대입}\\
              &= 0.0504                   && \text{반올림}
\end{align*}
$\lambda$가 작을수록 대기 시간(여기서는 노선 수)의 기댓값이 크다.
\end{example}

\paragraph{멱법칙(Power Law).}
지수분포는 노선 500개짜리 공항을 사실상 불가능하다고 본다($\lambda=0.05$면 확률이
약 $e^{-25}$). 그런데 실제로 500개가 넘는 공항이 5곳 있다(\Cref{sec:hw2-airport}).
꼬리가 지수보다 느리게 줄어드는 모형이 필요하다.

\term{멱법칙}(Power Law) 분포는 $x \geq x_{\min}$에서 $f(x) \propto x^{-\alpha}$
형태의 확률밀도함수를 갖는다. 지수분포처럼 큰 값일수록 드물지만, 꼬리가 지수적으로
줄지 않고 \textbf{거듭제곱으로만} 줄어들어 훨씬 천천히 감소한다. 네트워크의 연결 수,
도시 인구, 부의 분포처럼 소수가 대부분을 차지하는 데이터에서 흔히 등장한다.

\begin{pitfall}[책에 없는 분포]
멱법칙은 \Cref{ch:distributions}에 아직 절이 없다. 여기서 처음 정의한다.
\end{pitfall}

연속형 멱법칙의 MLE는 Clauset 등(2009)의 공식을 쓴다.
\[
  \hat\alpha = 1 + \frac{n}{\sum_{i} \ln(x_i / x_{\min})}
\]
$x_{\min}$은 멱법칙이 시작되는 하한이다. 이 과제에서는 $x_{\min} = \min_i x_i$로
정했다(데이터의 최솟값 자체를 하한으로 삼는 가장 단순한 선택). 이 선택이 공항
데이터에서 $x_{\min}=1$, 영화 데이터에서 $x_{\min}=1.9$를 준다.

\begin{example}[공항 데이터에 멱법칙 MLE 적용, 축약]
\label{ex:hw2-powerlaw-mle}
$n = 3409$, $x_{\min} = 1$일 때
\begin{align*}
  \sum_i \ln(x_i/x_{\min}) &= \sum_i \ln x_i    && x_{\min}=1\text{이므로 로그가 그대로 } \ln x_i \\
                            &= 5569.4            && \text{코드로 합산한 값}\\
  \hat\alpha &= 1 + n/\sum_i \ln(x_i/x_{\min})   && \text{공식}\\
             &= 1 + 3409/5569.4                  && \text{값 대입}\\
             &= 1 + 0.612                        && \text{나눗셈}\\
             &= 1.612                            && \text{반올림}
\end{align*}
\end{example}

\Cref{tab:hw2-params}는 네 분포 모두를 두 데이터셋에 한 번에 정리한다.

\begin{table}[htbp]
\centering
\begin{tabular}{lrr}
\toprule
모수(Parameter) & 공항 노선(Airport) & 영화 평점(Movie) \\
\midrule
$\alpha$ (멱법칙 Power Law, $x_{\min}$ 포함) & $1.612$ ($x_{\min}=1$) & $1.851$ ($x_{\min}=1.9$) \\
$\lambda$ (지수 Exponential)      & $0.0504$               & $0.1606$ \\
$[a, b]$ (균등 Uniform)           & $[1, 915]$             & $[1.9, 8.5]$ \\
$\mu, \sigma$ (정규 Normal, $n{-}1$) & $19.845,\ 53.506$   & $6.227,\ 0.893$ \\
\bottomrule
\end{tabular}
\caption{두 데이터셋에서 추정한 네 분포의 모수. $n = 3409$(공항), $n = 4392$(영화).}
\label{tab:hw2-params}
\end{table}

\section{네 모형을 비교하는 방법}
\label{sec:hw2-compare-how}

\begin{inbrief}
  \item[Why] 모수를 정했다고 모형이 맞는지 알 수 있는 것은 아니다. 데이터와 모형을
    같은 그림 위에 겹쳐야 어디가 맞고 어디가 틀리는지 눈으로 확인할 수 있다.
  \item[When] 후보 분포가 여럿이고 하나를 골라야 할 때.
  \item[Where] PDF, CDF, QQ 도표(\Cref{sec:cdf}) 중 과제는 하나만 요구했다. 세
    시각화가 서로 다른 정보를 준다.
  \item[How] 분포마다 데이터와 같은 개수 $n$의 난수를 뽑아, 데이터의 CCDF/PDF와
    모형 표본의 CCDF/PDF를 겹쳐 그리고, QQ 도표와 콜모고로프--스미르노프
    거리(Kolmogorov--Smirnov Distance) $D$로 요약한다.
\end{inbrief}

\paragraph{공항: CCDF를 로그-로그 축에.}
공항 노선 수는 정수이고 꼬리가 두 자릿수 넘게 퍼져 있다. 보통 히스토그램처럼 구간을
나누면(log bin) 큰 값 쪽에 관측치가 하나뿐인 구간이 많아 그래프가 듬성듬성해진다.
그래서 값별 확률 $P(X=k)$를 점으로 찍고, \term{여누적분포함수}(Complementary CDF,
CCDF) $P(X \geq x)$를 로그-로그(log-log) 축에 그렸다. 멱법칙은 로그-로그 축에서
\textbf{직선}이 된다는 성질이 있어(양변에 로그를 취하면 $\ln P(X\geq x) \approx
-(\alpha-1)\ln x + c$, 직선의 방정식) 이 축이 멱법칙 여부를 눈으로 바로 보여준다.

\paragraph{영화: 0.1 격자에 맞춘 히스토그램.}
영화 평점은 0.1 단위로 반올림된 값이다. 구간 경계가 0.1 격자와 어긋나면 히스토그램이
들쭉날쭉해진다(어떤 구간은 값 두 개를 담고, 바로 옆 구간은 하나도 못 담는 식). 그래서
구간 너비를 0.1 격자에 맞춰 PDF(히스토그램)를 그렸고, 꼬리가 좁아 로그 축 없이도
전체가 한 그림에 들어온다.

\paragraph{역변환 표본추출(Inverse Transform Sampling).}
네 분포 각각에서 데이터와 같은 개수 $n$의 난수를 뽑아야 한다. 지수·균등·정규는
\texttt{numpy}에 표본추출 함수가 있지만, 멱법칙은 없어서 직접 만들었다. 멱법칙의
CDF는
\[
  F(x) = 1 - (x/x_{\min})^{-(\alpha-1)}
\]
$u \sim \mathrm{Uniform}(0,1)$에 대해 $u = F(x)$를 $x$에 대해 풀면
\begin{align*}
  u &= 1 - (x/x_{\min})^{-(\alpha-1)}          && \text{CDF 정의}\\
  (x/x_{\min})^{-(\alpha-1)} &= 1-u            && \text{항 이항}\\
  x/x_{\min} &= (1-u)^{-1/(\alpha-1)}          && \text{양변을 } -(\alpha-1)\text{ 제곱근}\\
  x &= x_{\min}(1-u)^{-1/(\alpha-1)}           && \text{정리}
\end{align*}
$u$를 균등분포에서 뽑아 이 식에 넣으면 멱법칙을 따르는 $x$가 나온다. 난수 시드는
6330으로 고정해 그림을 재현할 수 있게 했다.

\paragraph{QQ 도표와 $D$.}
CCDF·PDF를 겹쳐 그리면 어디가 맞고 틀리는지는 보이지만, ``얼마나'' 다른지 숫자로
비교하기는 어렵다. 분포 전체를 한 번에 대조하면서도 숫자 하나로 요약할 도구가
필요하다.

\term{분위수 대 분위수 도표}(Quantile-Quantile Plot, QQ Plot)는 $x$축에 데이터
분위수, $y$축에 모형 표본 분위수를 찍는다. 점이 기준선 $y=x$ 위에 있으면 두 분포가
같은 모양이라는 뜻이고, 선에서 벗어난 방향이 어디서 어떻게 다른지 말해준다. 로그
축을 쓸 때는 표본에 음수나 0이 섞일 수 있어(정규 표본만 해당, 지수 표본은 항상
양수다) 그 값을 먼저 제외했다.
비교를 숫자 하나로도 요약했다. \term{콜모고로프--스미르노프 거리}(KS Distance)
$D$는 데이터의 경험적 CDF와 모형 \textbf{표본}의 경험적 CDF 사이 세로 격차의
최댓값이다(\Cref{sec:cdf}). 모형 자체의 이론적 CDF가 아니라 모형에서 뽑은 표본을
쓰므로, $D$는 그 표본을 뽑은 난수 시드에 약간 좌우된다. $D=0$이면 두 CDF가 동일하다.
이 $D$를 검정 통계량으로 쓰는 두 표본 KS 검정(Two-sample KS Test)은 \Cref{sec:cd-ks-two}에서 다룬다.

\section{공항 노선 수: 멱법칙, 상한이 있는}
\label{sec:hw2-airport}

\begin{figure}[htbp]
\centering
\includegraphics[width=0.95\linewidth]{p1a_0_data.png}
\caption{공항 노선 수의 경험적 분포(Empirical Distribution). 왼쪽은 값별 확률 $P(X=k)$, 오른쪽은 CCDF
$P(X\geq x)$, 둘 다 로그-로그 축. 38\%의 공항이 노선 1--2개뿐이고, 가장 큰
공항은 915개다.}
\label{fig:hw2-airport-data}
\end{figure}

\Cref{fig:hw2-airport-data}에서 왼쪽 PDF가 이미 직선에 가깝다. 이는 멱법칙의
특징이다. 38.0\%의 공항이 노선 1--2개뿐이고, 99번째 백분위수는 285개인데 최댓값은
915개다. 왜도(Skewness)는 6.0으로(\Cref{sec:skewness}) 강하게 오른쪽으로 치우쳐
있다.

\begin{figure}[htbp]
\centering
\includegraphics[width=0.95\linewidth]{p1a_1_powerlaw.png}
\caption{공항 노선 수 대 멱법칙 모형($\alpha=1.612$, $x_{\min}=1$). CCDF는 1에서
약 60개 노선까지 데이터를 따라가고, QQ 점도 그 구간에서 $y=x$ 위에 있다. 100개를
넘으면 모형은 계속 뻗어 $10^4$--$10^5$개짜리 공항을 만들어내지만, 데이터는 915에서
멈춘다. $D=0.209$.}
\label{fig:hw2-airport-powerlaw}
\end{figure}

\Cref{tab:hw2-airport-verdict}가 네 모형을 요약한다.

\begin{table}[htbp]
\centering
\begin{tabular}{lrrl}
\toprule
분포(Distribution) & 모수(Parameter) & KS 거리 $D$ & 판정(Verdict) \\
\midrule
멱법칙(Power Law)     & $\alpha=1.612,\ x_{\min}=1$ & \textbf{0.209} & 최선, 꼬리는 과대추정 \\
지수(Exponential)     & $\lambda=0.0504$            & 0.379          & 꼬리가 너무 빨리 줄어듦 \\
정규(Normal)           & $\mu=19.8,\ \sigma=53.5$   & 0.358          & 음수 36\% 허용, 구조적 실패 \\
균등(Uniform)          & $[1, 915]$                  & 0.857          & 최악, 형태 자체가 다름 \\
\bottomrule
\end{tabular}
\caption{공항 노선 수: 네 모형의 콜모고로프--스미르노프 거리(KS Distance) $D$와 판정.}
\label{tab:hw2-airport-verdict}
\end{table}

\textbf{공항 데이터는 멱법칙을 따르되, 꼬리에 상한이 있다.} 로그-로그 축에서 CCDF가
직선인 것은 멱법칙뿐이고, 데이터의 CCDF도 1에서 약 100개 노선까지 두 자릿수 넘게
직선에 가깝다. 나머지 세 모형은 모수를 조금 바꿔도 구조적으로 못 맞는다.

\begin{itemize}
  \item \textbf{균등}은 노선 1개짜리와 915개짜리를 똑같이 흔하다고 가정하지만,
    대부분의 공항은 작다.
  \item \textbf{정규}는 대칭이고 음수를 허용한다. $\sigma$가 $\mu$의 거의 3배라
    모형 질량의 36\%가 0 미만인데, 노선 수는 음수일 수 없다. 데이터는 하한 1에서
    시작해 오른쪽으로 치우쳐 있다(왜도 6.0, \Cref{sec:skew-sign}).
  \item \textbf{지수}는 모양은 맞다(작은 값이 많고 큰 값이 드물다). 그러나 꼬리가
    너무 빨리 줄어든다. $\lambda=0.05$면 노선 500개짜리 공항의 확률은 약
    $e^{-25}$인데, 실제로 500개를 넘는 공항이 5곳 있다.
\end{itemize}

멱법칙도 완벽하지는 않다. $\hat\alpha=1.61 < 2$이므로 이론적으로 평균이 발산하는
모형이고, 표본에는 $10^5$개짜리 공항이 등장한다(\Cref{fig:hw2-airport-powerlaw}).
실제로는 공항이 3409곳뿐이고 활주로·게이트·슬롯이 허브의 성장을 막는다. 그래서
데이터의 CCDF는 100개를 넘는 지점부터 아래로 꺾인다. 가장 적절한 설명은
\textbf{상한이 있는 멱법칙}이며, 이는 \term{선호적 연결}(Preferential Attachment)로
성장하는 네트워크의 전형적인 모양이다. 이미 노선이 많은 공항에 새 노선이 더 많이
붙는다는 뜻이다.

\section{영화 평점: 정규분포에 가깝다}
\label{sec:hw2-movie}

\begin{figure}[htbp]
\centering
\includegraphics[width=0.95\linewidth]{p1b_0_data.png}
\caption{영화 평점의 경험적 분포(Empirical Distribution). 왼쪽은 0.1 격자에 맞춘 히스토그램(PDF), 오른쪽은
CDF. 중앙값 6.3, 평균 6.23으로 6.3 근처에서 한 개의 봉우리를 이루며 거의 대칭이다.
왼쪽 꼬리가 약간 더 길다(왜도 $-0.48$).}
\label{fig:hw2-movie-data}
\end{figure}

\begin{figure}[htbp]
\centering
\includegraphics[width=0.95\linewidth]{p1b_4_normal.png}
\caption{영화 평점 대 정규분포 모형($\mu=6.227$, $\sigma=0.893$). 모형 표본이
데이터와 같은 종 모양을 그리고, QQ 점은 약 4.5에서 7.5까지 $y=x$ 위에 있다. 양 끝만
어긋난다. 데이터의 최저 평점이 모형 예측보다 더 낮고(왼쪽 꼬리가 더 김), 데이터는
8.5에서 끝나지만 모형은 계속된다. $D=0.055$로 네 모형 중 가장 작다.}
\label{fig:hw2-movie-normal}
\end{figure}

\begin{table}[htbp]
\centering
\begin{tabular}{lrrl}
\toprule
분포(Distribution) & 모수(Parameter) & KS 거리 $D$ & 판정(Verdict) \\
\midrule
정규(Normal)           & $\mu=6.227,\ \sigma=0.893$   & \textbf{0.055} & 최선, 양 끝만 이탈 \\
균등(Uniform)           & $[1.9, 8.5]$                  & 0.394          & 봉우리를 무시함 \\
멱법칙(Power Law)      & $\alpha=1.851,\ x_{\min}=1.9$ & 0.485          & 단조감소, 구조적 실패 \\
지수(Exponential)      & $\lambda=0.1606$              & 0.490          & 단조감소, 최악 \\
\bottomrule
\end{tabular}
\caption{영화 평점: 네 모형의 KS 거리 $D$와 판정.}
\label{tab:hw2-movie-verdict}
\end{table}

\textbf{영화 평점은 근사적으로 정규분포를 따른다.} $D=0.055$는 다음으로 좋은 모형의
약 7분의 1이다. 정규분포는 $\pm1\sigma$ 안에 68.3\%, $\pm2\sigma$ 안에 95.4\%를
예측하는데, 실제 데이터는 각각 70.1\%와 95.6\%로 거의 일치한다.

나머지 세 모형은 구조적으로 실패한다. 멱법칙과 지수분포는 둘 다
\textbf{단조감소}(monotonically decreasing) 형태라 가장 가능성이 높은 값이 가장
작은 값이다. 즉 대부분의 영화가 최악의 평점을 받고 소수만 매우 높은 평점을 받는다고
예측하는데, 실제 데이터는 6.3 근처에 봉우리 하나가 있다. 멱법칙 표본은 13.3\%가
20점을 넘어가는데, 평점은 1에서 10 사이만 가능하다. 균등분포는 봉우리 자체를
무시한다.

정규 적합도 완벽하지 않다. 첫째, 평점은 $[1, 10]$에 갇혀 있고 약간 왼쪽으로
치우쳐 있다(\Cref{sec:skew-sign}). 나쁜 영화는 평균보다 훨씬 더 낮게 갈 수 있지만
좋은 영화가 평균보다 더 높이 갈 여지는 상대적으로 작다. 둘째, 평점이 소수점
한 자리로 반올림되어 QQ 도표에 작은 계단이 생긴다.

\term{중심극한정리}(Central Limit Theorem, CLT)는 영화 한 편의 평점이 그 영화의
개별 투표 여러 개의 평균이라는 사실을 설명해 줄 수는 있지만, 그것이 설명하는 대상은
영화 \textbf{한 편} 안에서 평점이 진짜 평균 주변에 얼마나 좁게 모이는지이지, 서로
다른 영화 4392편의 평균 평점이 \textbf{영화 간}에 왜 종 모양으로 퍼지는지가 아니다.
후자는 영화마다 다른 진짜 작품성의 분포에서 온다. 개별 투표의 잡음이 평균에서
줄어드는 것과, 영화 간 평점 분포가 정규 모양을 띠는 것은 서로 다른 현상이며, 이
장에서 관찰한 것은 후자, 즉 \Cref{fig:hw2-movie-data}와 \Cref{tab:hw2-movie-verdict}가
보여주는 데이터 자체의 형태다.

\section{나이, 첫째 자리, 마지막 자리}
\label{sec:hw2-age}

\begin{pitfall}[생년월일 정제: 두 자리 연도와 결측 표시]
생년월일은 \texttt{dd-Mon-yy} 형식이 대부분이고 일부만 4자리 연도다. 모든 두 자리
연도를 \textbf{19\textit{yy}}로 읽었다. 이 데이터(IMDB-WIKI)는 2015년에 수집됐으므로
``20''--``24''는 2020--2024년일 수 없고, 연도별 개수를 세어 보면 1920년대부터
``00''까지 매끄럽게 줄어들 뿐 2000년대로 전환됐다고 볼 만한 단절이 없다. 이 확인을
거친 뒤에야 19\textit{yy} 규칙을 확정했다(짐작으로 정하지 않았다). 연도 ``0000''은
\term{결측값}(Missing Value) 표시로 보고 제외했다(129건). 나이는
\textbf{2015-01-01}(수집 연도) 기준
만 나이이고, 1에서 99 사이만 남겼다. 그 결과 \textbf{458{,}015명}이 남는다(평균
45.3, 중앙값 43, 범위 15--99). 최솟값이 15인 것은 데이터의 성질이 아니라
19\textit{yy} 규칙의 결과다. 가장 늦은 출생연도가 1999년이므로 2015년 기준
최소 나이가 15로 정해진다.
\end{pitfall}

\begin{figure}[htbp]
\centering
\includegraphics[width=0.8\linewidth]{p2_1_age.png}
\caption{2015년 기준 나이 분포($n=458{,}015$). 30대와 40대에 넓은 봉우리가
하나 있고, 99세까지 긴 오른쪽 꼬리가 이어진다. 대부분의 질량은 25에서 60 사이다.}
\label{fig:hw2-age}
\end{figure}

\begin{figure}[htbp]
\centering
\includegraphics[width=0.8\linewidth]{p2_2_first_digit.png}
\caption{나이의 첫째 자리. 점선은 균등 기준선 $1/9$, 주황 곡선은
벤포드 법칙(Benford's Law) $\log_{10}(1+1/d)$. 자릿수 1부터 9까지 비율은
1.3\%, 11.3\%, 26.2\%, 28.3\%, 16.7\%, 9.0\%, 4.4\%, 2.2\%, 0.7\%다. 데이터는
균등과도 벤포드와도 거리가 멀다. 벤포드는 1에서 꾸준히 감소하지만 데이터는 3과
4에서 정점을 찍는다.}
\label{fig:hw2-first-digit}
\end{figure}

\begin{figure}[htbp]
\centering
\includegraphics[width=0.8\linewidth]{p2_3_last_digit.png}
\caption{나이의 마지막 자리. 점선은 균등 기준선 $1/10$. 모든 자릿수가
8.95\%에서 11.01\% 사이, 즉 10\%에서 최대 1.1퍼센트포인트 안쪽에 있다.}
\label{fig:hw2-last-digit}
\end{figure}

\textbf{마지막 자리는 (근사적으로) 균등하고, 첫째 자리는 그렇지 않다.}

\paragraph{마지막 자리.} 이는 예상된 결과다. 나이의 마지막 자리는 나이를 10으로 나눈
나머지다. 20대, 30대처럼 10년 단위 구간마다 마지막 자리 0에서 9가 정확히 한 번씩
들어 있다. 나이 분포는 한 살 단위로 보면 완만하게 변하므로, 어느 10년 구간을 보든
그 안의 열 개 나이는 거의 같은 비율을 차지하고, 이를 모든 구간에 걸쳐 더하면 각
마지막 자리가 약 10\%씩 나온다. 남은 $\pm1.1$퍼센트포인트의 편차는 특정 자릿수를
선호해서가 아니라, 특정 출생연도에 배우가 몇 명 태어났는지의 연도별 잡음
(\Cref{fig:hw2-age}의 들쭉날쭉한 막대)에서 온다.

\paragraph{첫째 자리.} 이것도 사실은 예상된 결과다. 나이 10에서 99까지는 첫째
자리가 곧 \textbf{인생의 10년 단위 구간(decade)}이다. 그래서 첫째 자리의 히스토그램은
나이 분포를 10년 단위로 묶은 것과 같은 모양이고, 실제로 \Cref{fig:hw2-age}와 같은
형태(봉우리 하나 + 오른쪽 꼬리)를 보인다. 자릿수 1이 거의 비어 있는 것은 15--19세만
가능한 좁은 구간이라 그 나이대 배우가 적기 때문이고, 9가 거의 비어 있는 것은 90대에
이르는 배우가 드물기 때문이다. 첫째 자리가 균등하려면 10대부터 90대까지 모든
10년 구간에 배우 수가 같아야 하는데, 활동 중인 배우 집단은 그렇지 않다.

\paragraph{벤포드 법칙도 아니다.} \term{벤포드 법칙}(Benford's Law, 첫째 자리
$d$의 확률이 $\log_{10}(1+1/d)$이라 1이 30.1\%로 가장 흔함)은 값이 여러 자릿수
(order of magnitude)에 걸쳐 퍼져 있는 데이터, 예를 들어 인구, 가격, 또는
\Cref{sec:hw2-airport}의 공항 노선 수처럼 로그값이 고르게 퍼진 데이터에서 성립한다.
나이는 15에서 99까지 한 자릿수 범위 안에 있고 45 근처에 몰려 있어서, 첫째 자리는
벤포드의 패턴이 아니라 그 봉우리가 어디 있는지를 보여줄 뿐이다. 요약하면, 마지막
자리는 값이 10년 구간의 어디에 있는지만 남기고 나머지 정보를 버리기 때문에
균등하고, 첫째 자리는 바로 그 정보(어느 구간인가)를 그대로 담고 있어서 나이
분포의 모양을 따라간다.

\section{이 과제에서 배운 것}
\label{sec:hw2-learned}

\textbf{모형이 맞는지는 모수 하나가 아니라 꼬리의 모양으로 갈린다.} 공항은
$\sigma$(53.5)가 $\mu$(19.8)의 거의 3배라 평균과 표준편차만 봐도 이미 정규분포가
아니다(\Cref{sec:hw2-airport}). 영화는 평균과 중앙값이 거의 같고 왜도도 작아 정규분포
후보로 보였고, 실제로 그렇게 나왔다(\Cref{sec:hw2-movie}). 두 경우 모두 갈림을 만든
것은 꼬리의 모양이었다. 공항은 꼬리가 지수보다 느리게 줄어 멱법칙이 이겼고, 영화는
꼬리가 짧고 대칭이라 정규분포가 이겼다.

\textbf{KS 거리 하나로는 ``왜'' 지는지 알 수 없다.} \Cref{tab:hw2-airport-verdict}와
\Cref{tab:hw2-movie-verdict}의 $D$ 값은 순위를 매겨 줄 뿐, 그 순위의 이유는 말해주지
않는다. 공항에서는 정규($D=0.358$)가 지수($D=0.379$)보다 근소하게 \textbf{낫지만},
정규는 모형 질량의 36\%를 음수 영역에 두는 구조적 실패이고 지수는 꼬리가 너무 빨리
줄어드는 실패다. $D$가 더 작다고 더 타당한 모형은 아니다. 숫자 뒤에 그림을 반드시
같이 본 이유다.

\textbf{자릿수 분포는 원본 분포의 그림자다.} 마지막 자리가 균등한 것과 첫째 자리가
균등하지 않은 것은 서로 다른 현상이 아니라, 같은 나이 분포를 \textbf{정보를 버리는
정도가 다른 두 렌즈}로 본 것이다. 벤포드 법칙이 성립하려면 값이 여러 자릿수에
걸쳐 퍼져 있어야 한다는 조건이 있고, 나이는 그 조건을 만족하지 않는다.

마지막으로 자료 확인 하나. \texttt{hw2/report\_submit/hw2\_taeklim\_report.tex}는
채점 제출용 빈 템플릿으로, 문항마다 \texttt{STEP} 주석만 있고 그림과 답은 모두
주석 처리되어 있다. 실제 답과 논의, 그리고 그림·표는 전부
\texttt{hw2/report/hw2\_taeklim\_report.tex}에 있다. 이 장의 모든 수치와 그림은
후자를 출처로 삼았다.

\input{chapters/06-project-log}


\appendix
\input{parts/part3-appendix}

\backmatter
\bibliographystyle{plainnat}
\bibliography{bib/references}

\end{document}
